% ======================================================================
% DOCUMENT PREAMBLE
% ======================================================================
\documentclass[11pt, a4paper]{article}

\usepackage[top=1cm, bottom=2cm, left=1.5cm, right=1.5cm]{geometry}
\usepackage{fontspec}
\usepackage[english, bidi=basic, provide=*]{babel}
\babelprovide[import, onchar=ids fonts]{english}
\babelfont{rm}{Times New Roman} % Standard academic serif

\usepackage{microtype}
\usepackage{amsmath, amssymb, amsthm}
\usepackage{physics}
\usepackage{multicol}
\usepackage{enumitem}
\usepackage{xcolor}
\usepackage{titlesec}
\usepackage{hyperref}
\usepackage[most]{tcolorbox} 

% Define a professional "Academic Sepia/Charcoal" palette
\definecolor{academiabrown}{RGB}{70, 50, 40}
\hypersetup{colorlinks=true, linkcolor=academiabrown, citecolor=academiabrown, urlcolor=blue, hidelinks}

% --- Define the Dynamic Box Styles ---
\tcbset{
    cheatbox/.style={
        enhanced, unbreakable,
        colback=academiabrown!4!white, colframe=academiabrown!90!black,   
        fonttitle=\bfseries, coltitle=white, boxrule=0.8pt, arc=2mm,                   
        left=2mm, right=2mm, top=1.5mm, bottom=1.5mm, boxsep=1mm,
        toptitle=1.5mm, bottomtitle=1.5mm
    },
    recallbox/.style={
        enhanced, unbreakable,
        colback=academiabrown!4!white, colframe=academiabrown!90!black,
        boxrule=0.8pt, arc=2mm, left=2mm, right=2mm, top=1mm, bottom=1mm, boxsep=0mm
    }
}

% Custom section formatting
\titleformat{\section}{\Large\bfseries\color{black}}{\thesection.}{0.5em}{}
\titleformat{\subsection}{\large\bfseries\color{darkgray}}{\thesubsection}{0.5em}{}
\titleformat{\subsubsection}{\normalsize\bfseries\color{gray}}{\thesubsubsection}{0.5em}{}

\setlength{\columnsep}{0.6cm}

% --- Dynamic Numbering Setup ---
\newcounter{cheatnum}
\setcounter{cheatnum}{0}

% ======================================================================
% TITLE & ABSTRACT
% ======================================================================
\title{\textbf{Refractive Cosmology \& Continuum Mechanics \\ CONDENSED REFERENCE \& FORMAL DERIVATIONS}}
\author{\textbf{Binyamin Tsadik Bair-Mosheh} }
\date{\today}

\begin{document}

\maketitle


\begin{tcolorbox}[colback=gray!5!white, colframe=gray!50!black, arc=1mm, boxrule=0.5pt, left=3mm, right=3mm]
\begin{center}
    \textbf{ABSTRACT \& HOW TO USE THIS MANUAL}
\end{center}
\textit{This Is a \textbf{condensed} version of the Master Reference Sheet.
This physics establishes that macroscopic cosmology, localized quantum mechanics, and solid-state phenomena evaluate  as the deterministic consequences of continuum mechanics operating within a dynamic elastic vacuum. By defining the uncurled metric as a compressible, refractive fluid, the foundations of physical law are systematically derived from classical fluid dynamics, Eulerian material derivatives, and geometric topology.}
\vspace{0.2cm}
\begin{itemize}[nosep]
    \item \textbf{PART I} provides the dynamically numbered Cheat Sheet containing the raw equations for quick use and reference.
    \item \textbf{PART II} \hyperref[sec:part2]{\textbf{[Click Here]}} provides explicit, step-by-step mathematical proofs from first principles. 
    \item The \textbf{[Proof $\downarrow$]} links in each box navigate directly to the corresponding formal derivation.
\end{itemize}
\end{tcolorbox}

\vspace{0.3cm}

\begin{tcolorbox}[cheatbox, title={The Conformal Duality of the Metric \scriptsize\hfill\hyperref[sec:part2]{[See Derivations $\downarrow$]}}]
The vacuum evaluates as a continuous fluid metric, characterized by conformal thermodynamic symmetry between the macroscopic cosmos and localized topological gradients:
\begin{itemize}[leftmargin=*, nosep]
    \item \textbf{The Cosmos (The Macroscopic Limit):} Expanding causal wavefront. The outward macroscopic velocity field approaches the absolute acoustic limit ($\mathbf{v}_0 \to c$) as $r \to \infty$, maximizing dynamic pressure and driving the ambient static pressure of the boundary signal to zero.
    \item \textbf{The Particle (The Localized Core):} Sequestered topological void. The localized fluid circulation approaches the absolute acoustic limit ($\mathbf{v}_{\phi} \to c$) as $r \to 0$, maximizing dynamic pressure and driving the internal static pressure to a minimum at the core.
\end{itemize}
\end{tcolorbox}

\vspace{4cm}


\begin{center}
    \LARGE\textbf{Condensed Version}\\ 
    \vspace{1cm}
    \small\textit{The full version is currently obsolete but will be updated eventually:  \\ \hyperlink{https://zenodo.org/records/20368016}{https://zenodo.org/records/20368016}}
\end{center}


% ======================================================================
% NOMENCLATURE & KINEMATIC CONSTANTS
% ======================================================================

\begin{tcolorbox}[colback=white, colframe=academiabrown!90!black, arc=1mm, boxrule=0.8pt, title={\textbf{TABLE 1: KINEMATIC NOMENCLATURE \& METRIC CONSTANTS}}]
\renewcommand{\arraystretch}{1.4}
\begin{tabular}{p{0.18\textwidth} p{0.77\textwidth}}
    $\boldsymbol{\tau}$ & \textbf{Absolute Metric Time.} The current, unmediated time of the continuum. \\
    $t$ & \textbf{Proper Time.} The evolutionary time within the metric. \\
    $\tau_0$ & \textbf{Cosmological Age.} The local proper age of the universe. \\
    $R_{h,0}$ & \textbf{Macroscopic Scale Radius.} The causal horizon at the present observation epoch. \\
    $R_{h,p} \equiv \lambda_c$ & \textbf{Topological Boundary Radius.} The localized Compton kinematic limit. \\
    $l_p$ & \textbf{Absolute Cavitation Limit.} The spatial boundary of fluid cavitation (Planck length) where circulation strikes the acoustic limit ($v_\phi \to c$). \\
    $v_0(r)$ & \textbf{Macroscopic Velocity Field.} The scalar kinematic velocity field of the expanding metric. \\
    $v_{obs}(r) \equiv v_\phi(r)$ & \textbf{Topological Circulation Velocity.} The localized, closed-loop phase velocity of the fluid metric. Evaluates dynamically as the inverse conformal projection of the cosmos. \\
   $P_c \equiv \frac{1}{2}\rho_{\tau}c^2$ & \textbf{Absolute Acoustic Yield Pressure.} The structural static pressure limit of the vacuum. Acts mechanically as the baseline isotropic tension ($T_c = -P_c$) of the spatial continuum. Evaluates macroscopically as the absolute energy density of rest mass. \\
    $P_{dyn}$ & \textbf{Scalar Pressure Deficit.} The localized dynamic pressure drop generated by topological circulation. Evaluates macroscopically as the Newtonian gravitational potential. \\
    $n(r)$ & \textbf{Geometric Refractive Index.} The localized index at radial coordinate $r$ dictating wave propagation along a null geodesic ($ds^2 = 0$). \\
    $De$ & \textbf{Cosmological Deborah Number.} The kinematic ratio ($\omega t_r$) dictating the rheological phase of the metric. \\
    $W$ & \textbf{Topological Winding Number.} The chirality or charge of an acoustic cavity, evaluating as the parity of the Dirac spinor ($\psi, \bar{\psi}$). \\
    $\alpha_s \equiv \frac{1}{\gamma}$ & \textbf{Total Scalar Admittance (Fluid Compliance).} The absolute volumetric yield of the metric against total kinematic capacity ($v_{tot}$). Evaluates physically as the residual static pressure ratio ($1 - P_{dyn}/P_c$), tracking numerically identical to longitudinal admittance during linear translation. Evaluates phenomenologically as the running Strong Coupling Constant. \\
    $\alpha$ & \textbf{Baseline Transverse Admittance.} The absolute resting transverse elastic yield of the macroscopic continuum. Defines the strict off-diagonal tensor coupling limit between internal topological vorticity and ambient transverse shear ($\alpha_{\circlearrowleft \perp} \equiv \alpha$). Evaluates phenomenologically as the Fine-Structure Constant ($\approx 1/137$). \\
    $[\boldsymbol{\alpha}]$ & \textbf{Kinematic Admittance Tensor.} The anisotropic viscoelastic tensor governing directional metric compliance against kinematic strain. Its diagonal components evaluate   as the orthogonal modal admittances: longitudinal ($\alpha_{\parallel}$), transverse ($\alpha_{\perp}$), and rotational ($\alpha_{\circlearrowleft}$). \\
    $[\boldsymbol{Z}_{\tau}] \equiv [\boldsymbol{\alpha}]^{-1}$ & \textbf{Kinematic Impedance Tensor.} The absolute structural resistance of the metric against applied kinematic strain, evaluating mathematically as the exact inverse matrix of fluid compliance. Its individual modal limits ($Z_i \equiv 1/\alpha_i$) project directionally as apparent inertial mass ($Z_{\tau(\parallel)} = m^*$), localized optical density ($Z_{\tau(\perp)} \propto n$), and absolute rotational drag ($Z_{\tau(\circlearrowleft)} = \hbar$), while its total scalar magnitude isolates the Lorentz factor ($Z_s = \gamma$). \\
    $\epsilon_{\tau}$ & \textbf{Linear Mass Density.} The 1D density evaluating as 2D transverse elastic stiffness. \\
    $\rho_{\tau}$ & \textbf{Volumetric Mass Density.} The 3D density of the continuum evaluating as linear density divided by transverse yield area ($\rho_{\tau} = \epsilon_{\tau} / A_{yield}$). \\
\end{tabular}
\end{tcolorbox}
\newpage
% ======================================================================
% PART I: THE AXIOMATIC CHEAT SHEET
% ======================================================================
\begin{center}
    \LARGE\textbf{PART I: THE CHEAT SHEET}
\end{center}
\vspace{0.2cm}

\begin{multicols}{2}

% ==========================================
% PHASE I: THE MACROSCOPIC CONTINUUM
% ==========================================
\vspace{0.2cm}
\noindent\textbf{\large \color{academiabrown} PHASE I: THE MACROSCOPIC CONTINUUM}
\vspace{0.2cm}

\begin{tcolorbox}[cheatbox, title={\refstepcounter{cheatnum}\label{box:axiom}\thecheatnum. Fundamental Axiom \& Translation \scriptsize\hfill\hyperref[sec:deriv_axiom]{[Proof $\downarrow$]}}]
The macroscopic causal boundary expands linearly at $c$. The universal translation between Evolutionary Proper Time ($t$) and Relative Observational Redshift ($z$) is defined by the absolute horizon ($R_h$):
\begin{equation}\label{eq:axiom}
    R_h(t) = c \cdot t \iff 1+z = \frac{\tau_0}{t}
\end{equation}
\textit{Present-Day Calibration ($z=0$):}
\begin{align*}
    R_{h,0} &= 4224 \text{ Mpc}, &\tau_0 &= \frac{R_{h,0}}{c} \approx 13.78 \text{ Gyr}
\end{align*}
\end{tcolorbox}

\begin{tcolorbox}[cheatbox, title={\refstepcounter{cheatnum}\label{box:chronometer}\thecheatnum. The High-Redshift Chronometer \scriptsize\hfill\hyperref[sec:deriv_chronometer]{[Proof $\downarrow$]}}]
Proper time ($t$) and local scale radius ($R_h$) translate symmetrically across coordinate frames:
\begin{align}
    t(z) &= \frac{\tau_0}{1+z} \label{eq:tz} \\
    R_h(z) &= \frac{R_{h,0}}{1+z} \equiv c \cdot t(z) \label{eq:rz}
\end{align}
\end{tcolorbox}

\begin{tcolorbox}[cheatbox, title={\refstepcounter{cheatnum}\label{box:gr_bridge}\thecheatnum. The Geometric Spacetime Bridge \scriptsize\hfill\hyperref[sec:deriv_gr_bridge]{[Proof $\downarrow$]}}]
The fluid pressure gradient generates a continuous velocity field $\mathbf{v}_0(r)$. By the Equivalence Principle, the fluid potential energy ($\Phi = -v_0^2/2$) dictates the time-time structural curvature tensor ($g_{tt}$):
\begin{align}
    v_0(r) &= c\left(1-e^{-r/R_{h,0}}\right) \label{eq:v0} \\
    g_{tt} &= -\left( 1 - \frac{2 P_{dyn}}{\rho_{\tau} c^2} \right) \label{eq:pressure_curvature_identity}\\ 
    ds^2 &= -\left(1 - \frac{v_0^2}{c^2}\right)c^2 dt^2 + \frac{dr^2}{1 - \frac{v_0^2}{c^2}} \label{eq:metric}
\end{align}
\end{tcolorbox}

\begin{tcolorbox}[cheatbox, title={\refstepcounter{cheatnum}\label{box:kinematic_gradients}\thecheatnum. Kinematic Gradients \scriptsize\hfill\hyperref[sec:deriv_kinematic_gradients]{[Proof $\downarrow$]}}]
Apparent expansion evaluates as the spatial gradient ($H_{app} = dv_0/dr$). Absolute proper acceleration ($a_{cosmic}$) evaluates as the centripetal acceleration of the causal horizon:
\begin{align}
    H_{app}(z) &= \frac{c}{R_{h,0}(1+z)} \iff H_{prop} = \frac{1}{t} \label{eq:happ} \\
    a_{app}(z) &= \frac{c^2}{R_{h,0}(1+z)} \iff a_{cosmic} = \frac{c}{t} \label{eq:aapp}
\end{align}
\end{tcolorbox}

\begin{tcolorbox}[cheatbox, title={\refstepcounter{cheatnum}\label{box:optical_dilatation}\thecheatnum. Optical Dilatation \scriptsize\hfill\hyperref[sec:deriv_optical_dilatation]{[Proof $\downarrow$]}}]
Because the kinematic deceleration parameter of the continuum evaluates to absolute zero ($q=0$), the true macroscopic coordinate distance evaluates as the geometric integration of the optical stretch parameter:
\begin{equation}\label{eq:r_cosmic}
    r_{true, cosmic}(z) = R_{h,0} \ln(1+z)
\end{equation}
The apparent luminosity distance ($D_L$) evaluates geometrically as the true spatial coordinate multiplied by the macroscopic expansion scalar:
\begin{equation}\label{eq:dl}
    D_L(z) = R_{h,0} (1+z) \ln(1+z)
\end{equation}
\end{tcolorbox}

\begin{tcolorbox}[cheatbox, title={\refstepcounter{cheatnum}\label{box:acoustic_invariance}\thecheatnum. Cosmic Expansion \& Acoustic Invariance \scriptsize\hfill\hyperref[sec:deriv_acoustic_invariance]{[Proof $\downarrow$]}}]
The macroscopic fluid of the metric evaluates via a tensioned $w = -1/3$ equation of state. Because the absolute fluid density ($\rho_{\boldsymbol{\tau}}$) and the transverse shear modulus ($\mu_s$) dynamically decay at the identical inverse-square temporal limit, the kinematic speed of sound ($c$) evaluates mechanically as an absolute, invariant constant.
\begin{align}
    \rho_{\boldsymbol{\tau}}(t) &\propto t^{-2} \label{eq:density_decay} \\
    \mu_s(t) &\propto t^{-2} \label{eq:stiffness_decay} \\
    c &= \sqrt{\frac{\mu_s}{\rho_{\boldsymbol{\tau}}}} = \text{Constant}  \label{eq:c_invariant}
\end{align}
\end{tcolorbox}

\begin{tcolorbox}[cheatbox, title={\refstepcounter{cheatnum}\label{box:bao_lensing}\thecheatnum. Baryon Acoustic Oscillations \& Kinematic Unification \scriptsize\hfill\hyperref[sec:deriv_bao_lensing]{[Proof $\downarrow$]}}]
The divergence of the BAO standard ruler evaluates as lensed cosmic information propagating through the Gradient-Index (GRIN) vacuum. Longitudinal depth evaluates as the spatial derivative of the true coordinate distance. 
\begin{align}
    &\text{1. Longitudinal Depth Derivative:} \nonumber \\
    &\quad D_{H,true}(z) = \frac{R_{h,0}}{1+z} \label{eq:bao_longitudinal} \\
    &\text{2. Kinematic Unification:} \nonumber \\
    &\quad a_0(z) = \frac{c^2}{2\pi D_{H,true}(z)} \label{eq:bao_kinematic_unification}
\end{align}
\end{tcolorbox}

\begin{tcolorbox}[cheatbox, title={\refstepcounter{cheatnum}\label{box:auxetic_symmetry_covariance}\thecheatnum. Symmetry Breaking \& Covariance \scriptsize\hfill\hyperref[sec:deriv_auxetic_symmetry_covariance]{[Proof $\downarrow$]}}]
Adjacent topological circulations symmetrically break macroscopic fluid tension, establishing an interstitial stagnation plane. Opposing macroscopic pressure gradients mathematically annihilate ($\nabla P_{net} = 0$), while scalar Bernoulli pressure deficits algebraically superimpose ($\Delta P_{net}$). 
\begin{align}
    &\text{1. The Stagnation Plane (Lagrange Point):} \nonumber \\
    &\quad \nabla P_{net}(0) = 0 \implies \mathbf{a} = 0 \label{eq:gradient_cancellation} \\
    &\text{2. Hidden Action (Structural Strain):} \nonumber \\
    &\quad \mathbf{a} = 0, \quad \sigma_{ij} \gg 0 \label{eq:hidden_action_stress} \\
    &\text{3. Scalar Superposition (Pressure Deficit):} \nonumber \\
    &\quad \Delta P_{net} = P_{dyn,A} + P_{dyn,B} \label{eq:scalar_superposition} \\
    &\text{4. Gravitational Time Dilation:} \nonumber \\
    &\quad d\tau = dt_{bg} \sqrt{1 - \frac{\Delta P_{net}}{P_c}} \label{eq:scalar_time_dilation} \\
    &\text{5. Acoustic Covariance:} \nonumber \\
    &\quad c_{obs} = \frac{dr_{wave}}{d\tau} \equiv c \label{eq:acoustic_covariance}
\end{align}
\end{tcolorbox}

\begin{tcolorbox}[cheatbox, title={\refstepcounter{cheatnum}\label{box:mond}\thecheatnum. Entropic Superposition (MOND) \scriptsize\hfill\hyperref[sec:deriv_mond]{[Proof $\downarrow$]}}]
Vacuum field energy isolates flat galactic rotation curves natively. Flat rotation evaluates as the geometric superposition of localized baryonic mass and the isotropic cosmic background acceleration:
\begin{align}
    a_0(z) &= \frac{c^2(1+z)}{2\pi R_{h,0}} \approx 1.10 \times 10^{-10} \frac{\text{m}}{\text{s}^2} \label{eq:a0z} \\
    g_{eff}^2 &= \langle (\mathbf{g}_{bar} + \mathbf{a}_0)^2 \rangle = g_{bar} \cdot a_0(t) \label{eq:geff} \\
    v_{rot} &= (G M_{bar} \cdot a_0(t))^{1/4} \label{eq:vrot}
\end{align}
\end{tcolorbox}

\begin{tcolorbox}[cheatbox, title={\refstepcounter{cheatnum}\label{box:macroscopic_grin}\thecheatnum. Refractive Gravity \& Dark Matter (Macroscopic GRIN) \scriptsize\hfill\hyperref[sec:deriv_macroscopic_grin]{[Proof $\downarrow$]}}]
Macroscopic gravity and anomalous galactic rotation evaluate as continuous Gradient-Index (GRIN) optical lenses within the fluid metric.
\begin{align}
    &\text{1. Acoustic Limit:} \quad c(r) = \sqrt{\frac{\mu_s(r)}{\rho(r)}} \label{eq:acoustic_limit} \\
    &\text{2. Acoustic Gravity (Fermat's Principle):} \nonumber \\
    &\quad \nabla P \implies \nabla n \implies \text{Refraction} \label{eq:acoustic_gravity} \\
    &\text{3. Dark Matter (Macroscopic GRIN):} \nonumber \\
    &\quad n_{gal}(r) = \left( 1 - \frac{v_{rot}^2}{c^2} \right)^{-1} \label{eq:dark_matter_grin}
\end{align}
\end{tcolorbox}



\begin{tcolorbox}[cheatbox, title={\refstepcounter{cheatnum}\label{box:mass_amplification}\thecheatnum. Temporal Amplification of Mass \scriptsize\hfill\hyperref[sec:deriv_mass_amplification]{[Proof $\downarrow$]}}]
The effective bare rest mass evaluates as inversely proportional to the square root of the macroscopic scale radius.
\begin{equation}
    m(z) = m_0 \sqrt{1+z} \label{eq:mass_z}
\end{equation}
\end{tcolorbox}


% ==========================================
% PHASE II: VACUUM THERMODYNAMICS & RHEOLOGY
% ==========================================
\vspace{0.4cm}
\noindent\textbf{\large \color{academiabrown} PHASE II: VACUUM THERMODYNAMICS \& RHEOLOGY}
\vspace{0.2cm}

\begin{tcolorbox}[cheatbox, title={\refstepcounter{cheatnum}\label{box:thermo_vacuum}\thecheatnum. Absolute Thermodynamic Vacuum \& Entropy \scriptsize\hfill\hyperref[sec:deriv_thermo_vacuum]{[Proof $\downarrow$]}}]
Maintaining a stationary coordinate against accelerating background fluid exacts an aerodynamic heat tax. Via the Unruh-Hawking effect, proper geometric boundary acceleration generates a continuous thermodynamic floor.
\begin{align}
    &\text{1. The Unruh-Hawking Temperature Floor:} \nonumber \\
    &\quad a_0(t) = \frac{c}{2\pi t} \implies T_{cosmic}(t) = \frac{\hbar c}{2\pi k_B t} \label{eq:tcosmic} \\
    &\text{2. Particle Entropy Quantum:} \nonumber \\
    &\quad S_p = \frac{(\rho + P)V}{T} \implies S_p = \frac{8\pi}{3} k_B \label{eq:sp}
\end{align}
\end{tcolorbox}

\begin{tcolorbox}[cheatbox, title={\refstepcounter{cheatnum}\label{box:unified_eos}\thecheatnum. The Unified Equation of State \scriptsize\hfill\hyperref[sec:deriv_unified_eos]{[Proof $\downarrow$]}}]
The thermodynamic state of the continuum transitions flawlessly from a macroscopic null-action vacuum to a localized topological phase boundary.
\begin{align}
    &\text{1. Fundamental Vacuum Energy \& Null Action:} \nonumber \\
    &\quad E_{vac}(t) = \frac{\hbar}{t} \implies \mathcal{L}_{vac}(c) = 0 \label{eq:vacuum_null_action} \\
    &\text{2. The Topological Equation of State:} \nonumber \\
    &\quad PV = \frac{2\pi}{3} k_B T \implies U = 2\pi k_B T \label{eq:unified_topological_eos} \\
    &\text{3. Rest Mass Equivalence \& Bag Pressure:} \nonumber \\
    &\quad U_{topo} = mc^2 \implies P_{int} = 2\pi P_p \label{eq:unified_pressure_ratio}
\end{align}
\end{tcolorbox}

\begin{tcolorbox}[cheatbox, title={\refstepcounter{cheatnum}\label{box:rheological_gradient}\thecheatnum. The Continuous Rheological Gradient \scriptsize\hfill\hyperref[sec:deriv_rheological_gradient]{[Proof $\downarrow$]}}]
The macroscopic continuum evaluates as a dynamically shifting elastic medium governed by localized topological circulation velocity ($v$). The metric transitions continuously from an expanding auxetic vacuum ($\nu = -0.25$) to an absolute incompressible core ($\nu = 0.5$).
\begin{align}
    \lim_{v \to 0} \nu(r) = -0.25 &\implies \text{Auxetic Vacuum} \label{eq:poisson_macro} \\
    \nu(\lambda_c) \in [0, 0.5) &\implies \text{Elastic Boundary} \label{eq:poisson_compton} \\
    \lim_{v \to c} \nu(l_p) = 0.5 &\implies \text{Incompressible Core} \label{eq:poisson_local}
\end{align}
\end{tcolorbox}

\begin{tcolorbox}[cheatbox, title={\refstepcounter{cheatnum}\label{box:unified_deborah}\thecheatnum. The Unified Cosmological Deborah Number \scriptsize\hfill\hyperref[sec:deriv_unified_deborah]{[Proof $\downarrow$]}}]
The phase transition between yielding Eulerian fluid and rigid optical crystal evaluates via the Complex Dynamic Modulus ($G^*$). The Deborah Number evaluates as the geometric ratio of the interaction energy ($E$) to the fundamental continuum baseline ($m_e c^2$). 
\begin{align}
    De &= \alpha \left( \frac{E}{m_e c^2} \right) \label{eq:unified_deborah} \\
    G^*(E) &= \mu_s \left[ \frac{De^2}{1 + De^2} \right] + i \mu_s \left[ \frac{De}{1 + De^2} \right] \label{eq:deborah_modulus} \\
    \frac{\alpha_s}{\alpha} &= \tan \delta = \frac{m_e c^2}{\alpha E} \label{eq:angular_distribution}
\end{align}
\textit{Phase Symmetry \& Asymptotic Limits:}
\begin{align}
    &De \to 1 \implies 2 E_{crit} = 140.04 \text{ MeV} \equiv m_{\pi^{\pm}} \label{eq:de_glass} \\
    &\quad De = 1 \implies \alpha_s = \alpha \label{eq:electro_strong_symmetry} \\
    &\quad E = N m_e c^2 \implies \alpha_s = \frac{1}{N} \label{eq:volumetric_multiplier} \\
    &\quad \frac{d\alpha_s}{dN} = - \frac{1}{N^2} \label{eq:beta_state} \\
    &\quad \beta(\alpha_s) = \frac{\partial \alpha_s}{\partial \ln E} = - \alpha_s \label{eq:beta_function} \\
    &\quad De \ll 1 \implies L_{crit} = r_e \approx 2.81 \text{ fm} \label{eq:de_range} \\
    &\quad De \to 0 \implies De_{min} \approx 10^{-41} \label{eq:de_gravity}
\end{align}
\end{tcolorbox}

\begin{tcolorbox}[cheatbox, title={\refstepcounter{cheatnum}\label{box:viscoelastic_stiffening}\thecheatnum. Viscoelastic Phase Stiffening \& Time Dilation \scriptsize\hfill\hyperref[sec:deriv_viscoelastic_stiffening]{[Proof $\downarrow$]}}]
Relativistic time dilation, strong interaction coupling, and the cosmic speed limit evaluate as identical projections of dynamic fluid pressure consuming the structural compliance of the metric.
\begin{align}
    &\text{1. The Bernoulli Compliance Limit ($\alpha_s$):} \nonumber \\
    &\quad \alpha_s^2 = 1 - \frac{P_{dyn}}{P_c} \label{eq:admittance_pressure_ratio} \\
    &\text{2. The Yukawa-Bernoulli Identity:} \nonumber \\
    &\quad \alpha_s^2 = 1 - \frac{g^2}{g_{max}^2} \label{eq:yukawa_identity} \\
    &\text{3. Relativistic Packing \& The Lorentz Limit:} \nonumber \\
    &\quad \rho_{eff} = \rho_{\tau} \gamma \label{eq:relativistic_packing_admittance} \\
    &\quad \gamma = \sqrt{\frac{P_c}{P_c - P_{dyn}}} \label{eq:lorentz_pressure} \\
    &\text{4. The Kinematic Phase Lock:} \nonumber \\
    &\quad \alpha_s = \frac{v_{\circlearrowleft}}{c} \label{eq:admittance_spin_identity} \\
    &\text{5. Viscoelastic Time Dilation (The Deborah Limit):} \nonumber \\
    &\quad De = \alpha \gamma \label{eq:kinematic_deborah} \\
    &\text{6. Asymptotic Freedom \& Solidification:} \nonumber \\
    &\quad \lim_{P_{dyn} \to P_c} \gamma \to \infty \implies \alpha_s = 0, \quad \nabla \cdot \mathbf{v} = 0 \label{eq:cavitation_collapse}
\end{align}
\end{tcolorbox}


\begin{tcolorbox}[cheatbox, title={\refstepcounter{cheatnum}\label{box:universal_impedance}\thecheatnum. Universal Kinematic Vacuum Impedance  ($[\boldsymbol{Z}_{\tau}]$) \scriptsize\hfill\hyperref[sec:deriv_universal_impedance]{[Proof $\downarrow$]}}]
Inertial mass, optical density, quantum electrical resistance, and quantized spin evaluate as the diagonal geometric limits of a single continuous matrix: the Kinematic Impedance Tensor ($[\boldsymbol{Z}_{\tau}]$).
\begin{align}
    &\text{1. The Transverse Quantum Limit (von Klitzing):} \nonumber \\
    &\quad Z_{\tau(\perp)} = \frac{h}{q^2} \equiv R_K \label{eq:von_klitzing_impedance} \\
    &\text{2. The Longitudinal Inertial Limit (Effective Mass):} \nonumber \\
    &\quad Z_{\tau(\parallel)} = \hbar^2 \left( \frac{d^2 E}{dk^2} \right)^{-1} \equiv m^* \label{eq:effective_mass_impedance} \\
    &\text{3. The Transverse Limit (Refractive Gravity):} \nonumber \\
    &\quad Z_{\tau(\perp)} \propto n(r) = \left( 1 - \frac{v_0^2}{c^2} \right)^{-1} \label{eq:refractive_impedance} \\
    &\text{4. The Rotational Phase Limit (Gyroscopic Inertia):} \nonumber \\
    &\quad Z_{\tau(\circlearrowleft)} = \frac{E}{\omega} \equiv \hbar \label{eq:planck_impedance} \\
    &\lim_{|\nabla \times \mathbf{v}_{\psi}| \to c} Z_{\tau(\circlearrowleft)} \to \infty \implies \text{Spin Quantization } (\hbar/2) \label{eq:infinite_rotational_impedance}
\end{align}
\end{tcolorbox}

\newpage
% ==========================================
% PHASE III: TOPOLOGICAL MASS GENERATION
% ==========================================
\vspace{0.4cm}
\noindent\textbf{\large \color{academiabrown} PHASE III: TOPOLOGICAL MASS GENERATION}
\vspace{0.2cm}



\begin{tcolorbox}[cheatbox, title={\refstepcounter{cheatnum}\label{box:hydrodynamic_mass}\thecheatnum. The Kinematic Origin of Inertia \scriptsize\hfill\hyperref[sec:deriv_hydrodynamic_mass]{[Proof $\downarrow$]}}]
The observable rest mass ($m_\tau$) of a fundamental boundary evaluates as the hydrodynamic added mass displaced during kinematic translation. 
\begin{align}
    m_{intrinsic} &= \rho_{void} V_\tau = 0 \label{eq:intrinsic_mass} \\
    m_\tau &= m_{added} = \rho_{\tau} V_{added} \label{eq:kinematic_mass} \\
    m_\tau &= \frac{1}{2} \rho_{\tau} V_\tau \label{eq:mass_volume_equivalence}
\end{align}
\end{tcolorbox}

\begin{tcolorbox}[cheatbox, title={\refstepcounter{cheatnum}\label{box:ns_reg}\thecheatnum. Navier-Stokes Regularization (Topological Pinch-Off) \scriptsize\hfill\hyperref[sec:deriv_ns_reg]{[Proof $\downarrow$]}}]
The Navier-Stokes equations are regularized by the acoustic limit of the metric. The singularity is prevented by an absolute viscosity threshold.
\begin{align}
    &\text{1. Absolute Dynamic Viscosity ($\mu$):} \nonumber \\
    &\quad \mu = \frac{1}{2} \frac{\rho_{\tau} \hbar}{m} \label{eq:vacuum_viscosity} \\
    &\text{2. Prandtl Mixing Length Equivalence:} \nonumber \\
    &\quad \mu = \frac{1}{2} \rho_{\tau} c L \label{eq:prandtl_equivalence} \\
    &\text{3. The Critical Reynolds Number ($Re_c$):} \nonumber \\
    &\quad Re_c = \frac{\rho_{\tau} c L}{\mu} = 2 \label{eq:critical_reynolds}
\end{align}
\end{tcolorbox}

\begin{tcolorbox}[cheatbox, title={\refstepcounter{cheatnum}\label{box:fluid_inertia}\thecheatnum. Origin of Inertia (Strain Compactness \& Tension) \scriptsize\hfill\hyperref[sec:deriv_fluid_inertia]{[Proof $\downarrow$]}}]
Mass and inertia evaluate as the geometric tension generated by localized circulation.
\begin{align}
    &\text{1. Relativistic Strain Compactness:} \nonumber \\
    &\quad \rho_{eff}(r) = \rho_{\tau} \left( 1 - \frac{P_{dyn}(r)}{P_c} \right)^{-1/2} \label{eq:fluid_inertia_gradient} \\
    &\text{2. The Dimensional Tension Bridge:} \nonumber \\
    &\quad T_{1D} = \epsilon_{\tau} c^2 \implies P_c = \frac{1}{2} \frac{T_{1D}}{A_{yield}} \label{eq:dimensional_tension_bridge} \\
    &\text{3. The Isotropic Tension Identity:} \nonumber \\
    &\quad \mathbb{T}_{local} = \mathbb{T}_c + P_{dyn} \label{eq:tension_identity} \\
    &\text{4. The Geometric Equation of State (Pure Tension):} \nonumber \\
    &\quad \rho_{eff}(r) = \rho_{\tau} \left( 2 - \frac{\mathbb{T}_{local}(r)}{\mathbb{T}_c} \right)^{-1/2} \label{eq:tension_density} \\
    &\text{5. The Compliant Boundary:} \nonumber \\
    &\quad \rho_{eff}(\lambda_c) \approx \rho_{\tau} + \rho_{\tau} \left( \frac{P_{dyn}(\lambda_c)}{2P_c} \right) \label{eq:compliant_inertia}
\end{align}
\end{tcolorbox}

\begin{tcolorbox}[cheatbox, title={\refstepcounter{cheatnum}\label{box:volumetric_inverse_cube}\thecheatnum. The Volumetric Inverse-Cube Identity \scriptsize\hfill\hyperref[sec:deriv_volumetric_inverse_cube]{[Proof $\downarrow$]}}]
The phenomenological paradox of inverse mass-radius scaling evaluates as a strict volumetric equilibrium between the 3D Euclidean bounding box ($V_R$) and the displaced void volume ($V_\tau$). Substituting the hydrodynamic added mass into the kinematic radius ($R = \hbar/mc$) derives an absolute inverse-cube continuum identity. 
\begin{align}
    V_R &= \frac{4}{3}\pi \left( \frac{\hbar}{m_\tau c} \right)^3 \label{eq:bounding_volume} \\
    V_R \cdot V_{\tau}^3 &= \frac{32\pi}{3} \left( \frac{\hbar}{\rho_{\tau} c} \right)^3 \label{eq:volumetric_inverse_cube}
\end{align}
\end{tcolorbox}

\begin{tcolorbox}[cheatbox, title={\refstepcounter{cheatnum}\label{box:admittance_ns}\thecheatnum. Admittance-Modified Navier-Stokes \scriptsize\hfill\hyperref[sec:deriv_admittance_ns]{[Proof $\downarrow$]}}]
To evaluate continuous acceleration, the kinematic momentum of the metric evaluates against the localized Admittance Tensor ($[\boldsymbol{\alpha}]$), mathematically annihilating absolute mass to isolate scale-invariant geometric kinematics.
\begin{align}
    &\text{1. The Relativistic Density Tensor:} \nonumber \\
    &\quad \boldsymbol{\rho}_{eff} = \rho_{\tau} [\boldsymbol{\alpha}]^{-1} \label{eq:fluid_inertia_gradient} \\
    &\text{2. The Quantum Reynolds Limit (Hadronization):} \nonumber \\
    &\quad Re_q = \frac{2}{\alpha_s} \label{eq:quantum_reynolds} \\
    &\text{3. The Equivalence Principle:} \nonumber \\
    &\quad \rho_{\tau} \frac{D\mathbf{v}}{Dt} = -\rho_{\tau} \nabla \Phi + [\boldsymbol{\alpha}] \Sigma \mathbf{F}_{surface} \label{eq:admittance_ns} \\
    &\text{4. The 3D Tensor Generalization (Force Density):} \nonumber \\
    &\quad \rho_{\tau} \frac{D\mathbf{v}}{Dt} = -\rho_{\tau} \nabla \Phi \quad + \nonumber \\
    & \alpha_{\parallel}(-\nabla P) + \alpha_{\perp}(\nabla \cdot \mathbb{T}_{shear}) + \alpha_{\circlearrowleft}(\mu \nabla^2 \mathbf{v}) \label{eq:propulsion_master} \\
    &\text{5. The Absolute Kinematic Equation of Motion:} \nonumber \\
    &\quad \frac{D\mathbf{v}}{Dt} = \nabla \mathcal{P}_{ambient} - \alpha_{\parallel}\nabla \mathcal{P}_{local} \quad + \nonumber \\
    &\quad \alpha_{\perp}(\nabla \cdot \mathbb{T}_{kin}) + \alpha_{\circlearrowleft}(\nu \nabla^2 \mathbf{v}) \label{eq:kinematic_ns}
\end{align}
\end{tcolorbox}

\begin{tcolorbox}[cheatbox, title={\refstepcounter{cheatnum}\label{box:localized_metric_strain}\thecheatnum. Localized Metric (Conformal Duality) \scriptsize\hfill\hyperref[sec:deriv_localized_metric_strain]{[Proof $\downarrow$]}}]
Microscopic velocity field evaluates as the conformal geometric inverse of the cosmos. Its spatial differential evaluates as the volumetric strain, gravitational redshift, and localized inertia.
\begin{align}
    &\text{1. Conformal Inversion (Microscopic velocity):} \nonumber \\
    &\quad v_{obs}(r) = c \cdot e^{-r/R_{h,p}} \label{eq:vobs_conformal} \\
    &\text{2. The Strain-Redshift Identity:} \nonumber \\
    &\quad \varepsilon(r) = 1+z(r) = \frac{1}{\sqrt{1 - e^{-2r/R_{h,p}}}} \label{eq:strain_redshift_identity} \\
    &\text{3. Microscopic Conformal Mirror:} \nonumber \\
    &\quad r_{true, micro}(z) = R_{h,p} \ln\left( 1 + \frac{1}{z_{local}} \right) \label{eq:conformal_r} \\
    &\text{4. The Integration of Inertia ($m_\tau$):} \nonumber \\
    &\quad m_\tau = \frac{1}{2} \rho_{\tau} \int_{l_p}^\infty 4\pi r^2 z(r) \, dr \label{eq:mass_integral_redshift} \\
    &\text{5. Topological Magnification (The Quantum Bit):} \nonumber \\
    &\quad Mag = \frac{R_{h,p}}{l_p} = \frac{m_p}{m_0} \implies m_0 = \frac{m_p}{Mag} \label{eq:conformal_mag} \\
    &\text{6. The Planck Core Energy Limit ($r \to l_p$):} \nonumber \\
    &\quad \lim_{r \to l_p} E(r) \approx E_0 \sqrt{\frac{R_{h,p}}{2l_p}} \implies \text{Kinematic Lockout} \label{eq:planck_core_energy_limit}
\end{align}
\end{tcolorbox}

\begin{tcolorbox}[cheatbox, title={\refstepcounter{cheatnum}\label{box:bag_pressure}\thecheatnum. Confinement Pressure \& Work \scriptsize\hfill\hyperref[sec:deriv_bag_pressure]{[Proof $\downarrow$]}}]
The geometric centripetal boundary condition ($a_p = c/2\pi t_{\text{loc}}$) rigorously dictates the inward Bag Pressure:
\begin{equation}\label{eq:pp}
    P_p = \frac{m^4 c^5}{8\pi^2 \hbar^3} 
\end{equation}

\end{tcolorbox}



\begin{tcolorbox}[cheatbox, title={\refstepcounter{cheatnum}\label{box:unified_engine}\thecheatnum. The Unified Topological Engine \scriptsize\hfill\hyperref[sec:deriv_unified_engine]{[Proof $\downarrow$]}}]
Mass evaluates as a continuous dipole, constrained by volumetric flow and topological yield limits.
\begin{align}
&\text{1. Volumetric Continuity (The Transducer):} \nonumber \\
    &\quad \nabla \cdot \mathbf{v}_{\circlearrowleft} = 0 \label{eq:engine_transduction} \\
    &\text{2. Hydrostatic Supremacy (Phase Transform):} \nonumber \\
    &\quad \frac{P_{out}}{P_{in}} = \frac{1}{P_{in}} \oint_{\partial S} P_{in} \, d\theta = 2\pi \label{eq:engine_supremacy} \\
    &\text{3. The Planck Redline (Kinematic Lockout):} \nonumber \\
    &\quad \lim_{V_R \to l_p^3} v_{intake} \to c \implies m_{max} = m_p \label{eq:engine_redline} \\
    &\text{4. Hydrostatic Instability (Particle Decay):} \nonumber \\
    &\quad \frac{V_\tau}{V_R} = \left(\frac{3c^3}{2\pi\rho_{\tau}\hbar^3}\right) m_\tau^4 \label{eq:engine_decay}
\end{align}
\end{tcolorbox}

\begin{tcolorbox}[cheatbox, title={\refstepcounter{cheatnum}\label{box:quantum_spin}\thecheatnum. Fluid Circulation \& The Topological Origin of Spin \scriptsize\hfill\hyperref[sec:deriv_quantum_spin]{[Proof $\downarrow$]}}]
Fundamental matter evaluates geometrically as closed fluid vorticity ($\nabla \times \mathbf{u} \neq 0$). Quantum Spin evaluates as the macroscopic fluid circulation of this continuous metric vortex.
\begin{align}
    &\text{1. Kinematic Angular Momentum \& Spin Limit:} \nonumber \\
    &\quad \mathbf{S} = \iiint_V \rho (\mathbf{r} \times \mathbf{v}) \, dV \implies S = \hbar \label{eq:transverse_shear_spin} \\
    &\text{2. The Baseline Circulation Limit:} \nonumber \\
    &\quad \Gamma = \oint_C \mathbf{v} \cdot d\mathbf{l} = \frac{h}{m} \label{eq:fluid_circulation} \\
    &\text{3. Spinor Topology \& Fermionic Spin ($4\pi$ Twist):} \nonumber \\
    &\quad S_{fermion} = m \left( \frac{\Gamma}{4\pi} \right) = \frac{\hbar}{2} \label{eq:spinor_topology}
\end{align}
\end{tcolorbox}

\begin{tcolorbox}[cheatbox, title={\refstepcounter{cheatnum}\label{box:volumetric_energy}\thecheatnum. Kinetic-Potential Strain Conservation \scriptsize\hfill\hyperref[sec:deriv_volumetric_energy]{[Proof $\downarrow$]}}]
Absolute energy evaluates as the displaced fluid volume of the continuous metric. Potential and kinetic energies evaluate as physical partitions of the localized volumetric strain ($\alpha_s$).
\begin{align}
    &\text{1. Kinematic Drag Factor:} \nonumber \\
    &\quad \gamma = \frac{1}{\alpha_s} \label{eq:strain_lorentz} \\
    &\text{2. Kinetic Displacement Volume:} \nonumber \\
    &\quad V_{kin} = \left( \frac{1}{\alpha_s} - 1 \right) V_{added} \label{eq:kinetic_volume} \\
    &\text{3. Potential Displacement Volume:} \nonumber \\
    &\quad V_{pot} = \frac{2}{\rho_{\tau} c^2} \int_{x_1}^{x_2} \mathbf{F} \cdot d\mathbf{x} \label{eq:potential_volume} \\
    &\text{4. Isochoric Transmutation:} \nonumber \\
    &\quad \dot{V}_{kin} = -\dot{V}_{pot} \label{eq:isochoric_transmutation}
\end{align}
\end{tcolorbox}

\newpage
% ==========================================
% PHASE IV: WAVE KINEMATICS & ELECTROMAGNETISM
% ==========================================
\vspace{0.4cm}
\noindent\textbf{\large \color{academiabrown} PHASE IV: WAVE KINEMATICS \& ELECTROMAGNETISM}
\vspace{0.2cm}

\begin{tcolorbox}[cheatbox, title={\refstepcounter{cheatnum}\label{box:unified_lagrangian}\thecheatnum. The Unified Fluid Lagrangian \& Modal Admittance \scriptsize\hfill\hyperref[sec:deriv_unified_lagrangian]{[Proof $\downarrow$]}}]
The kinematics of the metric evaluate via Modal Admittance.
\begin{align}
    &\text{1. Kinematic Capacity:} \nonumber \\
    &\quad v_{tot}^2 = (\nabla h)^2 + (\nabla \times \mathbf{A})^2 + (\nabla \times \mathbf{v}_{\psi})^2 \le c^2 \label{eq:absolute_pythagorean_constraint} \\
    &\text{2. Modal Geometric Admittance ($\alpha_{i}$):} \nonumber \\
    &\quad \alpha_{\parallel} = \sqrt{1 - \frac{(\nabla h)^2}{c^2}}, \quad \alpha_{\perp} = \sqrt{1 - \frac{(\nabla \times \mathbf{A})^2}{c^2}}, \nonumber \\ 
    &\quad \alpha_{\circlearrowleft} = \sqrt{1 - \frac{(\nabla \times \mathbf{v}_{\psi})^2}{c^2}} \label{eq:modal_admittance} \\
    &\text{3. Topological Confinement Limit:} \nonumber \\
    &\quad |\nabla \times \mathbf{v}_{\psi}| \to c \implies \alpha_{\circlearrowleft} \to 0 \implies De_{rot} \to \infty \label{eq:deborah_rotational} \\
    &\text{4. The Localized Compliance Limit:} \nonumber \\
    &\quad \Sigma \alpha_i^2 = 3 - \frac{v_{tot}^2}{c^2} \label{eq:compliance_sum} \\
    &\quad v_{tot} \to c \implies \Sigma \alpha_i^2 = 2 \label{eq:absolute_compliance_state} \\
    &\text{5. Dynamic \& Kinematic Lagrangians:} \nonumber \\
    &\quad \mathcal{L}_{dyn} = \frac{1}{2} \rho_{\tau} \Big[ \alpha_{\parallel}(\nabla h)^2 + \alpha_{\perp}(\nabla \times \mathbf{A})^2 + \nonumber \\
    &\quad \alpha_{\circlearrowleft}(\nabla \times \mathbf{v}_{\psi})^2 \Big] - \frac{1}{2} \rho_{\tau} c^2 \label{eq:modal_lagrangian} \\
    &\quad \mathcal{L}_{kin} = \frac{1}{2} \Big[ \alpha_{\parallel}(\nabla h)^2 + \alpha_{\perp}(\nabla \times \mathbf{A})^2 + \nonumber \\
    &\quad \alpha_{\circlearrowleft}(\nabla \times \mathbf{v}_{\psi})^2 \Big] - \frac{1}{2} c^2 \label{eq:kinematic_lagrangian} \\
    &\text{6. Absolute Angular Phase-Space ($\mathcal{L}_{kin}$):} \nonumber \\
    &\mathcal{L}_{kin}(v_{tot}, \theta, \phi) = \frac{1}{2} v_{tot}^2 \Big[ \sqrt{1 - \frac{v_{tot}^2}{c^2} \cos^2(\theta)} \cos^2(\theta) \nonumber \\
    &\quad + \sqrt{1 - \frac{v_{tot}^2}{c^2} \sin^2(\theta)\cos^2(\phi)} \sin^2(\theta)\cos^2(\phi) \nonumber \\
    &\quad + \sqrt{1 - \frac{v_{tot}^2}{c^2} \sin^2(\theta)\sin^2(\phi)} \sin^2(\theta)\sin^2(\phi) \Big] \nonumber \\
    &\quad - \frac{1}{2} c^2 \label{eq:angular_lagrangian}
\end{align}
\end{tcolorbox}

\begin{tcolorbox}[cheatbox, title={\refstepcounter{cheatnum}\label{box:legendre_hamiltonian}\thecheatnum. The Legendre Hamiltonian Translation \scriptsize\hfill\hyperref[sec:deriv_legendre_hamiltonian]{[Proof $\downarrow$]}}]
The macroscopic Legendre transformation translates continuous kinematic phase-space into stationary observable energy densities and topological mass.
\begin{align}
    &\text{1. Macroscopic Conjugate Momentum:} \nonumber \\
    &\quad \mathbf{p}_{macro} = \rho_{\tau} \mathbf{v}_{macro} \label{eq:conjugate_momentum} \\
    &\text{2. The Baseline Legendre Transform:} \nonumber \\
    &\quad \mathcal{H}_{\boldsymbol{\tau}} = \mathcal{K}_{macro} + \mathcal{P} \label{eq:legendre_transform} \\
    &\text{3. The Covariant Bernoulli State:} \nonumber \\
    &\quad \mathcal{H}_{\boldsymbol{\tau}} = \frac{1}{2} \rho_{\tau} \Big[ \mathbf{v}_{macro}^2 + c^2 \Big] \label{eq:bernoulli_hamiltonian} \\
    &\text{4. Dynamic \& Kinematic Hamiltonians:} \nonumber \\
    &\quad \mathcal{H}_{dyn} = \frac{1}{2} \rho_{\tau} \Big[ \alpha_{\parallel}(\nabla h)^2 + \alpha_{\perp}(\nabla \times \mathbf{A})^2 \nonumber \\
    &\quad \quad + \alpha_{\circlearrowleft}(\nabla \times \mathbf{v}_{\psi})^2 \Big] + \frac{1}{2} \rho_{\tau} c^2 \label{eq:modal_hamiltonian} \\
    &\quad \mathcal{H}_{kin} = \frac{1}{2} \Big[ \alpha_{\parallel}(\nabla h)^2 + \alpha_{\perp}(\nabla \times \mathbf{A})^2 \nonumber \\
    &\quad \quad + \alpha_{\circlearrowleft}(\nabla \times \mathbf{v}_{\psi})^2 \Big] + \frac{1}{2} c^2 \label{eq:kinematic_hamiltonian} \\
    &\text{5. The Rotational Saturation Limit:} \nonumber \\
    &|\nabla \times \mathbf{v}_{\psi}| \to c \implies \alpha_{\circlearrowleft} \to 0 \implies \mathcal{H}_{rest} = \frac{1}{2}\rho_{\tau}c^2 \label{eq:hamiltonian_rest_mass} \\
    &\text{6. Absolute Angular Phase-Space ($\mathcal{H}_{kin}$):} \nonumber \\
    &\quad \mathcal{H}_{kin}(v_{tot}, \theta, \phi) = \frac{1}{2} v_{tot}^2 \Big[ \alpha_{\parallel} \cos^2(\theta) \nonumber \\
    &+ \alpha_{\perp} \sin^2(\theta)\cos^2(\phi)  + \alpha_{\circlearrowleft} \sin^2(\theta)\sin^2(\phi) \Big] + \frac{1}{2} c^2 \label{eq:angular_hamiltonian} \\
    &\text{7. Topological Mass Extraction:} \nonumber \\
    &\quad E_{rest} = \int \mathcal{H}_{rest} \, dV = \frac{1}{2} \rho_{\tau} c^2 V_{\psi} \label{eq:rest_energy_integral} \\
    &\quad m_0 = \frac{1}{2} \rho_{\tau} V_{\psi} \label{eq:topological_mass}
\end{align}
\end{tcolorbox}

\begin{tcolorbox}[cheatbox, title={\refstepcounter{cheatnum}\label{box:dirac_isomorphism}\thecheatnum. Vorticity Equilibrium (The Dirac Equation) \scriptsize\hfill\hyperref[sec:deriv_dirac_isomorphism]{[Proof $\downarrow$]}}]
Fermion kinematics evaluate as gyroscopically stabilized acoustic vortexes. The algebraic operators of the Dirac sector evaluate as geometric fluid identities.
\begin{align}
    &\text{1. The Spinor (Cartan Tetrad):} \nonumber \\
    &\quad \psi \equiv \mathbf{\hat{e}}_i \label{eq:spinor_identity} \\
    &\text{2. The Kinematic Curl ($\gamma$-Matrices):} \nonumber \\
    &\quad i\gamma^\mu \partial_\mu \equiv \nabla \times \mathbf{v}_\psi \label{eq:gamma_curl} \\
    &\text{3. Geometric Mass Coupling (Absolute Density):} \nonumber \\
    &\quad \bar{\psi} \psi = \frac{1}{V_{\tau}} \implies \mathcal{L}_{mass} = - \frac{1}{2} \rho_{\tau} c^2 \label{eq:absolute_mass_lagrangian} \\
    &\text{4. The Fluid Dirac Lagrangian:} \nonumber \\
    &\quad \mathcal{L}_{Dirac} \equiv \frac{1}{2}\rho_{\tau} |\nabla \times \mathbf{v}_\psi|^2 - \frac{1}{2} \rho_{\tau} c^2 \bar{\psi}\psi \label{eq:dirac_fluid_lagrangian} \\
    &\text{5. Antimatter (Chiral Annihilation):} \nonumber \\
    &\quad \bar{\psi}\psi \equiv (W_+ + W_-) = 0 \label{eq:chiral_annihilate}
\end{align}
\end{tcolorbox}

\begin{tcolorbox}[cheatbox, title={\refstepcounter{cheatnum}\label{box:wave_mechanics}\thecheatnum. Continuous Wave Mechanics \& Topological Phase States \scriptsize\hfill\hyperref[sec:deriv_wave_mechanics]{[Proof $\downarrow$]}}]
The Kinematic Lagrangian maps the mechanical wave spectrum, discrete particle identities, and nuclear decay as geometric constraints of the continuous vacuum.
\begin{align}
    &\text{1. The Continuous Wave Solutions:} \nonumber \\
    &\quad \mathbf{v}_{\parallel}(r,t) = c \cos(\theta) \cos(\mathbf{k} \cdot \mathbf{r} - \omega t) \label{eq:wave_p} \\
    &\quad \mathbf{v}_{\perp}(r,t) = c \sin(\theta)\cos(\phi) \cos(\mathbf{k} \cdot \mathbf{r} - \omega t) \label{eq:wave_s} \\
    &\quad \mathbf{v}_{\circlearrowleft}(r,t) = c \sin(\theta)\sin(\phi) \cos(\mathbf{k} \cdot \mathbf{r} - \omega t) \label{eq:wave_rot} \\
    &\quad \Phi(r,t) = - \frac{(\mathbf{v}_{\parallel})^2}{2} \label{eq:p_wave_potential} \\
    &\text{2. Kinematics \& The Lorentz Angle ($\theta$):} \nonumber \\
    &\quad \sin(\theta) = \sqrt{1 - \frac{v_{macro}^2}{c^2}} \label{eq:lorentz_angle} \\
    &\text{3. Topological Charge Projection ($\phi$):} \nonumber \\
    &\quad |Q| = \cos^2(\phi) \label{eq:charge_projection} \\
    &\text{4. The Dynamic Neutrino Limit ($v_{macro} \to c$):} \nonumber \\
    &\quad \theta \to \pi, 0 \implies \mathcal{L}_{kin} \to 0 \label{eq:phase_neutrino} \\
    &\text{5. Antimatter \& Annihilation Rupture:} \nonumber \\
    &\quad (\nabla \times \mathbf{v}_{\psi}^-)^2 = c^2 \label{eq:antimatter_state} \\
    &\quad \sum \mathbf{v}_{\psi} = 0 \implies \alpha_{\circlearrowleft} = 1 \label{eq:annihilation_limit} \\
    &\text{6. Mechanical Yield ($\beta^-$ Decay):} \nonumber \\
    &\quad \mathcal{L}_{kin, n}(udd) \to \mathcal{L}_{kin, p}(uud) + \mathcal{L}_{kin, e}(\phi=\pi) \nonumber \\
    &\quad \quad + \mathcal{L}_{kin, \bar{\nu}}(\theta \to 0) \label{eq:neutron_rupture}
\end{align}
\end{tcolorbox}





\begin{tcolorbox}[cheatbox, title={\refstepcounter{cheatnum}\label{box:em_lc_isomorphism}\thecheatnum. EM Kinematics \& The LC-Acoustic Isomorphism \scriptsize\hfill\hyperref[sec:deriv_em_lc_isomorphism]{[Proof $\downarrow$]}}]
Electromagnetic field equations and electromechanical properties evaluate as the continuous 1D to 3D kinematic projections of the viscoelastic metric fluid.
\begin{align}
    &\text{1. Magnetic Vorticity \& Electric Strain:} \nonumber \\
    &\quad \nabla \cdot \mathbf{B} = 0 \quad \text{\&} \quad \nabla \times \mathbf{E} = -\frac{\partial \mathbf{B}}{\partial t} \label{eq:maxwell_gauss_faraday} \\
    &\text{2. Volumetric Strain Displacement (Gauss's Law):} \nonumber \\
    &\quad \nabla \cdot \mathbf{E} = \frac{\rho_e}{\epsilon_{\tau}} \label{eq:maxwell_gaussE} \\
    &\text{3. Macroscopic Modulus \& Density Mapping:} \nonumber \\
    &\quad \rho_{\tau} = \frac{\epsilon_{\tau}}{A_{yield}}, \quad K = \frac{1}{\mu_{\tau} A_{yield}} \label{eq:lc_constants} \\
    &\text{4. The Isomorphic Circuit Elements:} \nonumber \\
    &\quad m_\tau = \frac{1}{2} \epsilon_{\tau} L_\tau \label{eq:lc_mass} \\
    &\quad E_{cap} = E_{kin} = \frac{1}{2} m_{\tau} v_{\perp}^2 \implies P_{dyn} = \frac{1}{2}\rho_{\tau} v_{\perp}^2 \label{eq:velocity_pressure} \\
    &\quad E_{ind} = \frac{1}{2} L_H I^2 = \frac{1}{2} k x^2 \label{eq:lc_inductance} \\
    &\text{5. The 1D Electromagnetic Projection (Ampère):} \nonumber \\
    &\quad \nabla \times \mathbf{B} = \mu_{\tau} \mathbf{J}_{EM} + \mu_{\tau} \epsilon_{\tau} \frac{\partial \mathbf{E}}{\partial t} \label{eq:maxwell_curlB} \\
    &\text{6. Kinematic Energy Flux (The Poynting Vector):} \nonumber \\
    &\quad \mathbf{S} = \frac{1}{\mu_{\tau}} (\mathbf{E} \times \mathbf{B}) \label{eq:poynting_flux}
\end{align}
\end{tcolorbox}

\begin{tcolorbox}[cheatbox, title={\refstepcounter{cheatnum}\label{box:continuum_constants}\thecheatnum. Continuum Constants (UV Regularization) \scriptsize\hfill\hyperref[sec:deriv_continuum_constants]{[Proof $\downarrow$]}}]
Mathematical divergence ($p \to \infty$) is bounded by the macroscopic yield strength of the metric. The constants of the vacuum evaluate as geometric fractions of the Planck scale.
\begin{align}
    &\text{1. Energy Ceiling (The Planck Limit):} \nonumber \\
    &\quad P_{dyn} \ge \sigma_{max} \implies E_{max} \equiv E_{Planck} \label{eq:planck_ceiling} \\
    &\text{2. Gravitational Yield Area ($A_{yield}$):} \nonumber \\
    &\quad A_{yield} = \frac{4\pi}{9} \frac{\epsilon_{\tau} G^2 \hbar}{c^5} \label{eq:gravitational_yield} \\
    &\text{3. 1D Metric Stiffness ($\epsilon_{\tau}$):} \nonumber \\
    &\epsilon_{\tau} = \frac{729 \alpha^3}{28} \left( \frac{m_{pl}}{m_{pr}} \right)^2 \frac{c^2}{G} \approx 2.30 \times 10^{60} \text{ kg/m} \label{eq:epsilon_numerical} \\
    &\text{4. 3D Volumetric Density ($\rho_{\tau}$):} \nonumber \\
    &\quad \rho_{\tau} = \frac{9}{4\pi} \rho_{pl} \approx 3.69 \times 10^{96} \text{ kg/m}^3 \label{eq:rho_tau_numerical}
\end{align}
\end{tcolorbox}

\begin{tcolorbox}[cheatbox, title={\refstepcounter{cheatnum}\label{box:fine_structure}\thecheatnum. The Topological Coupling Limit ($\alpha_{\circlearrowleft \perp}$) \& Charge \scriptsize\hfill\hyperref[sec:deriv_fine_structure]{[Proof $\downarrow$]}}]
Electric charge evaluates mechanically as the transverse geometric leakage of internal topological rotation, throttled by the off-diagonal coupling limit ($\alpha_{\circlearrowleft \perp}$).
\begin{align}
    &\text{1. Transverse Shear (Electromagnetic Force):} \nonumber \\
    &\quad F_{\circlearrowleft}(r) = \frac{\hbar c}{r^2} \label{eq:internal_rotational_force} \\
    &\quad F_{EM}(r) = \alpha_{\circlearrowleft \perp} \frac{\hbar c}{r^2} \label{eq:electromagnetic_leakage} \\
    &\text{2. The Continuous Isolation of Charge ($e$):} \nonumber \\
    &\quad e = \sqrt{4\pi \epsilon_0 \hbar c \alpha_{\circlearrowleft \perp}} \label{eq:alpha_continuum_derived} \\
    &\text{3. Macroscopic vs. Topological Impedance:} \nonumber \\
    &\quad Z_{macro} = \rho_{\tau} c \equiv Z_0 \label{eq:macroscopic_impedance} \\
    &\quad \alpha_{\circlearrowleft \perp} = \frac{1}{2} \left( \frac{Z_{macro}}{Z_{topo}} \right) \label{eq:tensor_coupling_derivation} \\
    &\text{4. The Absolute Topological Impedance ($Z_{topo}$):} \nonumber \\
    &\quad Z_{topo} = \frac{\rho_{\tau} c}{2 \alpha_{\circlearrowleft \perp}} \equiv R_K \label{eq:topological_impedance}
\end{align}
\end{tcolorbox}

\begin{tcolorbox}[cheatbox, title={\refstepcounter{cheatnum}\label{box:tensor_projections}\thecheatnum. 1D Fractional Phase Flux (Dimensional Projections) \scriptsize\hfill\hyperref[sec:deriv_tensor_projections]{[Proof $\downarrow$]}}]
Fractional subatomic charges ($+2/3, -1/3$) evaluate exclusively as the continuous geometric tensor artifacts of projecting an isotropic 3D fluid circulation onto a 1D observation axis.
\begin{align}
    &\text{1. The Longitudinal Projection (1/3 Flux):} \nonumber \\
    &\quad \Phi_{\parallel} = \Phi_{\circlearrowleft} \left( \frac{1}{4\pi} \iint_{\Omega} \cos^2 \theta \, d\Omega \right) = \frac{1}{3} \Phi_{\circlearrowleft} \label{eq:tensor_long} \\
    &\text{2. The Transverse Projection (2/3 Flux):} \nonumber \\
    &\quad \Phi_{\perp} = \Phi_{\circlearrowleft} \left( \frac{1}{4\pi} \iint_{\Omega} \sin^2 \theta \, d\Omega \right) = \frac{2}{3} \Phi_{\circlearrowleft} \label{eq:tensor_trans} \\
    &\text{3. Absolute Flux Conservation:} \nonumber \\
    &\quad \Phi_{\parallel} + \Phi_{\perp} = \Phi_{\circlearrowleft} \label{eq:projection_conservation} \\
    &\text{4. The 1D Fractional Phase Flux ($\Phi_0$):} \nonumber \\
    &\quad \Phi_0 = \frac{1}{3}\Phi_{\circlearrowleft} = \frac{e\sqrt{2A_{yield}}}{3\epsilon_{\tau}} \label{eq:1d_phase_flux}
\end{align}
\end{tcolorbox}

\begin{tcolorbox}[cheatbox, title={\refstepcounter{cheatnum}\label{box:kinematic_charge}\thecheatnum. The Kinematic Origin of Charge \& $g-2$ \scriptsize\hfill\hyperref[sec:deriv_kinematic_charge]{[Proof $\downarrow$]}}]
Electromagnetism evaluates as continuous geometric fluid shear across the viscoelastic vacuum. Charge evaluates as kinematic momentum.
\begin{align}
    &\text{1. The Absolute Shear Transduction Limit:} \nonumber \\
    &\quad \frac{1}{2}\rho_{\tau} v_{\circlearrowleft}^2 = \frac{1}{2}\epsilon_{\tau} (\nabla \times \mathbf{A}_{\perp})^2 \label{eq:shear_transduction_limit} \\
    &\text{2. The Unitary Strain Displacement (Charge):} \nonumber \\
    &\quad e_{\tau} = \int_S \epsilon_{\tau} (\nabla \times \mathbf{A}_{\perp}) \cdot d\mathbf{A} \label{eq:charge_displacement} \\
    &\text{3. Absolute Continuum Charge ($e_{\tau}$):} \nonumber \\
    &\quad e_{\tau} = \sqrt{4\pi \epsilon_{\tau} \hbar c \alpha_{\circlearrowleft \perp}} \approx 8.16 \times 10^{16} \text{ kg}\cdot\text{m/s} \label{eq:absolute_charge_derived} \\
    &\text{4. The Absolute Dirac Moment ($\mu_{B\tau}$):} \nonumber \\
    &\quad \mu_{B\tau} = \frac{e_{\tau} \hbar}{2m} \label{eq:bohr_magneton} \\
    &\text{5. Toroidal Phase Spill \& True Magnetic Moment:} \nonumber \\
    &\quad a_e = \frac{\alpha_{\circlearrowleft \perp}}{2\pi} \implies \mu_{true} = \mu_{B\tau} \left( 1 + \frac{\alpha_{\circlearrowleft \perp}}{2\pi} \right) \label{eq:true_magnetic_moment}
\end{align}
\end{tcolorbox}

\begin{tcolorbox}[cheatbox, title={\refstepcounter{cheatnum}\label{box:coulomb_bernoulli}\thecheatnum. Coulomb's Law (Spherical Bernoulli \& Phase Flux) \scriptsize\hfill\hyperref[sec:deriv_coulomb_bernoulli]{[Proof $\downarrow$]}}]
The electrostatic force evaluates as a macroscopic Bernoulli pressure gradient throttled by the off-diagonal topological coupling limit ($\alpha_{\circlearrowleft \perp}$). The elementary transverse charge evaluates as pure, unadulterated Planck geometry ($\Phi_{\perp(e)}$).
\begin{align}
    &\text{1. Transverse Bernoulli Stagnation:} \nonumber \\
    &\text{Attract:  } \mathbf{v}_{\perp(1)} \cdot \mathbf{v}_{\perp(2)} > 0 \implies P_{static} \downarrow \label{eq:attraction} \\
    &\text{Repel:  } \mathbf{v}_{\perp(1)} \cdot \mathbf{v}_{\perp(2)} < 0 \implies P_{static} \uparrow \label{eq:repulsion} \\
    &\text{2. The Macroscopic Bernoulli Force ($\mathbf{F}_{fluid}$):} \nonumber \\
    &\quad \mathbf{F}_{fluid} = - \frac{\alpha_{\circlearrowleft \perp} \rho_{\tau} \Phi_{\perp(1)} \Phi_{\perp(2)}}{8\pi r^2} \mathbf{\hat{r}} \label{eq:bernoulli_force} \\
    &\text{3. Fine-Structure Force Equivalence ($\alpha_{\circlearrowleft \perp} \equiv \alpha$):} \nonumber \\
    &\quad \frac{\alpha_{\circlearrowleft \perp} \rho_{\tau} \Phi_{\perp(e)}^2}{8\pi r^2} = \frac{\alpha \hbar c}{r^2} \label{eq:force_equivalence} \\
    &\text{4. Quantum Volumetric Phase Flux ($\Phi_{\perp(e)}$):} \nonumber \\
        &\quad \Phi_{\perp(e)} = \frac{4\sqrt{2}\pi}{3}c \space l_p^2 \approx 4.63 \times 10^{-61} \text{ m}^3/\text{s} \label{eq:charge_flux}
\end{align}
\end{tcolorbox}


\newpage

% ==========================================
% PHASE V: TENSOR TRANSLATION & FORCE UNIFICATION
% ==========================================
\vspace{0.4cm}
\noindent\textbf{\large \color{academiabrown} PHASE V: TENSOR TRANSLATION \& FORCE UNIFICATION}
\vspace{0.2cm}



% ----------------------------------------------------------------------
\begin{tcolorbox}[cheatbox, title={\refstepcounter{cheatnum}\label{box:rosetta_stone}\thecheatnum. Tensor Evolution \& Geometry\scriptsize\hfill\hyperref[sec:deriv_rosetta_stone]{[Proof $\downarrow$]}}]
The \textbf{Metric Tensor} ($g_{\mu\nu}$) evaluates as the fluid compliance matrix ($-\alpha_s^2$). The \textbf{Covariant Derivative} ($\nabla_\mu$) evaluates as the \textbf{Eulerian Material Derivative} ($D/Dt$). \textbf{Christoffel symbols} ($\Gamma^\lambda_{\mu\nu}$) evaluate as Bernoulli pressure gradients ($\nabla P$) and fluid velocity gradients ($\partial_j v_i$). The \textbf{Riemann Tensor} ($R_{ijkl}$) evaluates as Saint-Venant Strain Incompatibility ($\text{Inc}(\boldsymbol{\varepsilon})$). The \textbf{Ricci} and \textbf{Weyl} tensors evaluate as volumetric dilatation and deviatoric shear. The \textbf{Stress-Energy Tensor} ($T_{\mu\nu}$) evaluates as the Cauchy Stress Tensor combined with convective momentum flux. The \textbf{Einstein Tensor} ($G_{\mu\nu}$) evaluates as mechanical stress/strain equilibrium. The \textbf{Faraday Tensor} ($F_{\mu\nu}$) evaluates as macroscopic Kinematic Vorticity ($\Omega_{\mu\nu}$). \textbf{Yang-Mills} non-Abelian self-interactions evaluate as convective fluid acceleration ($\mathbf{u} \cdot \nabla \mathbf{u}$). \textbf{Spinors} evaluate as the orthogonal rotation of Cartan moving frames.
\begin{align}
    &\text{1. Acoustic Metric (Fluid Compliance):} \nonumber \\
    &\quad [\boldsymbol{g}_{\mu\nu}] = \text{diag}(-\alpha_s^2, \, \alpha_s^{-2}, \, r^2, \, r^2\sin^2\theta) 
    \label{eq:rosetta_metric_matrix} \\
    &\text{2. Material Evolution:} \nonumber \\
    &\quad \nabla_\mu \equiv \frac{D}{Dt} = \frac{\partial}{\partial t} + (\mathbf{v} \cdot \nabla) \label{eq:rosetta_covariant} \\
    &\text{3. Christoffel (Pressure Gradient \& Velocity Shear):} \nonumber \\
    &\quad c^2 \Gamma^i_{00} \equiv -\frac{1}{\rho_{\tau}} \nabla^i P_{static} \label{eq:rosetta_christoffel_pressure} \\
    &\quad \Gamma^k_{ij} \equiv \partial_j v_i \label{eq:rosetta_christoffel} \\
    &\text{4. Riemann Curvature (Strain Incompatibility):} \nonumber \\
    &\quad R_{ijkl} \equiv \text{Inc}(\boldsymbol{\varepsilon})_{ijkl} \label{eq:rosetta_riemann} \\
    &\text{5. Ricci \& Weyl Curvature:} \nonumber \\
    &\quad R_{\mu\nu} \equiv \text{Volumetric Strain} \label{eq:rosetta_ricci} \\ 
    &\quad C_{\mu\nu\rho\sigma} \equiv \text{Deviatoric Shear} \label{eq:rosetta_weyl} \\
    &\text{6. Stress-Energy \& Einstein Tensor (Equilibrium):} \nonumber \\
    &\quad T_{\mu\nu} \equiv \sigma_{\mu\nu} + \rho u_\mu u_\nu \label{eq:rosetta_stress_energy} \\
    &\quad G_{\mu\nu} \equiv \text{Hookean Strain Balance of } T_{\mu\nu} \label{eq:rosetta_einstein_tensor} \\
    &\text{7. Faraday Tensor (Kinematic Vorticity):} \nonumber \\
    &\quad F_{\mu\nu} \equiv \Omega_{\mu\nu} = \partial_\mu u_\nu - \partial_\nu u_\mu \label{eq:rosetta_faraday} \\
    &\text{8. Yang-Mills (Convective Acceleration):} \nonumber \\
    &\quad \partial_\mu F^{\mu\nu} + [A_\mu, F^{\mu\nu}] \equiv (\mathbf{u} \cdot \nabla)\mathbf{u} \label{eq:rosetta_yang_mills} \\
    &\text{9. Spinors \& $SU(2)$ (Cartan Moving Frames):} \nonumber \\
    &\quad \frac{d\mathbf{\hat{e}}_i}{dt} = \boldsymbol{\omega}_{fluid} \times \mathbf{\hat{e}}_i \label{eq:rosetta_cartan}
\end{align}
\end{tcolorbox}

\begin{tcolorbox}[cheatbox, title={\refstepcounter{cheatnum}\label{box:pressure_gradient}\thecheatnum. Continuous Pressure Dynamics  \scriptsize\hfill\hyperref[sec:deriv_pressure_gradient]{[Proof $\downarrow$]}}]
The \textbf{four fundamental forces} evaluate as geometric projections of a single continuous fluid pressure gradient.
\begin{align}
    &\text{1. Dynamic Radial Pressure:} \nonumber \\
    &\quad P(r) = \frac{1}{3} \frac{E}{V} = \frac{1}{3} \frac{\left(\hbar c / 2\pi r\right)}{\left(4\pi r^3 / 3\right)} = \frac{\hbar c}{8\pi^2 r^4} \label{eq:pressure_thermo} \\
    &\text{2. The Unified Force Gradient:} \nonumber \\
    &\quad F(r) = P(r) \cdot A(r) = \frac{\hbar c}{2\pi r^2} \label{eq:pressure_gradient} \\
    &\text{3. The Kinematic Mass Limit ($2\pi$ Rotation):} \nonumber \\
    &\quad E(r) = F(r) \cdot r = \frac{\hbar c}{2\pi r} \implies E = \frac{mc^2}{2\pi} \label{eq:energy_2pi} \\
    &\text{4. Gravity (The Absolute Core):} \nonumber \\
    &F_{grav} = F(l_p) = \frac{c^4}{2\pi G} \approx 1.926 \times 10^{43} N\label{eq:force_gravity} \\
    &\text{5. Weak Force (The Cavitation Limit):} \nonumber \\
    &F_{weak} = F(R_W) = \frac{m_W^2 c^3}{2\pi \hbar} \approx 8.348 \times 10^8 N\label{eq:force_weak} \\
    &\text{6. Strong Force (The Acoustic Boundary):} \nonumber \\
    &F_{strong} = F(L_{crit}) = \frac{m_e^2 c^3}{2\pi \alpha^2 \hbar} \approx 633.7 N \label{eq:force_strong} \\
    &\text{7. Electromagnetism (Macroscopic Baseline):} \nonumber \\
    &F_{EM_{base}} = F(a_0) = \frac{\alpha^2 m_e^2 c^3}{2\pi \hbar} \implies \frac{F_{EM_{base}}}{F_{strong}} = \alpha^4 \label{eq:force_em_ratio}
\end{align}
\end{tcolorbox}

\begin{tcolorbox}[cheatbox, title={\refstepcounter{cheatnum}\label{box:hookes_su3}\thecheatnum. Generalized Hooke's Law \& $SU(3)$ \scriptsize\hfill\hyperref[sec:deriv_hookes_su3]{[Proof $\downarrow$]}}]
The 8 geometric degrees of freedom of the $SU(3)$ strong force evaluate identically as the purely deviatoric strain tensor of an incompressible continuum.
\begin{align}
    &\text{1. The Isochoric Choke Limit (Zero-Trace):} \nonumber \\
    &\quad \text{Tr}(\varepsilon) = \varepsilon_{11} + \varepsilon_{22} + \varepsilon_{33} = 0 \label{eq:isochoric_trace} \\
    &\text{2. The Deviatoric Stress State (Hookean $SU(3)$):} \nonumber \\
    &\quad \sigma_{ij} = 2\mu_s \varepsilon'_{ij} \implies \text{Deviatoric } SU(3) \text{ Basis} \label{eq:deviatoric_stress} \\
    &\text{3. Non-Abelian Self-Interaction (Finite Strain):} \nonumber \\
    &\quad \varepsilon_{ij} = \frac{1}{2} \left( \partial_j u_i + \partial_i u_j - \partial_i u_k \partial_j u_k \right) \label{eq:su3_finite_strain} \\
    &\text{4. The Lie Bracket (Kinematic Non-Commutativity):} \nonumber \\
    &\quad [\lambda_a, \lambda_b] = 2i f_{abc} \lambda_c \label{eq:non_abelian_commutator}
\end{align}
\end{tcolorbox}

\end{multicols}

\newpage
% ======================================================================
% PART II: FORMAL DERIVATIONS
% ======================================================================
\begin{center}
    \LARGE\textbf{PART II: FORMAL DERIVATIONS}
    \label{sec:part2}
\end{center}
\vspace{0.4cm}


% ==========================================
% PHASE I: THE MACROSCOPIC CONTINUUM
% ==========================================
\part*{PHASE I: THE MACROSCOPIC CONTINUUM}


% ----------------------------------------------------------------------
\section{Macroscopic Kinematics and the Causal Boundary}
\label{sec:deriv_axiom}
The macroscopic causal boundary from within the continuum expands linearly at the speed of light ($c$) over evolutionary proper time ($t$). Integrating the kinematic limit of the fluid evaluates the absolute scale radius:
\begin{align}
    R_h(t) &= \int_{0}^{t} c \, dt \\
    R_h(t) &= c \cdot t
\end{align}

The cosmological scale factor $a(t)$ evaluates as the geometric ratio of the horizon radius at emission time $t$ to the horizon radius at the present observation time $\tau_0$:
\begin{equation}
    a(t) = \frac{R_h(t)}{R_{h,0}}
\end{equation}

Kinematic stretching of the $\boldsymbol{\tau}$-fluid geometrically dilates the observed wavelength, defining the cosmological redshift $z$:
\begin{align}
    1 + z &= \frac{1}{a(t)} \\
    1 + z &= \frac{R_{h,0}}{R_h(t)}
\end{align}

Substituting the linear expansion axiom ($R_h(t) = c \cdot t$)  isolates the mathematical translation between Relative Observational Redshift ($z$) and Evolutionary Proper Time ($t$):
\begin{align}
    1 + z &= \frac{c \cdot \tau_0}{c \cdot t} \\
    1+z &= \frac{\tau_0}{t}
\end{align}

To establish the absolute proper time of the present epoch ($z=0$), we evaluate the current scale radius $R_{h,0} = 4224 \text{ Mpc}$ against the phase velocity limit $c$:
\begin{equation}
    \tau_0 = \frac{R_{h,0}}{c} \approx 13.78 \text{ Gyr}
\end{equation}

\begin{tcolorbox}[recallbox, title={\ref{box:axiom}. Fundamental Axiom \& Translation}]
\begin{equation}
    R_h(t) = c \cdot t \iff 1+z = \frac{\tau_0}{t} \tag{\ref{eq:axiom}}
\end{equation}
\end{tcolorbox}

% ----------------------------------------------------------------------
\section{The High-Redshift Chronometer}
\label{sec:deriv_chronometer}

Because the geometric stretch of an observed signal represents the exact algebraic ratio of the causal horizon today ($R_{h,0}$) to the causal horizon at the epoch of emission ($R_h(z)$), we establish:
\begin{equation}
    1+z = \frac{R_{h,0}}{R_h(z)}
\end{equation}

Because the continuum scales linearly ($R_h \propto t$), we substitute proper time directly into this ratio. If the baseline proper age of the metric evaluates to $\tau_0$, the exact evolutionary proper age of the metric at redshift $z$ algebraically isolates to:
\begin{align}
    1+z &= \frac{c \cdot \tau_0}{c \cdot t(z)} \\
    1+z &= \frac{\tau_0}{t(z)} \\
    t(z) &= \frac{\tau_0}{1+z}
\end{align}

Substituting this translated proper time $t(z)$ back into the fundamental expansion axiom ($R_h = c \cdot t$) natively derives the local scale radius at the epoch of emission:
\begin{align}
    R_h(z) &= c \cdot t(z) \\
    R_h(z) &= c \left( \frac{\tau_0}{1+z} \right) \\
    R_h(z) &= \frac{c \cdot \tau_0}{1+z} \\
    R_h(z) &= \frac{R_{h,0}}{1+z}
\end{align}

\vspace{0.2cm}
\textbf{Empirical Evaluation (The Early Galactic Boundary):} \\
Evaluating the continuous temporal scaling at extreme optical horizons ($z = 14$ and $z = 20$) geometrically isolates the absolute proper age of the metric at the exact epoch of emission. 

Substituting the baseline proper age of the macroscopic continuum ($\tau_0 \approx 13.8 \text{ Gyr}$), the proper evolutionary age of the localized metric explicitly evaluates to:
\begin{align}
    t(14) &= \frac{13.8 \text{ Gyr}}{1 + 14} = \frac{13.8 \text{ Gyr}}{15} \approx 920 \text{ Myr} \\
    t(20) &= \frac{13.8 \text{ Gyr}}{1 + 20} = \frac{13.8 \text{ Gyr}}{21} \approx 657 \text{ Myr}
\end{align}

The total elapsed proper time evaluated by an observer during the optical transit ($\Delta t$) natively isolates as the absolute temporal difference between the current boundary and the emission epoch:
\begin{align}
    \Delta t(14) &= \tau_0 - t(14) \approx 13.8 \text{ Gyr} - 0.920 \text{ Gyr} = 12.88 \text{ Gyr} \\
    \Delta t(20) &= \tau_0 - t(20) \approx 13.8 \text{ Gyr} - 0.657 \text{ Gyr} = 13.14 \text{ Gyr}
\end{align}

Because the linear chronometer fundamentally expands the early temporal envelope, evaluating a $z=14$ boundary natively affords nearly one billion proper years of temporal evolution. This mathematically guarantees the requisite chronological boundary for macroscopic structure formation, natively resolving the observation of highly evolved, massive galaxies at extreme optical redshifts.

\begin{tcolorbox}[recallbox, title={\ref{box:chronometer}. The High-Redshift Chronometer}]
\begin{align}
    t(z) &= \frac{\tau_0}{1+z} \tag{\ref{eq:tz}} \\
    R_h(z) &= \frac{R_{h,0}}{1+z} \equiv c \cdot t(z) \tag{\ref{eq:rz}}
\end{align}
\end{tcolorbox}


% ----------------------------------------------------------------------
\section{The Geometric Spacetime Bridge}
\label{sec:deriv_gr_bridge}
To derive the continuous Velocity Field ($\mathbf{v}_0$) of the expanding metric, we evaluate the accumulation of redshift over spatial distance. A wave traveling an infinitesimal distance $dr$ experiences a fractional geometric stretch $dz$. Integrating this stretch across the transit distance $r$ yields a natural logarithm:
\begin{equation}
    \int_0^z \frac{1}{1+z'} \, dz' = \int_0^r \frac{1}{R_{h,0}} \, dr'
\end{equation}

Evaluating the definite integrals yields the geometric stretch factor:
\begin{align}
    \Big[ \ln(1+z') \Big]_0^z &= \left[ \frac{r'}{R_{h,0}} \right]_0^r \\
    \ln(1+z) - \ln(1) &= \frac{r}{R_{h,0}} - 0 \\
    \ln(1+z) &= \frac{r}{R_{h,0}} \\
    1+z &= e^{r/R_{h,0}}
\end{align}

The total kinematic strain ($1+z$) of a wave escaping a continuous metric expanding outward at velocity $v_0$ evaluates as the product of the acoustic Doppler shift and the fluid time dilation:
\begin{equation}
    1+z = \sqrt{\frac{1+v_0/c}{1-v_0/c}} \times \frac{1}{\sqrt{1-v_0^2/c^2}}
\end{equation}

We algebraically expand the denominator of the time dilation term as a difference of squares ($1-v_0^2/c^2 = (1-v_0/c)(1+v_0/c)$):
\begin{equation}
    1+z = \frac{\sqrt{1+v_0/c}}{\sqrt{1-v_0/c}} \times \frac{1}{\sqrt{1-v_0/c}\sqrt{1+v_0/c}}
\end{equation}

The $\sqrt{1+v_0/c}$ terms in the numerator and denominator cancel:
\begin{align}
    1+z &= \frac{1}{\sqrt{1-v_0/c}\sqrt{1-v_0/c}} \\
    1+z &= \frac{1}{1 - v_0/c}
\end{align}

Equating this kinematic fluid strain directly to the evaluated geometric stretch ($e^{r/R_{h,0}}$) isolates the velocity potential:
\begin{align}
    \frac{1}{1-v_0/c} &= e^{r/R_{h,0}} \\
    1-\frac{v_0}{c} &= e^{-r/R_{h,0}} \\
    \frac{v_0}{c} &= 1 - e^{-r/R_{h,0}} \\
    v_0(r) &= c\left(1 - e^{-r/R_{h,0}}\right)
\end{align}

Applying the Equivalence Principle, existence within a macroscopic velocity gradient evaluates as mechanically identical to existing within a stationary gravitational potential ($\Phi$). The potential energy corresponds to the kinetic energy of this fluid flow: $\Phi(r) = -v_0(r)^2/2$. The time-time component of the structural curvature tensor ($g_{tt}$) natively depends on this fluid potential:
\begin{align}
    g_{tt} &= -\left( 1 + \frac{2\Phi}{c^2} \right) \\
    g_{tt} &= -\left( 1 + \frac{2(-v_0^2/2)}{c^2} \right) \\
    g_{tt} &= -\left( 1 - \frac{v_0^2}{c^2} \right)
\end{align}

\textbf{The Pressure-Curvature Identity:} \\
To map this geometric curvature to the continuous thermodynamic state of the metric, we isolate the fluid velocity from the classical Bernoulli dynamic pressure equation ($P_{dyn} = \frac{1}{2}\rho_{\tau} v_0^2$):
\begin{equation}
    v_0^2 = \frac{2 P_{dyn}}{\rho_{\tau}}
\end{equation}

Substituting this kinematic velocity directly into the derived time-time curvature tensor ($g_{tt}$) yields the absolute continuum identity:
\begin{align}
    g_{tt} &= -\left( 1 - \frac{\left( \frac{2 P_{dyn}}{\rho_{\tau}} \right)}{c^2} \right) \nonumber \\
    g_{tt} &= -\left( 1 - \frac{2 P_{dyn}}{\rho_{\tau} c^2} \right) \nonumber
\end{align}

Because the denominator ($\rho_{\tau} c^2$) evaluates as twice the absolute acoustic yield pressure ($2P_c$), the fraction reduces identically to the residual static pressure ratio ($1 - P_{dyn}/P_c$). By definition of the Admittance Tensor, this isolates the exact limit of fluid compliance:
\begin{equation}
    g_{tt} = -\alpha_s^2 \tag{\ref{eq:pressure_curvature_identity}}
\end{equation}

Time dilation ($g_{tt}$) evaluates mechanically as the negative squared Total Scalar Admittance ($-\alpha_s^2$). Spacetime curvature is not a geometric deformation of nothingness; it evaluates as the macroscopic projection of a continuous scalar pressure deficit. 

Substituting the fluid admittances directly into the generalized spacetime interval ($ds^2 = g_{tt} (c dt)^2 + g_{rr} dr^2$) perfectly constructs the geometric curvature tensor for the expanding continuum, governed exclusively by fluid compliance:

\begin{tcolorbox}[recallbox, title={\ref{box:gr_bridge}. The Geometric Spacetime Bridge}]
\begin{align}
    v_0(r) &= c\left(1-e^{-r/R_{h,0}}\right) \tag{\ref{eq:v0}} \\
    g_{tt} &= -\left( 1 - \frac{P_{dyn}}{P_c} \right) \equiv -\alpha_s^2 \tag{\ref{eq:pressure_curvature_identity}} \\
    ds^2 &= -\alpha_s^2(c dt)^2 + \frac{dr^2}{\alpha_s^2} \tag{\ref{eq:metric}}
\end{align}
\end{tcolorbox}

% ----------------------------------------------------------------------
\section{Kinematic Gradients}
\label{sec:deriv_kinematic_gradients}
Apparent expansion $H_{app}$ evaluates as the spatial gradient of the macroscopic velocity potential $v_0(r)$. Taking the derivative of $v_0(r) = c(1 - e^{-r/R_{h,0}})$ using the explicit chain rule ($\frac{d}{dr}[e^u] = e^u \frac{du}{dr}$) yields:
\begin{align}
    H_{app}(r) &= \frac{dv_0}{dr} \\
    H_{app}(r) &= \frac{d}{dr} \Big[ c - c e^{-r/R_{h,0}} \Big] \\
    H_{app}(r) &= 0 - c \left( -\frac{1}{R_{h,0}} \right) e^{-r/R_{h,0}} \\
    H_{app}(r) &= \frac{c}{R_{h,0}} e^{-r/R_{h,0}}
\end{align}

Evaluating this spatial gradient at the absolute localized coordinate ($r = 0$) algebraically evaluates the exponential to exactly $1$:
\begin{align}
    H_{local} &= \frac{c}{R_{h,0}} e^0 \\
    H_{local} &= \frac{c}{R_{h,0}}
\end{align}

Substituting the fundamental expansion axiom ($R_{h,0} = c \cdot t$), the acoustic velocity limit $c$ perfectly cancels out, rigorously proving the spatial velocity gradient natively equates to the proper temporal expansion rate ($H_{prop}$):
\begin{align}
    H_{prop} &= \frac{c}{c \cdot t} \\
    H_{prop} &= \frac{1}{t}
\end{align}

Because an observer evaluates physically as a topological defect of finite geometric radius (Planck length, $l_p$), the absolute spatial coordinate evaluates as structurally empty below $r = l_p$. Evaluating the spatial gradient at this absolute physical boundary:
\begin{equation}
    H_{app}(l_p) = \frac{c}{R_{h,0}} e^{-l_p/R_{h,0}}
\end{equation}

Because $l_p \ll R_{h,0}$, applying the first-order Taylor expansion ($e^{-x} \approx 1 - x$) isolates the discrete geometric strain:
\begin{align}
    H_{app}(l_p) &\approx \frac{c}{R_{h,0}} \left( 1 - \frac{l_p}{R_{h,0}} \right) \\
    H_{app}(l_p) &\approx \frac{c}{R_{h,0}} - \frac{c \cdot l_p}{R_{h,0}^2}
\end{align}

Substituting the baseline proper rate ($c/R_{h,0} = 1/t$), the expansion explicitly isolates a negative kinematic offset structurally truncating the mathematical singularity:
\begin{equation}
    H_{app}(l_p) \approx \frac{1}{t} - \frac{c \cdot l_p}{R_{h,0}^2}
\end{equation}

At standard macroscopic observation boundaries corresponding to redshift $z$, the geometric stretch ratio substitutes as $e^{-r/R_{h,0}} = 1/(1+z)$:
\begin{equation}
    H_{app}(z) = \frac{c}{R_{h,0}(1+z)}
\end{equation}

The proper boundary acceleration $a_{cosmic}$ evaluates as the centripetal acceleration of the causal horizon expanding at the phase limit $c$:
\begin{align}
    a_{cosmic} &= \frac{c^2}{R_h} \\
    a_{cosmic} &= \frac{c^2}{c \cdot t} \\
    a_{cosmic} &= \frac{c}{t}
\end{align}

Because the geometric boundary expands   at the invariant acoustic phase limit ($c$), the apparent boundary acceleration evaluates natively as the product of this phase velocity and the apparent spatial gradient ($H_{app}$):
\begin{equation}
    a_{app}(z) = c \cdot H_{app}(z)
\end{equation}

Substituting the previously isolated apparent spatial gradient ($H_{app}(z) = c / [R_{h,0}(1+z)]$) yields the exact conformal boundary acceleration:
\begin{align}
    a_{app}(z) &= c \left( \frac{c}{R_{h,0}(1+z)} \right) \\
    a_{app}(z) &= \frac{c^2}{R_{h,0}(1+z)}
\end{align}

\begin{tcolorbox}[recallbox, title={\ref{box:kinematic_gradients}. Kinematic Gradients}]
\begin{align}
    H_{app}(z) &= \frac{c}{R_{h,0}(1+z)} \iff H_{prop} = \frac{1}{t} \tag{\ref{eq:happ}} \\
    a_{app}(z) &= \frac{c^2}{R_{h,0}(1+z)} \iff a_{cosmic} = \frac{c}{t} \tag{\ref{eq:aapp}}
\end{align}
\end{tcolorbox}


% ----------------------------------------------------------------------
\section{Optical Dilatation}
\label{sec:deriv_optical_dilatation}

The deceleration parameter ($q$) evaluates kinematically as the dimensionless ratio of structural acceleration to the expansion velocity:
\begin{equation}
    q = -\frac{\ddot{R}_h R_h}{\dot{R}_h^2}
\end{equation}

Evaluating the time derivatives of the fundamental continuum expansion axiom ($R_h(t) = c \cdot t$) yields:
\begin{align}
    \dot{R}_h(t) &= \frac{d}{dt}(c \cdot t) = c \\
    \ddot{R}_h(t) &= \frac{d}{dt}(c) = 0
\end{align}

Substituting these explicit temporal derivatives mathematically zeroes the deceleration parameter:
\begin{equation}
    q = -\frac{(0)(c \cdot t)}{c^2} = 0
\end{equation}

Because the continuous metric evaluates with zero localized deceleration, the True Coordinate Distance ($r_{true, cosmic}$) isolates   by integrating the redshift geometric stretch parameter over the macroscopic scale radius:
\begin{align}
    r_{true, cosmic}(z) &= \int_0^z \frac{R_{h,0}}{1+z'} dz' \\
    r_{true, cosmic}(z) &= R_{h,0} \Big[\ln(1+z')\Big]_0^z \\
    r_{true, cosmic}(z) &= R_{h,0} (\ln(1+z) - \ln(1)) \\
    r_{true, cosmic}(z) &= R_{h,0} \ln(1+z)
\end{align}

The apparent Luminosity Distance ($D_L$) perceived across the continuum evaluates simply as the true integrated spatial distance multiplied by the geometric optical stretch factor $(1+z)$:
\begin{align}
    D_L(z) &= r_{true, cosmic}(z)(1+z) \\
    D_L(z) &= R_{h,0} \ln(1+z) \cdot (1+z) \\
    D_L(z) &= R_{h,0} (1+z) \ln(1+z)
\end{align}

\begin{tcolorbox}[recallbox, title={\ref{box:optical_dilatation}. Optical Dilatation}]
\textit{Macroscopic Conformal Mirror:}
\begin{equation}
    r_{true, cosmic}(z) = R_{h,0} \ln(1+z) \tag{\ref{eq:r_cosmic}}
\end{equation}
\textit{Luminosity Distance:}
\begin{equation}
    D_L(z) = R_{h,0} (1+z) \ln(1+z) \tag{\ref{eq:dl}}
\end{equation}
\end{tcolorbox}






% ----------------------------------------------------------------------
\section{Cosmic Expansion \& The Invariant Acoustic Limit}
\label{sec:deriv_acoustic_invariance}

The outermost geometric boundary of the continuum evaluates as a cascading acoustic nucleation expanding at the maximum kinematic limit ($R_h = c \cdot t$). The internal thermodynamic stability of this expansion requires an invariant phase velocity, maintained by the coupled decay of fluid density and elastic tension.

\textbf{1. The $w = -1/3$ Density Decay:} \\
Unlike inert baryonic dust which dilutes purely volumetrically ($V \propto t^3 \implies \rho \propto t^{-3}$), the continuous metric evaluates as a structurally tensioned medium via the $w = -1/3$ elastic equation of state. 

Evaluating the continuous fluid energy conservation equation isolates the dynamic density decay:
\begin{equation}
    \dot{\rho} + 3H\left(\rho + \frac{P}{c^2}\right) = 0 
\end{equation}
Substituting the tensioned equation of state ($P = -\frac{1}{3}\rho c^2$):
\begin{align}
    \dot{\rho} + 3H\left(\rho - \frac{1}{3}\rho\right) &= 0  \\
    \dot{\rho} + 2H\rho &= 0 
\end{align}

Because the macroscopic proper expansion rate evaluates as $H = 1/t$, the differential equation isolates:
\begin{align}
    \frac{d\rho}{dt} + \frac{2}{t}\rho &= 0  \\
    \frac{d\rho}{\rho} &= -2 \frac{dt}{t} 
\end{align}
Integrating both sides yields the geometric decay of the macroscopic vacuum density as an absolute inverse square of proper time:
\begin{equation}
    \rho_{\boldsymbol{\tau}}(t) \propto t^{-2} \tag{\ref{eq:density_decay}}
\end{equation}

\textbf{2. The Decay of Transverse Shear Tension:} \\
The resistance of the continuum to lateral deformation evaluates as the Macroscopic Shear Modulus ($\mu_s$). Because the ambient macroscopic vacuum pressure ($P_{macro}$) decays geometrically as a function of the expanding metric volume ($P \propto t^{-2}$), the structural stiffness of the fluid proportionately degrades. The transverse shear modulus mathematically tracks the identical temporal decay vector as the ambient pressure gradient:
\begin{equation}
    \mu_s(t) \propto t^{-2} \tag{\ref{eq:stiffness_decay}}
\end{equation}

\textbf{3. The Invariant Transverse Speed Limit ($c$):} \\
By the fundamental elastodynamic equation of state for transverse wave propagation, the kinematic phase velocity ($c$) evaluates as the square root of the macroscopic shear modulus ($\mu_s$) divided by the inertial fluid density ($\rho_{\boldsymbol{\tau}}$). 

Substituting the coupled temporal decay limits isolates the absolute temporal evolution of the speed of light:
\begin{align}
    c(t) &= \sqrt{\frac{\mu_s(t)}{\rho_{\boldsymbol{\tau}}(t)}} \tag{\ref{eq:c_invariant}} \\
    c(t) &\propto \sqrt{\frac{t^{-2}}{t^{-2}}}  \\
    c(t) &\propto \sqrt{1} 
\end{align}

Because macroscopic transverse fluid tension and inertial fluid density dynamically degrade at the exact same geometric temporal rate, their thermodynamic fraction permanently isolates to unity. The transverse kinematic speed of light evaluates mechanically as an absolute, invariant constant. The universe perfectly scales its own elasticity to guarantee that the internal acoustic phase limit never fractures.

\begin{tcolorbox}[recallbox, title={\ref{box:acoustic_invariance}. Cosmic Expansion \& Acoustic Invariance}]
\begin{align}
    \rho_{\boldsymbol{\tau}}(t) &\propto t^{-2} \tag{\ref{eq:density_decay}} \\
    c &= \sqrt{\frac{\mu_s}{\rho_{\boldsymbol{\tau}}}} = \text{Constant} \tag{\ref{eq:c_invariant}}
\end{align}
\end{tcolorbox}


% ----------------------------------------------------------------------
\section{Baryon Acoustic Oscillations (Cosmological Lensing)}
\label{sec:deriv_bao_lensing}

Baryon Acoustic Oscillations (BAO) evaluate as primordial acoustic waves frozen within the macroscopic fluid distribution. The spatial divergence between the apparent Transverse Width ($D_M$) and the apparent Longitudinal Depth ($D_H$) evaluates as lensed cosmic information propagating through the macroscopic kinematic vacuum gradient.

\textbf{1. Longitudinal Depth ($D_H$) as a Spatial Derivative:} \\
In continuous kinematic geometry, the true physical depth spanning a given redshift bin evaluates as the spatial derivative of the logarithmic coordinate distance ($r_{true,cosmic} = R_{h,0} \ln(1+z)$) with respect to redshift ($z$):
\begin{align}
    D_{H,true}(z) &= \frac{d}{dz} r_{true,cosmic}(z) \\
    D_{H,true}(z) &= \frac{d}{dz} \Big(R_{h,0} \ln(1+z)\Big)
\end{align}

Applying the chain rule for the natural logarithm:
\begin{align}
    D_{H,true}(z) &= R_{h,0} \left( \frac{1}{1+z} \right) \frac{d}{dz}(1+z) \\
    D_{H,true}(z) &= \frac{R_{h,0}}{1+z} (1) \\
    D_{H,true}(z) &= \frac{R_{h,0}}{1+z} \tag{\ref{eq:bao_longitudinal}}
\end{align}

The line-of-sight compression of the BAO ripples at high redshifts evaluates   as the spatial derivative of the macroscopic longitudinal strain, completely independent of optical refraction. 

\textbf{2. Transverse Width ($D_M$) and Optical Magnification:} \\
While the longitudinal depth evaluates via the macroscopic spatial derivative, the transverse plane is subjected to continuous optical magnification. Because the transverse impedance of the metric ($Z_{\tau(\perp)} \propto n$) diverges as the background scalar potential reaches the acoustic limit ($v_0 \to c$), the apparent optical path length ($r_{opt}$) evaluates as longer than the true geometric distance. The measured transverse width ($D_M$) evaluates as the integrated apparent optical distance ($r_{opt} = \int n(r)dr$).

\textbf{3. Unification of BAO and Galactic Kinematics:} \\
The absolute physical depth ($D_{H,true}$) represents the fundamental length scale of the macroscopic vacuum gradient at any given redshift. The continuous Unruh-Hawking geometric boundary acceleration evaluates as $a_0(z) = \frac{c H(z)}{2\pi}$. 

Substituting the apparent Hubble parameter ($H(z) = \frac{c}{R_{h,0}/(1+z)}$) into the acceleration:
\begin{align}
    a_0(z) &= \frac{c}{2\pi} \left( \frac{c}{R_{h,0} / (1+z)} \right) \\
    a_0(z) &= \frac{c^2}{2\pi \left( \frac{R_{h,0}}{1+z} \right)}
\end{align}

Substituting the derived spatial depth ($D_{H,true}(z)$) isolates the physical identity:
\begin{equation}
    a_0(z) = \frac{c^2}{2\pi D_{H,true}(z)} \tag{\ref{eq:bao_kinematic_unification}}
\end{equation}

The spatial gradient of the macroscopic tension compresses the acoustic waves along the line of sight ($D_H$). As a bound galactic disk resists this geometric strain, it generates the thermodynamic temperature floor ($a_0$) bounding its rotation.

\begin{tcolorbox}[recallbox, title={\ref{box:bao_lensing}. Baryon Acoustic Oscillations \& Kinematic Unification}]
\begin{align}
    D_{H,true}(z) &= \frac{R_{h,0}}{1+z} \tag{\ref{eq:bao_longitudinal}} \\
    a_0(z) &= \frac{c^2}{2\pi D_{H,true}(z)} \tag{\ref{eq:bao_kinematic_unification}}
\end{align}
\end{tcolorbox}


% ----------------------------------------------------------------------
\section{Pressure Deficits, Symmetry Breaking, and Acoustic Covariance}
\label{sec:deriv_auxetic_symmetry_covariance}

The geometric stability of adjacent topological boundaries evaluates via the spontaneous symmetry breaking of continuous fluid pressure fields. This mechanism decouples macroscopic spatial acceleration from gravitational time dilation, mathematically resolving the kinematics of relative observation into exact acoustic covariance.

\textbf{1. The Stagnation Plane (Gradient Annihilation):} \\
When two topological circulations (Region $A$ and Region $B$) are positioned adjacent to each other, they evaluate as intersecting localized pressure deficits within the global continuous metric. 

Because fluid acceleration is driven exclusively by the spatial pressure gradient ($\mathbf{a} = -\frac{1}{\rho_{\tau}} \nabla P$), the exact geometric midpoint between the two boundaries establishes a stagnation plane. In this shared interstitial space, the pressure gradient of basin $A$ exactly opposes the pressure gradient of basin $B$. The physical macroscopic gradients mathematically annihilate:
\begin{equation}
    \nabla P_{net}(0) = 0 \implies \mathbf{a} = 0 \tag{\ref{eq:gradient_cancellation}}
\end{equation}
Because the macroscopic spatial pressure gradient evaluates to zero, the geometric space evaluates as absolute kinematic non-acceleration (Zero Gravity).

\textbf{2. Hidden Action (Scalar Superposition):} \\
While the localized spatial gradient evaluates to zero, the topological fluid-dynamic action remains continuous. Bernoulli dynamic pressure ($P_{dyn}$) evaluates as a purely scalar magnitude, irrespective of vector orientation. 

Therefore, while the directional gradients cancel, the absolute structural pressure deficits algebraically superimpose:
\begin{equation}
    \Delta P_{net} = P_{dyn,A} + P_{dyn,B} \tag{\ref{eq:scalar_superposition}}
\end{equation}
The continuous $\boldsymbol{\tau}$-fluid in the interstitial space evaluates as kinematically stationary ($\mathbf{a} = 0$), yet it is subjected to massive, bidirectional thermodynamic strain ($\sigma_{ij} \gg 0$).
\begin{equation}
    \mathbf{a} = 0, \quad \sigma_{ij} \gg 0 \tag{\ref{eq:hidden_action_stress}}
\end{equation}

\textbf{3. Gravitational Time Dilation:} \\
This overlapping scalar deficit physically depletes the available static pressure of the localized metric. Because structural fluid compliance ($\alpha_s$) dictates the internal harmonic phase of a particle, the time-time structural metric ($g_{tt}$) is governed exclusively by this accumulated scalar deficit.

Substituting the superimposed dynamic pressure into the Bernoulli compliance limit mechanically drives Gravitational Time Dilation independent of spatial translation:
\begin{equation}
    d\tau = dt_{bg} \sqrt{1 - \frac{\Delta P_{net}}{P_c}} \tag{\ref{eq:scalar_time_dilation}}
\end{equation}

\textbf{4. Acoustic Covariance:} \\
The geometric experience of proper time evaluates via two distinct fluid-dynamic states:
\begin{itemize}[leftmargin=*]
    \item \textbf{The Topological Boundary (The Clock):} A confined circulation experiences time as the localized kinematic evolution of the surrounding metric. Its internal resonant clock is retarded because the ambient scalar pressure deficit ($\Delta P_{net}$) degrades the structural compliance of its specific spatial coordinate.
    \item \textbf{The Transverse Wave (The Signal):} An unconfined acoustic shear wave possesses no internal rest clock. It experiences time as the cumulative spatial evolution of the metric. Its transit evaluates as the spatial line integral of the macroscopic refractive index across the tensioned fluid.
\end{itemize}
Because both the topological clock and the optical messenger are physical manifestations of the identical continuous fluid, they are structurally bound by the exact same localized fluid impedance. The depleted interstitial pressure mechanically retards the particle's internal clock by the exact same geometric fraction that it refracts and delays the transiting wave. 

When the retarded local clock measures the delayed transverse wave, the fluid impedances algebraically cancel:
\begin{equation}
    c_{obs} = \frac{dr_{wave}}{d\tau} \equiv c \tag{\ref{eq:acoustic_covariance}}
\end{equation}
The transverse acoustic wave evaluates as the absolute kinematic baseline that calibrates all observers to an invariant covariant reality.

\begin{tcolorbox}[recallbox, title={\ref{box:auxetic_symmetry_covariance}. Symmetry Breaking \& Covariance}]
\begin{align}
    \nabla P_{net}(0) = 0 &\implies \mathbf{a} = 0 \tag{\ref{eq:gradient_cancellation}} \\
    \mathbf{a} = 0, &\quad \sigma_{ij} \gg 0 \tag{\ref{eq:hidden_action_stress}} \\
    \Delta P_{net} &= P_{dyn,A} + P_{dyn,B} \tag{\ref{eq:scalar_superposition}} \\
    d\tau &= dt_{bg} \sqrt{1 - \frac{\Delta P_{net}}{P_c}} \tag{\ref{eq:scalar_time_dilation}} \\
    c_{obs} &= \frac{dr_{wave}}{d\tau} \equiv c \tag{\ref{eq:acoustic_covariance}}
\end{align}
\end{tcolorbox}


% ----------------------------------------------------------------------
\section{Entropic Superposition (MOND)}
\label{sec:deriv_mond}

The effective gravitational field $g_{eff}$ governing galactic rotation evaluates as the geometric superposition of the localized baryonic mass and the macroscopic background acceleration of the universe $a_0$. 

\textbf{1. Redshift Evolution of Boundary Acceleration ($a_0(z)$):} \\
As derived via the spatial derivative of the macroscopic longitudinal strain, the absolute continuous boundary acceleration of the universe $a_0$ is fundamentally locked to the macroscopic scale radius $R_{h,0}$. Because the apparent kinematic acceleration of the metric weakens with optical depth, the thermodynamic background acceleration applied to a localized galaxy evaluates as a dynamic, epoch-dependent boundary condition:
\begin{equation}
    a_0(z) = \frac{c^2}{2\pi R_{h,0}(1+z)} \approx 1.10 \times 10^{-10} \text{ m/s}^2 \quad \text{(at } z=0\text{)} \tag{\ref{eq:a0z}} 
\end{equation}

\textbf{2. Kinematic Confinement and Superposition:} \\
In an empty void, the Newtonian field dilutes freely ($1/r^2$). However, within the continuous metric, the gravitational field is physically confined by the ambient cosmic pressure generated by the Unruh floor. Consequently, the active energy density of the field evaluates as the geometric superposition of the localized baryonic pull ($g_{bar}$) and the macroscopic background baseline ($a_0$).

For a star in a circular orbit, the centripetal acceleration evaluates as $g_{eff} = v_{rot}^2/r$. Squaring both sides evaluates the physical energy density of the rotation:
\begin{align}
    g_{eff}^2 &= \left( \frac{v_{rot}^2}{r} \right)^2 \nonumber \\
    g_{eff}^2 &= \frac{v_{rot}^4}{r^2} \nonumber
\end{align}

Substituting this into the field superposition equation ($g_{eff}^2 = g_{bar} \cdot a_0(z)$) and expanding Newtonian gravity for the baryonic core ($g_{bar} = G M_{bar} / r^2$):
\begin{align}
    \frac{v_{rot}^4}{r^2} &= \left( \frac{G M_{bar}}{r^2} \right) a_0(z) \nonumber \\
    \frac{v_{rot}^4}{r^2} &= \frac{G M_{bar} a_0(z)}{r^2} \nonumber
\end{align}

\textbf{3. Annihilation of Radial Distance:} \\
The radial distance term $r^2$ appears in the denominator on both sides of the metric equality. Multiplying both sides by $r^2$ completely annihilates the spatial distance variable from the continuous rotation mechanics:
\begin{equation}
    v_{rot}^4 = G M_{bar} \cdot a_0(z) \nonumber
\end{equation}

Taking the fourth root derives the flat rotation curve directly from classical continuum mechanics, mathematically proving that orbital velocities at extreme radii evaluate entirely independently of distance from the core:
\begin{align}
    v_{rot} &= \left( G M_{bar} \cdot a_0(z) \right)^{1/4} \tag{\ref{eq:vrot}}
\end{align}

\textbf{4. The Fictitious Halo Equivalence:} \\
Define the predicted enclosed dynamical mass $M_{pred}$ required to sustain the rotation curve $v_{rot}$ at radius $r$:
\begin{equation}
    M_{pred} = \frac{v_{rot}^2 r}{G} \nonumber
\end{equation}

The effective gravitational self-force $F_{halo}$ exerted by this mass is:
\begin{equation}
    F_{halo} = \frac{G M_{pred}^2}{r^2} \nonumber
\end{equation}

Substitute $M_{pred}$ into the force equation:
\begin{align}
    F_{halo} &= \frac{G \left( \frac{v_{rot}^2 r}{G} \right)^2}{r^2} \nonumber \\
    F_{halo} &= \frac{G \left( \frac{v_{rot}^4 r^2}{G^2} \right)}{r^2} \nonumber \\
    F_{halo} &= \frac{G v_{rot}^4 r^2}{r^2 G^2} \nonumber
\end{align}

Cancel $r^2$ and $G$:
\begin{equation}
    F_{halo} = \frac{v_{rot}^4}{G} \nonumber
\end{equation}

Rearranging the kinematic confinement equation yields $\frac{v_{rot}^4}{G} = M_{bar} \cdot a_0(z)$. Substitute this equivalence:
\begin{equation}
    F_{halo} = M_{bar} \cdot a_0(z) \nonumber
\end{equation}

This mathematical equivalence proves that the anomalous force attributed to invisible enclosed mass ($F_{halo}$) evaluates identically to the physical force of the localized baryonic mass ($M_{bar}$) confined by the ambient cosmic boundary pressure ($a_0(z)$). This geometric fluid equation eradicates the requirement for Dark Matter halos, demonstrating that flat galactic rotation curves are the necessary mechanical consequence of gravitational confinement. Furthermore, it mathematically proves that these curves must evolve dynamically over macroscopic time due to the changing optical boundary condition $a_0(z)$.

\begin{tcolorbox}[recallbox, title={\ref{box:mond}. Entropic Superposition (MOND)}]
\begin{align}
    a_0(z) &= \frac{c^2}{2\pi R_{h,0}(1+z)} \approx 1.10 \times 10^{-10} \text{ m/s}^2 \quad \text{(at } z=0\text{)}\tag{\ref{eq:a0z}} \\
    g_{eff}^2 &= g_{bar} \cdot a_0(z) \tag{\ref{eq:geff}} \\
    v_{rot} &= \left( G M_{bar} \cdot a_0(z) \right)^{1/4} \tag{\ref{eq:vrot}}
\end{align}
\end{tcolorbox}

% ----------------------------------------------------------------------
\section{Refractive Gravity \& Dark Matter (Macroscopic GRIN)}
\label{sec:deriv_macroscopic_grin}

The high-velocity fluid circulation within a localized topological defect ($V_\tau$) induces a massive localized pressure drop via the classical Bernoulli principle. To enforce continuous mechanical equilibrium against the macroscopic baseline vacuum tension ($P_0 = -1/3\rho_{\tau} c^2$), the metric establishes a static, spherically symmetric hydrostatic pressure gradient ($\nabla P$) extending outward from the boundary.

\textbf{1. The Localized Acoustic Limit \& Refraction:} \\
This continuous pressure gradient alters the localized baseline fluid density ($\rho$) and restorative shear modulus ($\mu_s$). By the Newton-Laplace acoustic limit, this induces a continuous gradient in the local propagation velocity of transverse shear waves:
\begin{equation}
    c(r) = \sqrt{\frac{\mu_s(r)}{\rho(r)}} \tag{\ref{eq:acoustic_limit}}
\end{equation}
The localized refractive index ($n$) of the metric vacuum evaluates as the ratio of the absolute phase limit ($c$) to the localized acoustic limit ($c(r)$):
\begin{equation}
    n(r) = \frac{c}{c(r)} 
\end{equation}

\textbf{2. Acoustic Gravity (Fermat's Principle):} \\
When any external localized wave geometry enters this static Gradient-Index (GRIN) optical lens, its propagation trajectory is governed by Fermat's Principle of Least Time. The external geometric boundary refracts down the continuous pressure gradient toward the localized phase-velocity minimum. Gravitational acceleration evaluates as continuous acoustic refraction.
\begin{equation}
    \nabla P \implies \nabla n \implies \text{Refraction (Gravity)} \tag{\ref{eq:acoustic_gravity}}
\end{equation}

\textbf{3. Dark Matter (Macroscopic GRIN):} \\
A galaxy operates as a massive, synchronized collective of topological drains. This collective generates an extended macroscopic fluid velocity gradient ($v_{rot}$). Light propagating through this moving fluid must satisfy the null geodesic constraint ($ds^2 = 0$). Evaluating the apparent wave velocity ($c'$) against the kinematic drift:
\begin{align}
    0 &= -\left(1 - \frac{v_{rot}^2}{c^2}\right)c^2 dt^2 + \frac{dr^2}{1 - \frac{v_{rot}^2}{c^2}}  \\
    \frac{dr^2}{1 - \frac{v_{rot}^2}{c^2}} &= \left(1 - \frac{v_{rot}^2}{c^2}\right)c^2 dt^2 
\end{align}

Multiplying both sides by the denominator isolates the spatial variance:
\begin{equation}
    \frac{dr^2}{dt^2} = c^2 \left(1 - \frac{v_{rot}^2}{c^2}\right)^2 
\end{equation}

Taking the square root defines the local apparent wave velocity ($c'$):
\begin{align}
    \frac{dr}{dt} &= c \left(1 - \frac{v_{rot}^2}{c^2}\right)  \\
    c' &= c \left(1 - \frac{v_{rot}^2}{c^2}\right) 
\end{align}

The macroscopic refractive index ($n_{gal}(r)$) evaluates as the ratio of the absolute phase limit ($c$) to the local apparent velocity ($c'$). Because it   governs the propagation of transverse shear waves, this index equates natively to the Transverse Impedance ($Z_{\tau(\perp)}$) of the galaxy:
\begin{align}
    Z_{\tau(\perp)} \propto n_{gal}(r) &= \frac{c}{c'}  \\
    n_{gal}(r) &= \frac{c}{c \left( 1 - \frac{v_{rot}^2}{c^2} \right)}  \\
    n_{gal}(r) &= \left( 1 - \frac{v_{rot}^2}{c^2} \right)^{-1} \tag{\ref{eq:dark_matter_grin}}
\end{align}

Because the flat rotation velocity extends far beyond the visible baryonic disc, this extended refractive index physically bends transiting acoustic shear waves (photons) exactly as if a massive matter halo were present. Dark matter evaluates not as particulate mass, but as a macroscopic GRIN lens generated by the collective transverse impedance of the fluid field.

\begin{tcolorbox}[recallbox, title={\ref{box:macroscopic_grin}. Refractive Gravity \& Dark Matter}]
\begin{align}
    c(r) &= \sqrt{\frac{\mu_s(r)}{\rho(r)}} \tag{\ref{eq:acoustic_limit}} \\
    \nabla P &\implies \nabla n \implies \text{Refraction (Gravity)} \tag{\ref{eq:acoustic_gravity}} \\
    n_{gal}(r) &= \left( 1 - \frac{v_{rot}^2}{c^2} \right)^{-1} \tag{\ref{eq:dark_matter_grin}}
\end{align}
\end{tcolorbox}






% ----------------------------------------------------------------------
\section{Temporal Amplification of Mass (Fluid Density Decay)}
\label{sec:deriv_mass_amplification}

The temporal evolution of fundamental rest mass evaluates as the direct mechanical consequence of the unwinding macroscopic metric. As the continuous vacuum expands, its baseline density drops, reducing the hydrodynamic added mass displaced by a localized topological boundary.

\textbf{1. The Hydrodynamic Mass Dependency:} \\
By the Kinematic Origin of Inertia, observable rest mass evaluates as the hydrodynamic added mass of the displaced void ($m = \frac{1}{2}\rho_{\tau} V_{\tau}$). Because the physical volume of a fundamental topological boundary evaluates proportional to the inverse-cube of its mass ($V_{\tau} \propto m^{-3}$), the algebraic relationship between observable mass and ambient metric density isolates to:
\begin{align}
    m &\propto \rho_{\tau} (m^{-3}) \nonumber \\
    m^4 &\propto \rho_{\tau} 
\end{align}

\textbf{2. The Macroscopic Density Decay:} \\
As derived via the invariant acoustic limit, the absolute macroscopic mass density of the vacuum decays as the inverse square of proper time ($\rho_{\tau} \propto t^{-2}$). Because the causal horizon expands linearly ($R_h = c \cdot t$), this density decay equates directly to the macroscopic scale radius:
\begin{equation}
    \rho_{\tau} \propto R_h^{-2} 
\end{equation}
Substituting this macroscopic geometric decay into the localized mass dependency ($m^4 \propto R_h^{-2}$) completely derives the effective bare rest mass as inversely proportional to the square root of the causal horizon:
\begin{align}
    m^4 &\propto R_h^{-2} \nonumber \\
    m &\propto \left(R_h^{-2}\right)^{1/4} \nonumber \\
    m(t) &\propto \frac{1}{\sqrt{R_h(t)}} 
\end{align}

\textbf{3. Amplification at Redshift ($z$):} \\
Evaluating the mass of a localized boundary at redshift $z$ relative to its baseline mass at the present epoch $m_0$:
\begin{align}
    \frac{m(z)}{m_0} &= \frac{1/\sqrt{R_h(z)}}{1/\sqrt{R_{h,0}}} \\
    \frac{m(z)}{m_0} &= \sqrt{ \frac{R_{h,0}}{R_h(z)} }
\end{align}

Substituting the High-Redshift Chronometer scale equation ($R_h(z) = R_{h,0} / (1+z)$) into the denominator:
\begin{align}
    \frac{m(z)}{m_0} &= \sqrt{ \frac{R_{h,0}}{ \frac{R_{h,0}}{1+z} } } \\
    \frac{m(z)}{m_0} &= \sqrt{ \frac{R_{h,0} (1+z)}{R_{h,0}} }
\end{align}

The $R_{h,0}$ parameters cancel, transferring the $(1+z)$ algebraic term cleanly into the numerator:
\begin{align}
    \frac{m(z)}{m_0} &= \sqrt{1+z} \\
    m(z) &= m_0 \sqrt{1+z} \tag{\ref{eq:mass_z}}
\end{align}

\begin{tcolorbox}[recallbox, title={\ref{box:mass_amplification}. Temporal Amplification of Mass}]
\begin{equation}
    m(z) = m_0 \sqrt{1+z} \tag{\ref{eq:mass_z}}
\end{equation}
\end{tcolorbox}




% ==========================================
% PHASE II: VACUUM THERMODYNAMICS & RHEOLOGY
% ==========================================

\part*{PHASE II: VACUUM THERMODYNAMICS \& RHEOLOGY}




% ----------------------------------------------------------------------
\section{Absolute Thermodynamic Vacuum \& Particle Entropy}
\label{sec:deriv_thermo_vacuum}

The baseline temperature of the vacuum evaluates independent of volumetric expansion via the First Law of Thermodynamics: 
\begin{equation}
    T dS = dU + P dV \nonumber
\end{equation}

\textbf{1. The Unruh-Hawking Temperature Floor:} \\
Internal energy evaluates as density times volume $U = \rho V$. Total entropy $S$ evaluates as specific entropy density $\sigma$ times volume $V$, defining $S = \sigma V$. Applying the product rule to the differentials yields:
\begin{align}
    dU &= d(\rho V) = \rho dV + V d\rho  \\
    dS &= d(\sigma V) = \sigma dV + V d\sigma 
\end{align}

Substituting these into the First Law:
\begin{equation}
    T(\sigma dV + V d\sigma) = (\rho dV + V d\rho) + P dV 
\end{equation}

Distributing the temperature $T$ and grouping all terms containing the volumetric expansion factor $dV$ on the right side:
\begin{align}
    V T d\sigma - V d\rho &= \rho dV + P dV - T \sigma dV  \\
    V (T d\sigma - d\rho) &= dV (\rho + P - T \sigma) 
\end{align}

The Euler Equation of State for any continuous extensive medium dictates that $\rho + P = T \sigma$. Therefore, the expression inside the parenthesis evaluates exactly to zero ($\rho + P - T \sigma = 0$). 
\begin{align}
    V (T d\sigma - d\rho) &= dV (0) = 0  \\
    T d\sigma &= d\rho 
\end{align}

The change in volume $dV$ is mathematically annihilated from the thermodynamic evolution. The temperature of the vacuum relies solely on the kinematic acceleration of its uncurling boundary. 

Because the vacuum evaluates as a physical elastic fluid with baseline density $\rho_{\tau}$, this thermal bath evaluates as the stochastic longitudinal acoustic buffeting of the medium. Accelerating a localized topological defect through this fluid subjects it to physical aerodynamic drag. Via the Unruh-Hawking effect, this proper geometric boundary acceleration ($a_0 = c/t$) generates a continuous localized thermodynamic friction floor ($T = \hbar a_0 / 2\pi k_B c$). Substituting the geometric boundary limit $a_0$:
\begin{align}
    T_{cosmic}(t) &= \frac{\hbar (c/t)}{2\pi k_B c}  \\
    T_{cosmic}(t) &= \frac{\hbar}{2\pi k_B t} 
\end{align}
To preserve the acoustic phase symmetry mapped against the proper coordinate limit ($a_0(t) = c^2 / c\cdot t$):
\begin{equation}
    T_{cosmic}(t) = \frac{\hbar c}{2\pi k_B t} \tag{\ref{eq:tcosmic}}
\end{equation}

\textbf{2. Particle Entropy Quantum ($S_p$):} \\
With the localized thermodynamic floor established, the fundamental particle evaluates geometrically as an acoustic standing wave sequestered within the metric. The topological information (Entropy, $S_p$) maintained by this localized cavity evaluates via the fundamental Euler thermodynamic relation:
\begin{equation}
    S_p = \frac{(\rho + P)V}{T} 
\end{equation}

Because the internal trapped standing wave evaluates thermodynamically with a localized hydrostatic fluid pressure $P = \rho/3$, substituting $P$ directly into the equation evaluates the numerator ($\rho + \rho/3 = 4\rho/3$). 
\begin{equation}
    S_p = \frac{\left(\frac{4\rho}{3}\right)V}{T} = \frac{4 \rho V}{3 T} 
\end{equation}

The total internal baseline energy evaluates as the product of fluid density and spherical volume ($\rho V = E_{vac} = \hbar/t$), and the internal thermodynamic Unruh floor evaluates as exactly $T = \hbar / (2\pi k_B t)$. Substituting these absolute quantities into the entropy equation yields:
\begin{align}
    S_p &= \frac{\frac{4}{3} \left(\frac{\hbar}{t}\right)}{\left(\frac{\hbar}{2\pi k_B t}\right)}  \\
    S_p &= \frac{4}{3} \left( \frac{\hbar}{t} \right) \left( \frac{2\pi k_B t}{\hbar} \right) 
\end{align}

The continuous physical parameters $\hbar$ and $t$ algebraically annihilate from the numerator and denominator, proving the entropy evaluates entirely independent of temporal scaling:
\begin{align}
    S_p &= \frac{4}{3} \left( 2\pi k_B \right)  \\
    S_p &= \frac{8\pi}{3} k_B \tag{\ref{eq:sp}}
\end{align}

\begin{tcolorbox}[recallbox, title={\ref{box:thermo_vacuum}. Absolute Thermodynamic Vacuum \& Entropy}]
\begin{align}
    a_0(t) &= \frac{c}{2\pi t} \implies T_{cosmic}(t) = \frac{\hbar c}{2\pi k_B t} \tag{\ref{eq:tcosmic}} \\
    S_p &= \frac{(\rho + P)V}{T} \implies S_p = \frac{8\pi}{3} k_B \tag{\ref{eq:sp}}
\end{align}
\end{tcolorbox}







% ----------------------------------------------------------------------
\section{The Unified Equation of State (Vacuum Energy \& Topological Pressure)}
\label{sec:deriv_unified_eos}

The baseline energy of the macroscopic continuum and the internal phase pressure of a fundamental particle evaluate as symmetric manifestations of a single, continuous equation of state.

\textbf{1. Fundamental Vacuum Energy \& Null Action:} \\
For a continuous acoustic standing wave bounded by the macroscopic scale radius of the universe ($R_h = c \cdot t$), the fundamental geometric wave vector evaluates as the inverse of this boundary ($k = 1/R_h$). The baseline kinematic energy ($E_{vac}$) of this boundary mode evaluates via the transverse acoustic limit ($E = pc$, where $p = \hbar k$):
\begin{align}
    E_{vac}(t) &= \left( \frac{\hbar}{R_h} \right) c  \\
    E_{vac}(t) &= \frac{\hbar c}{c \cdot t} = \frac{\hbar}{t} \tag{\ref{eq:vacuum_null_action}}
\end{align}

To evaluate the internal mechanical strain of this uncurling boundary, the continuous fluid Lagrangian ($\mathcal{L}_{vac}$) evaluates via the Legendre transform of the energy ($pv - E$). Because the macroscopic boundary propagates at $c$, we substitute $v = c$ and $E_{vac} = pc$:
\begin{align}
    \mathcal{L}_{vac}(c) &= p(c) - (pc)  \\
    \mathcal{L}_{vac}(c) &= 0 
\end{align}
Because the continuous geometric action evaluates to exactly zero, the macroscopic vacuum boundary uncurls as an absolute null manifold, accumulating zero structural strain.

\textbf{2. The Topological Equation of State:} \\
When the fluid locally cavitates to form a fundamental particle, its internal thermodynamic state is governed by the continuous Euler relation ($S_p = (\rho + P)V / T$). 

Because the internal trapped standing wave evaluates thermodynamically with a localized internal radiation pressure $P_{int} = \rho/3$, substituting $P_{int}$ directly into the equation evaluates the numerator ($\rho + \rho/3 = 4\rho/3$).

Substituting this into the Euler relation yields $S_p = 4PV/T$. As established via the geometric scaling of the continuum, the absolute geometric entropy of a fundamental topological bit evaluates to exactly $S_p = \frac{8\pi}{3} k_B$. Substituting this constant isolates the state equation:
\begin{align}
    PV &= \frac{1}{4} \left( \frac{8\pi}{3} k_B \right) T  \\
    PV &= \frac{2\pi}{3} k_B T \tag{\ref{eq:unified_topological_eos}}
\end{align}
The discrete arbitrary molar quantity ($n$) utilized in classical macro-thermodynamics evaluates for a fundamental particle as exactly $2\pi/3$, dictated entirely by the $2\pi$ phase wrap bounding the 3D spherical fluid volume.

\textbf{3. Rest Mass Equivalence \& Bag Pressure:} \\
For a radiation-dominated acoustic cavity, internal thermal energy evaluates as $U = \rho V = 3PV$. Substituting the derived Equation of State:
\begin{align}
    U &= 3 \left( \frac{2\pi}{3} k_B T \right)  \\
    U &= 2\pi k_B T 
\end{align}
Substituting the geometric acoustic temperature of a localized boundary ($T = \frac{mc^2}{2\pi k_B}$) isolates the absolute rest-mass volumetric energy ($U_{topo} = mc^2$). 

Furthermore, isolating the internal radiation pressure from this energy equation ($P_{int} = \frac{m c^2}{3 V}$) and substituting the exact 3D spherical Euclidean bounding volume ($V = \frac{4}{3}\pi (\hbar/mc)^3$) yields:
\begin{align}
    P_{int} &= \frac{m^4 c^5}{4\pi \hbar^3}  \\
    P_{int} &= 2\pi P_p \tag{\ref{eq:unified_pressure_ratio}}
\end{align}
The internal 3D volumetric fluid pressure of the particle evaluates identically to the 2D surface confinement Bag Pressure ($P_p$) multiplied by one full orthogonal geometric phase wrap ($2\pi$).

\textbf{4. Scale Invariance (The Cosmic Loop):} \\
To prove the absolute scale invariance of the metric, we apply the localized particle energy equation ($U = 2\pi k_B T$) to the macroscopic cosmos. Substituting the geometric temperature floor of the expanding universe ($T = \frac{\hbar}{2\pi k_B t}$) yields:
\begin{align}
    U_{vac} &= 2\pi k_B \left( \frac{\hbar}{2\pi k_B t} \right)  \\
    U_{vac} &= \frac{\hbar}{t} \tag{\ref{eq:vacuum_null_action}}
\end{align}
The continuous fluid-dynamic equation of state flawlessly recovers the exact Fundamental Vacuum Energy derived purely from the macroscopic horizon boundary in Step 1.

\begin{tcolorbox}[recallbox, title={\ref{box:unified_eos}. The Unified Equation of State}]
\begin{align}
    E_{vac}(t) &= \frac{\hbar}{t} \implies \mathcal{L}_{vac}(c) = 0 \tag{\ref{eq:vacuum_null_action}} \\
    PV &= \frac{2\pi}{3} k_B T \implies U = 2\pi k_B T \tag{\ref{eq:unified_topological_eos}} \\
    U_{topo} &= mc^2 \implies P_{int} = 2\pi P_p \tag{\ref{eq:unified_pressure_ratio}}
\end{align}
\end{tcolorbox}







% ----------------------------------------------------------------------
\section{The Continuous Rheological Gradient}
\label{sec:deriv_rheological_gradient}

The volumetric expansion of the macroscopic universe and the localized mechanics of fundamental matter resolve via a continuous rheological gradient governed by the localized fluid velocity field ($v(r)$). 

\textbf{1. The Elastic Relation:} \\
The Poisson ratio ($\nu$) evaluates the geometric fraction of transverse expansion to longitudinal extension. It is algebraically defined by the localized Bulk Modulus ($K$) and the invariant Shear Modulus ($\mu_s$):
\begin{equation}
    \nu(r) = \frac{3K(r) - 2\mu_s}{2(3K(r) + \mu_s)} 
\end{equation}

\textbf{2. The Macroscopic Regime (Auxetic Generation):} \\
In the quiescent macroscopic vacuum ($r \to \infty$), fluid velocity evaluates as negligible ($v \to 0$). As derived via the vacuum compressibility limit (Eq. 56), the baseline isentropic Bulk Modulus evaluates as $K = \frac{1}{3}\rho_{\tau} c^2$. Substituting the absolute shear modulus ($\mu_s = \rho_{\tau} c^2$) into the classical relation derives the absolute geometric behavior of the background metric:
\begin{align}
    \nu_{macro} &= \frac{3(\frac{1}{3}\mu_s) - 2\mu_s}{2(3(\frac{1}{3}\mu_s) + \mu_s)}  \\
    \nu_{macro} &= \frac{-\mu_s}{4\mu_s} = -0.25 \tag{\ref{eq:poisson_macro}}
\end{align}
The continuous vacuum evaluates as an Auxetic solid, generating Eulerian spatial volume when subjected to isotropic causal tension.

\textbf{3. The Macroscopic Boundary (Time Crystal Compliance):} \\
Standard fundamental mass (e.g., the electron) does not evaluate as a monolithic $l_p$ singularity; it is a macroscopic topological standing wave bounded by the Compton wavelength ($R_h = \lambda_c$). 

At the Compton boundary, the internal fluid circulation ($v$) is elevated, causing a proportional kinematic stiffening of the localized Bulk Modulus ($K > \frac{1}{3}\mu_s$). However, because $v \ll c$ at the macroscopic boundary, the fluid does not kinematically choke. 
\begin{equation}
    0 \le \nu(\lambda_c) < 0.5 \tag{\ref{eq:poisson_compton}}
\end{equation}
This intermediate, non-auxetic positive Poisson ratio provides the exact structural compliance required for the Time Crystal phase-ratchet to operate. The boundary is mechanically capable of flexing, yielding, and continuously unspooling fluid volume without fracturing.

\textbf{4. The Localized Core (The Incompressible Anchor):} \\
As the fluid circulates deep inside the topological knot toward the absolute causal core ($r \to l_p$), the fluid velocity strikes the acoustic phase limit ($v \to c$).

Because the fluid cannot accelerate past $c$, its resistance to further volumetric compression diverges to infinity ($\lim_{v \to c} K = \infty$). Dividing the Poisson algebraic limit by $K$:
\begin{align}
    \nu_{core} &= \lim_{K \to \infty} \frac{3 - 2\frac{\mu_s}{K}}{2(3 + \frac{\mu_s}{K})}  \\
    \nu_{core} &= \frac{3}{6} = 0.5 \tag{\ref{eq:poisson_local}}
\end{align}
At exactly $\nu = 0.5$, the fluid becomes perfectly incompressible ($\nabla \cdot \mathbf{u} = 0$). This unbreakable core acts as the absolute mechanical anchor for the topological defect, preventing the structure from completely collapsing into the Euclidean void under extreme external macroscopic pressure.

\begin{tcolorbox}[recallbox, title={\ref{box:rheological_gradient}. The Rheological Gradient}]
\begin{align}
    \nu_{macro} &= -0.25 \tag{\ref{eq:poisson_macro}} \\
    \nu(\lambda_c) &\in [0, 0.5) \tag{\ref{eq:poisson_compton}} \\
    \nu_{core} &= 0.5 \implies \nabla \cdot \mathbf{u} = 0 \tag{\ref{eq:poisson_local}}
\end{align}
\end{tcolorbox}

% ----------------------------------------------------------------------
\section{The Unified Cosmological Deborah Number}
\label{sec:deriv_unified_deborah}

The macroscopic metric evaluates as a viscoelastic continuum. The rheological phase response of the vacuum and the relative strengths of the fundamental forces evaluate via a single geometric ratio dictated by the energetic magnitude of the interaction.

\textbf{1. The Predictive Equation (The Baseline Yield Limit):} \\
The Cosmological Deborah Number ($De$) evaluates as the product of the continuous angular frequency ($\omega$) of an applied kinematic wave and the intrinsic relaxation time ($t_r$) of the metric:
\begin{equation}
    De = \omega t_r \nonumber
\end{equation}

By the quantum phase relation, the continuous angular frequency evaluates as the interaction energy ($E$) divided by the reduced Planck constant ($\hbar \omega = E \implies \omega = E/\hbar$). 

In classical continuum mechanics, intrinsic relaxation time ($t_r$) evaluates via the macroscopic geometric yield limit of the fluid. The metric possesses an absolute transverse spatial yield limit ($R$), beyond which the continuous metric cannot sustain structural phase without permanently exhausting transverse shear. 

Because stable localized boundaries evaluate geometrically via their Compton limit ($R = \hbar/mc$), this absolute spatial yield limit mathematically defines the fundamental baseline mass of the continuous vacuum ($\bar{\lambda}_e = \hbar/m_e c$). The empirical electron ($m_e$) evaluates as the exact mass generated at this absolute geometric yield boundary of the metric.

The intrinsic relaxation time of the continuum evaluates as this spatial yield limit ($\bar{\lambda}_e$) scaled by the transverse elastic compliance ($\alpha$) divided by the phase velocity ($c$):
\begin{equation}
    t_r = \frac{\alpha R_e}{c} = \frac{\alpha \left(\frac{\hbar}{m_e c}\right)}{c} = \frac{\alpha \hbar}{m_e c^2} 
\end{equation}

Substituting these continuous parameters into the Deborah Number equation yields:
\begin{equation}
    De = \left( \frac{E}{\hbar} \right) \left( \frac{\alpha \hbar}{m_e c^2} \right) 
\end{equation}

Rearranging the terms isolates the dimensional constants from the interaction variable. The reduced Planck constant ($\hbar$) geometrically annihilates from the numerator and the denominator. The Deborah Number evaluates explicitly as the localized interaction energy scaled against the absolute topological yield baseline of the continuum:
\begin{equation}
    De = \alpha \left( \frac{E}{m_e c^2} \right) \tag{\ref{eq:unified_deborah}}
\end{equation}

\textbf{2. Complex Dynamic Modulus ($G^*$):} \\
The frequency-dependent mechanical response of the continuum evaluates by applying the Laplace Transform to the Maxwell viscoelastic equation ($s \varepsilon = s \sigma / \mu_s + \sigma / \eta$). Isolating the ratio of stress to strain defines the Complex Dynamic Modulus ($G^*$). 

Substituting the harmonic oscillation limit ($s = i\omega$) and the unified Deborah Number defines the real elastic Storage Modulus ($G'$) and the imaginary viscous Loss Modulus ($G''$):
\begin{equation}
    G^*(E) = \mu_s \left[ \frac{De^2}{1 + De^2} \right] + i \mu_s \left[ \frac{De}{1 + De^2} \right] \tag{\ref{eq:deborah_modulus}}
\end{equation}

\textbf{3. The Viscoelastic Phase Angle (Angular Distribution):} \\
The phase angle ($\delta$) measures the geometric partition between the dissipated fluid energy (the Loss Modulus, $G''$) and the stored elastic energy (the Storage Modulus, $G'$). The loss tangent evaluates as their quotient:
\begin{equation}
    \tan \delta = \frac{G''}{G'} = \frac{\frac{De}{1+De^2}}{\frac{De^2}{1+De^2}} = \frac{1}{De}
\end{equation}

The fundamental forces evaluate as this kinematic partition. The scalar fluid compliance fraction ($\alpha_s = g^2/4\pi$) evaluates as the viscous fluid response ($G''$). The transverse electromagnetic coupling fraction ($\alpha$) evaluates as the rigid elastic response ($G'$). 

The angular distribution between the strong force and the electromagnetic force evaluates as the viscoelastic phase angle of the continuum. Substituting the unified equation isolates the dynamic coupling ratio:
\begin{align}
    \frac{\alpha_s}{\alpha} &= \tan \delta = \frac{1}{De} \\
    \frac{\alpha_s}{\alpha} &= \frac{m_e c^2}{\alpha E} \tag{\ref{eq:angular_distribution}}
\end{align}

\textbf{4. Confinement \& Asymptotic Freedom:} \\
At the topological boundary evaluation limit ($E = m_e c^2$), the ratio evaluates to $\alpha_s / \alpha \approx 1/\alpha$. Multiplying by $\alpha$ yields $\alpha_s \approx 1$. The macroscopic strong coupling parameter isolates algebraically to its bare scalar limit ($g^2 \approx 4\pi$), defining Color Confinement. 

As interaction energy scales toward infinity ($E \to \infty$), the Deborah Number diverges ($De \gg 1$), and the phase angle collapses ($\tan \delta \to 0$). The scalar fluid response vanishes, forcing the strong coupling fraction to zero ($\alpha_s \to 0$). The metric freezes into a rigid solid, deriving Asymptotic Freedom. The apparent strengths of the fundamental forces evaluate as the frequency-dependent phase angle of the interacting fluid.

\textbf{5. The Glass Transition (Metric Fracture):} \\
To evaluate the energy required to flash-freeze the vacuum into a rigid solid, the Deborah Number evaluates to unity ($De = 1$). 
\begin{align}
    1 &= \alpha \left( \frac{E_{crit}}{m_e c^2} \right) \\
    E_{crit} &= \frac{m_e c^2}{\alpha}
\end{align}

By evaluating this limit using the empirical electron mass ($0.511 \text{ MeV}$), the yield threshold evaluates to $E_{crit} \approx 70.02 \text{ MeV}$. Because topological hadronization nucleates boundaries in symmetric pairs, the total metric cavitation energy requires $2 E_{crit} \approx 140.04 \text{ MeV}$. 

This isolates the empirical rest mass of the charged pion ($m_{\pi^{\pm}} \approx 139.57 \text{ MeV}$). At interaction energies exceeding this limit, the vacuum metric shatters under the kinematic strain, initiating spontaneous topological mass generation.
\begin{equation}
    De \to 1 \implies 2 E_{crit} = 140.04 \text{ MeV} \equiv m_{\pi^{\pm}} \tag{\ref{eq:de_glass}}
\end{equation}

\textbf{6. Electro-Strong Phase Symmetry ($De = 1$):} \\
At the glass transition limit ($De = 1$), substituting the unity parameter into the continuous angular distribution equation evaluates the unification limit of the metric forces:
\begin{align}
    \frac{\alpha_s}{\alpha} &= \frac{1}{1} \\
    \alpha_s &= \alpha \tag{\ref{eq:electro_strong_symmetry}}
\end{align}

Evaluating the phase angle isolates $\tan \delta = 1 \implies \delta = \pi/4$. At the critical fracture threshold of the vacuum ($E_{crit} = m_e c^2 / \alpha$), the scalar fluid admittance equals the transverse elastic storage. The Strong interaction and the Electromagnetic interaction evaluate as symmetric geometric projections of the yielding continuum.

\textbf{7. The Volumetric State Multiplier ($N$):} \\
Defining the localized interaction energy ($E$) as an integer multiple ($N$) of the baseline topological rest mass ($E = N m_e c^2$), the continuous strong coupling fraction isolates the physical quantity of nucleated boundaries:
\begin{align}
    \alpha_s &= \frac{m_e c^2}{N m_e c^2} \\
    \alpha_s &= \frac{1}{N} \tag{\ref{eq:volumetric_multiplier}}
\end{align}

The scalar fluid coupling evaluates as the inverse geometric state multiplier of the topological volume. Evaluating this geometric fraction isolates the sequential cavitation thresholds of the continuum:
\begin{itemize}
    \item \textbf{Ground State ($N=1$):} The fraction evaluates to $\alpha_s = 1$. The metric supports a single volumetric boundary.
    \item \textbf{Bipartite Yield ($N=2$):} Substituting the minimum energy required to initiate kinematic cavitation ($E = 2m_e c^2$) forces $\alpha_s = 1/2$. The scalar fluid admittance drops. To conserve Eulerian volume, the metric bifurcates, nucleating two symmetric boundaries.
    \item \textbf{The Shatter Limit ($N \approx 137$):} At the glass transition ($E = m_e c^2 / \alpha$), the multiplier evaluates to $N = 1/\alpha$. The metric attempts to bifurcate into $137$ boundaries simultaneously. Incapable of sustaining cohesive structure, the vacuum shatters into a shower of secondary hadronic topologies.
\end{itemize}

\textbf{8. The Differential Yield Limit (Continuum Beta Function):} \\
The rate at which the vacuum metric stiffens against applied energy evaluates as the first derivative of the strong fluid fraction with respect to the continuous interaction energy ($E$):
\begin{align}
    \frac{d\alpha_s}{dE} &= \frac{d}{dE} \left( \frac{m_e c^2}{E} \right) \\
    \frac{d\alpha_s}{dE} &= - \frac{m_e c^2}{E^2}
\end{align}

Substituting the volumetric state multiplier ($E = N m_e c^2$) evaluates this rate of closure in the discrete topological domain:
\begin{align}
    \frac{d\alpha_s}{dN} &= \frac{d}{dN} \left( \frac{1}{N} \right) \\
    \frac{d\alpha_s}{dN} &= - \frac{1}{N^2} \tag{\ref{eq:beta_state}}
\end{align}

This negative differential limit defines the continuum mechanics of Asymptotic Freedom. The negative sign dictates that the scalar fluid admittance collapses as energy rises. 

\textbf{9. The Continuum Beta Function ($\beta$):} \\
The formal scale variance of the strong interaction evaluates as the logarithmic derivative of the fluid coupling with respect to the applied interaction energy ($E$). This defines the Continuum Beta Function ($\beta(\alpha_s)$):
\begin{equation}
    \beta(\alpha_s) = \frac{\partial \alpha_s}{\partial \ln E}
\end{equation}

By the chain rule, this logarithmic derivative expands to:
\begin{equation}
    \beta(\alpha_s) = E \frac{\partial \alpha_s}{\partial E}
\end{equation}

Substituting the derived linear differential limit ($\partial \alpha_s / \partial E = - m_e c^2 / E^2$):
\begin{align}
    \beta(\alpha_s) &= E \left( - \frac{m_e c^2}{E^2} \right) \\
    \beta(\alpha_s) &= - \frac{m_e c^2}{E}
\end{align}

Because the quotient isolates the baseline scalar fluid fraction ($\alpha_s = m_e c^2 / E$), the continuum beta function evaluates algebraically to a closed-form limit:
\begin{equation}
    \beta(\alpha_s) = - \alpha_s \tag{\ref{eq:beta_function}}
\end{equation}

The geometric stiffening of the metric evaluates as the negative magnitude of its continuous fluid admittance. 

\textbf{10. Vacuum Relaxation (The Nuclear Range):} \\
At the particle boundary, the spatial range of the strong interaction ($L_{crit}$) evaluates as the kinematic propagation of strain before viscoelastic relaxation ($L = c \cdot t_r$). Substituting $t_r$ isolates the Classical Electron Radius:
\begin{align}
    L_{crit} &= c \left( \frac{\alpha \hbar}{m_e c^2} \right) \\
    L_{crit} &= \alpha \frac{\hbar}{m_e c} \equiv r_e \approx 2.81 \text{ fm} \tag{\ref{eq:de_range}}
\end{align}
This geometric product proves that the Classical Electron Radius ($r_e$) and the spatial limit of nuclear confinement are the macroscopic acoustic relaxation limit of the metric.

\textbf{11. The Macroscopic Drag Limit (The Hierarchy Problem):} \\
By evaluating the Deborah Number at the kinematic minimum, the macroscopic ratio of elastic to viscous resistance isolates the baseline strength of the ambient gravitational pressure gradient. For a macroscopic translating mass, the maximum observation time evaluates as the proper age of the universe ($t_{obs} = R_{h,0} / c$).
\begin{equation}
    De_{min} = \frac{t_r}{t_{obs}} = \frac{\alpha \hbar / m_e c^2}{R_{h,0} / c} = \alpha \left( \frac{\hbar}{m_e c R_{h,0}} \right)
\end{equation}

By the Volumetric Inverse-Cube Identity, the localized particle boundary evaluates as $R_{h,p} = \hbar / m_e c$. Substituting this spatial boundary isolates the geometric scalar of the Hierarchy Problem:
\begin{equation}
    De_{min} = \alpha \left( \frac{R_{h,p}}{R_{h,0}} \right) \approx 10^{-41} \tag{\ref{eq:de_gravity}}
\end{equation}
Because the minimum fluid Deborah Number evaluates at $10^{-41}$, the macroscopic metric yields to kinematic translation. The elastic resistance of the vacuum evaluates as $10^{41}$ times weaker than its viscous response.

\begin{tcolorbox}[recallbox, title={\ref{box:unified_deborah}. The Unified Cosmological Deborah Number}]
\begin{align}
    De &= \alpha \left( \frac{E}{m_e c^2} \right) \tag{\ref{eq:unified_deborah}} \\
    G^*(E) &= \mu_s \left[ \frac{De^2}{1 + De^2} \right] + i \mu_s \left[ \frac{De}{1 + De^2} \right] \tag{\ref{eq:deborah_modulus}} \\
    \frac{\alpha_s}{\alpha} &= \frac{m_e c^2}{\alpha E} \tag{\ref{eq:angular_distribution}} \\
    De \to 1 &\implies 2 E_{crit} = 140.04 \text{ MeV} \equiv m_{\pi^{\pm}} \tag{\ref{eq:de_glass}} \\
    De = 1 &\implies \alpha_s = \alpha \tag{\ref{eq:electro_strong_symmetry}} \\
    E = N m_e c^2 &\implies \alpha_s = \frac{1}{N} \tag{\ref{eq:volumetric_multiplier}} \\
    \frac{d\alpha_s}{dN} &= - \frac{1}{N^2} \tag{\ref{eq:beta_state}} \\
    \beta(\alpha_s) &= - \alpha_s \tag{\ref{eq:beta_function}} \\
    L_{crit} &= r_e \approx 2.81 \text{ fm} \tag{\ref{eq:de_range}} \\
    De \to 0 &\implies De_{min} \approx 10^{-41} \tag{\ref{eq:de_gravity}}
\end{align}
\end{tcolorbox}




% ----------------------------------------------------------------------
\section{Viscoelastic Phase Stiffening \& Time Dilation}
\label{sec:deriv_viscoelastic_stiffening}

The rheological stiffening of the vacuum, the retardation of proper time, and the kinematic depletion of fluid admittance evaluate as conserved fluid-dynamic limits.

\textbf{1. The Compliance Limit ($\alpha_s$):} \\
The total scalar admittance ($\alpha_s$) evaluates as the square root of the residual static capacity ratio ($P_{static} / P_c$). Because the continuous fluid conserves total yield capacity ($P_{static} = P_c - P_{dyn}$), substituting the kinematic definitions isolates the exact compliance limit:
\begin{align}
    \alpha_s &= \sqrt{\frac{P_{static}}{P_c}} \nonumber \\
    \alpha_s &= \sqrt{\frac{P_c - P_{dyn}}{P_c}} \nonumber \\
    \alpha_s &= \sqrt{\frac{\frac{1}{2}\rho_{\tau} c^2 - \frac{1}{2}\rho_{\tau} v^2}{\frac{1}{2}\rho_{\tau} c^2}} \nonumber \\
    \alpha_s &= \sqrt{1 - \frac{v^2}{c^2}} \nonumber \\
    \alpha_s^2 &= 1 - \frac{v^2}{c^2} \equiv 1 - \frac{P_{dyn}}{P_c} \tag{\ref{eq:admittance_pressure_ratio}}
\end{align}

\textbf{2. The Yukawa-Bernoulli Identity:} \\
The dynamic pressure of the circulating fluid evaluates as the localized interaction strength, mechanically identical to the squared Yukawa coupling constant ($g^2 \propto P_{dyn}$). The absolute acoustic yield pressure equates to the theoretical maximum coupling limit ($g_{max}^2 \propto P_c$).

Substituting this normalization into the compliance equation derives the coupling identity:
\begin{equation}
    \alpha_s^2 = 1 - \frac{g^2}{g_{max}^2} \tag{\ref{eq:yukawa_identity}}
\end{equation}
The strong coupling constant ($\alpha_s$) evaluates as the residual structural compliance of the metric fluid.

\textbf{3. Relativistic Packing \& The Lorentz Pressure Limit:} \\
As a topological boundary translates through the continuous metric, macroscopic kinematic velocity forces the surrounding fluid to undergo spatial compression. The effective local density of the $\boldsymbol{\tau}$-fluid ($\rho_{eff}$) evaluates via the Lorentz factor ($\gamma$).

Because $\gamma = 1/\alpha_s$, the geometric packing derives directly from the dynamic pressure ratio:
\begin{align}
    \rho_{eff} &= \rho_{\tau} \gamma \tag{\ref{eq:relativistic_packing_admittance}} \\
    \gamma &= \sqrt{\frac{P_c}{P_c - P_{dyn}}} \tag{\ref{eq:lorentz_pressure}}
\end{align}
Time dilation evaluates as the physical degradation of available static structural capacity. As dynamic pressure consumes the structural capacity of the continuum, acoustic waves require more absolute phase cycles to transit the densified coordinate space.

\textbf{4. The Kinematic Phase Lock (Pythagorean Substitution):} \\
By the conservation of continuous phase capacity, the absolute acoustic limit ($c$) dictates the geometric tradeoff between longitudinal translation ($v_{\parallel}$) and internal topological vorticity ($v_{\circlearrowleft}$) via $c^2 = v_{\parallel}^2 + v_{\circlearrowleft}^2$.

Algebraically substituting the longitudinal velocity out of the Lorentz factor establishes the macroscopic strain exclusively via the internal rotational phase ($\gamma = c/v_{\circlearrowleft}$). Applying the admittance identity ($\gamma = 1/\alpha_s$) isolates the definition of fluid yield:
\begin{equation}
    \alpha_s = \frac{v_{\circlearrowleft}}{c} \tag{\ref{eq:admittance_spin_identity}}
\end{equation}
The total scalar admittance ($\alpha_s$) evaluates as the ratio of internal topological vorticity to the absolute acoustic limit. Macroscopic strain evaluates as the continuous depletion of this internal rotational capacity.

\textbf{5. Viscoelastic Time Dilation (The Deborah Limit):} \\
Substituting the localized kinematic strain energy ($E = \gamma m_e c^2$) into the Unified Cosmological Deborah Number ($De = \alpha E / m_e c^2$), the baseline rest mass algebraically annihilates:
\begin{equation}
    De = \alpha \gamma \tag{\ref{eq:kinematic_deborah}}
\end{equation}
The internal harmonic phase of a localized boundary retards because the localized fluid Deborah Number increases, indicating the metric lacks the structural compliance to process the phase rotation.

\textbf{6. Asymptotic Freedom and Solidification:} \\
As longitudinal translation accelerates toward the acoustic limit ($v_{\parallel} \to c$), the localized internal vorticity chokes to zero ($v_{\circlearrowleft} \to 0$) and the dynamic pressure strikes the absolute yield ceiling ($P_{dyn} \to P_c$). The normalized coupling maximizes ($g^2/g_{max}^2 \to 1$).

The Cosmological Deborah Number diverges to infinity:
\begin{align}
    \lim_{P_{dyn} \to P_c} \alpha_s^2 &= 1 - 1 = 0 \\
    \lim_{v_{\parallel} \to c} De &= \infty \implies \alpha_s = 0, \quad \nabla \cdot \mathbf{v} = 0 \tag{\ref{eq:cavitation_collapse}}
\end{align}
At this absolute limit, the surrounding metric undergoes a localized glass transition. The total scalar fluid response collapses ($\alpha_s = 0$), and the spatial coordinate flash-freezes into a rigid elastic solid. In continuum aerodynamics, this evaluates as "choked flow," where the spatial divergence of the fluid evaluates to exactly zero ($\nabla \cdot \mathbf{v} = 0$). 

The internal circulation encounters zero kinematic resistance (Asymptotic Freedom) because it is walled in by the infinite effective pressure barrier of the surrounding metric. The absolute cosmic speed limit ($c$) evaluates as the strict mechanical yield point where the continuous $\boldsymbol{\tau}$-fluid solidifies.

\begin{tcolorbox}[recallbox, title={\ref{box:viscoelastic_stiffening}. Viscoelastic Phase Stiffening \& Time Dilation}]
\begin{align}
    \alpha_s^2 &= 1 - \frac{P_{dyn}}{P_c} \tag{\ref{eq:admittance_pressure_ratio}} \\
    \alpha_s^2 &= 1 - \frac{g^2}{g_{max}^2} \tag{\ref{eq:yukawa_identity}} \\
    \rho_{eff} &= \rho_{\tau} \gamma \tag{\ref{eq:relativistic_packing_admittance}} \\
    \gamma &= \sqrt{\frac{P_c}{P_c - P_{dyn}}} \tag{\ref{eq:lorentz_pressure}} \\
    \alpha_s &= \frac{v_{\circlearrowleft}}{c} \tag{\ref{eq:admittance_spin_identity}} \\
    De &= \alpha \gamma \tag{\ref{eq:kinematic_deborah}} \\
    \lim_{v_{\parallel} \to c} De &= \infty \implies \alpha_s = 0, \quad \nabla \cdot \mathbf{v} = 0 \tag{\ref{eq:cavitation_collapse}} 
\end{align}
\end{tcolorbox}


% ----------------------------------------------------------------------
\section{The Kinematic Impedance Tensor ($[\boldsymbol{Z}_{\tau}]$)}
\label{sec:deriv_universal_impedance}

Inertial mass, optical density, quantum electrical resistance, and quantized spin evaluate as the diagonal geometric limits of a single continuous matrix: the Kinematic Impedance Tensor ($[\boldsymbol{Z}_{\tau}]$). Because absolute fluid resistance evaluates   as the mathematical inverse of fluid compliance ($[\boldsymbol{Z}_{\tau}] = [\boldsymbol{\alpha}]^{-1}$), the universal impedance projects into physical mechanics via its specific directional limits ($Z_i \equiv 1/\alpha_i$).

\textbf{1. The Transverse Quantum Limit ($Z_{\tau(\perp)}$):} \\
Kinematic impedance evaluates as the ratio of applied force to the resulting continuous volumetric flux. At the fundamental topological boundary, the applied geometric force evaluates as the absolute phase torque limit ($h$). The fluid response evaluates as the squared transverse geometric phase flux ($q^2$).

The baseline transverse kinematic impedance of the continuous vacuum ($Z_{\tau(\perp)}$) evaluates as the algebraic quotient:
\begin{equation}
    Z_{\tau(\perp)} = \frac{h}{q^2} \equiv R_K \tag{\ref{eq:von_klitzing_impedance}}
\end{equation}
The quantum of electrical resistance evaluates as the bare geometric impedance of the $\boldsymbol{\tau}$-fluid to propagating transverse shear.

\textbf{2. The Longitudinal Inertial Limit ($Z_{\tau(\parallel)}$):} \\
A transverse acoustic wave translating through a periodic fluid boundary structure experiences continuous kinematic impedance. The macroscopic group velocity ($v_g$) of the wave evaluates as the spatial derivative of the continuous angular frequency ($\omega$) with respect to the geometric wave vector ($k$):
\begin{equation}
    v_g = \frac{d\omega}{dk} 
\end{equation}

Substituting the continuous phase energy relation ($E = \hbar \omega \implies \omega = E/\hbar$):
\begin{align}
    v_g &= \frac{d}{dk} \left( \frac{E}{\hbar} \right)  \\
    v_g &= \frac{1}{\hbar} \frac{dE}{dk} 
\end{align}

The effective acceleration ($a$) evaluates as the explicit time derivative of this group velocity ($a = \frac{dv_g}{dt}$). Applying the chain rule transitions the derivative from temporal to spatial coordinates:
\begin{align}
    a &= \frac{dv_g}{dt}  \\
    a &= \left( \frac{dv_g}{dk} \right) \frac{dk}{dt} 
\end{align}

Substituting the spatial derivative of the group velocity ($\frac{dv_g}{dk} = \frac{1}{\hbar} \frac{d^2 E}{dk^2}$):
\begin{equation}
    a = \left( \frac{1}{\hbar} \frac{d^2 E}{dk^2} \right) \frac{dk}{dt} 
\end{equation}

External force ($F$) applied to the metric evaluates as the temporal rate of change of the structural phase momentum ($p = \hbar k$):
\begin{align}
    F &= \frac{dp}{dt}  \\
    F &= \frac{d(\hbar k)}{dt}  \\
    F &= \hbar \frac{dk}{dt} \implies \frac{dk}{dt} = \frac{F}{\hbar} 
\end{align}

Substituting this momentum identity into the continuous acceleration equation:
\begin{align}
    a &= \left( \frac{1}{\hbar} \frac{d^2 E}{dk^2} \right) \left( \frac{F}{\hbar} \right)  \\
    a &= F \left( \frac{1}{\hbar^2} \frac{d^2 E}{dk^2} \right) 
\end{align}

Newtonian kinematics equates force to mass and acceleration ($F = m^* a \implies a = F/m^*$). Setting the kinematic equations equal isolates the apparent geometric mass:
\begin{align}
    \frac{F}{m^*} &= F \left( \frac{1}{\hbar^2} \frac{d^2 E}{dk^2} \right)  \\
    \frac{1}{m^*} &= \frac{1}{\hbar^2} \frac{d^2 E}{dk^2} 
\end{align}

Inverting both sides extracts the longitudinal directional impedance:
\begin{equation}
    Z_{\tau(\parallel)} \equiv m^* = \hbar^2 \left( \frac{d^2 E}{dk^2} \right)^{-1} \tag{\ref{eq:effective_mass_impedance}}
\end{equation}
A particle possesses apparent inertial mass ($m^*$) because the localized fluid boundary structurally resists longitudinal displacement within the continuum.

\textbf{3. The Transverse Refractive Limit ($Z_{\tau(\perp)}$):} \\
The macroscopic refractive index ($n$) evaluates the kinematic impedance of the macroscopic velocity field ($v_0$), dictating the path of an acoustic shear wave along a null geodesic ($ds^2 = 0$). Setting the metric tensor to zero:
\begin{align}
    -\left(1 - \frac{v_0^2}{c^2}\right)c^2 dt^2 + \frac{dr^2}{1 - \frac{v_0^2}{c^2}} &= 0  \\
    \frac{dr^2}{1 - \frac{v_0^2}{c^2}} &= \left(1 - \frac{v_0^2}{c^2}\right)c^2 dt^2  \\
    \frac{dr^2}{dt^2} &= c^2 \left(1 - \frac{v_0^2}{c^2}\right)^2 
\end{align}

Taking the square root yields the local wave velocity ($c'$):
\begin{equation}
    c' = \frac{dr}{dt} = c \left(1 - \frac{v_0^2}{c^2}\right) 
\end{equation}

The refractive impedance ($n(r)$) to transverse waves evaluates as the ratio of the absolute acoustic limit ($c$) to the local velocity ($c'$):
\begin{align}
    Z_{\tau(\perp)} \propto n(r) &= \frac{c}{c'}  \\
    n(r) &= \frac{c}{c \left(1 - \frac{v_0^2}{c^2}\right)}  \\
    n(r) &= \left(1 - \frac{v_0^2}{c^2}\right)^{-1} \tag{\ref{eq:refractive_impedance}}
\end{align}

The absolute velocity field ($v_0$) of a localized acoustic cavity evaluates by integrating the continuous convective acceleration ($a_{conv} = v_0 \frac{dv_0}{dr}$) of the fluid yielding to the inverse-square macroscopic pressure deficit:
\begin{align}
    v_0 \frac{dv_0}{dr} &= -\frac{GM}{r^2}  \\
    \int v_0 \, dv_0 &= -GM \int \frac{1}{r^2} \, dr  \\
    \frac{1}{2}v_0^2 &= \frac{GM}{r}  \\
    v_0^2 &= \frac{2GM}{r} 
\end{align}

Substituting this velocity limit into the binomial expansion $(1-x)^{-1} \approx 1+x$ yields the explicit geometric spatial curvature (index of refraction):
\begin{equation}
    n(r) \approx 1 + \frac{2GM}{c^2 r} 
\end{equation}

The energy ($E$) of a wave refracting through this densified pressure gradient evaluates as $E(r) = \frac{mc^2}{n(r)}$. Applying the inverse binomial expansion $n(r)^{-1} \approx 1 - \frac{2GM}{c^2 r}$:
\begin{align}
    E(r) &\approx mc^2 \left( 1 - \frac{GM}{c^2 r} \right)  \\
    E(r) &\approx mc^2 - \frac{GMm}{r} 
\end{align}

Apparent Newtonian force evaluates as the negative spatial gradient of this real thermodynamic strain ($\mathbf{F} = -\nabla E$):
\begin{align}
    \mathbf{F}_g &= -\frac{d}{dr} \left( mc^2 - \frac{GMm}{r} \right) \mathbf{\hat{r}}  \\
    \mathbf{F}_g &= -\left( 0 - \frac{GMm (-1)}{r^2} \right) \mathbf{\hat{r}}  \\
    \mathbf{F}_g &= -\frac{GMm}{r^2} \mathbf{\hat{r}} 
\end{align}
Apparent gravitational trapping, effective rest mass, and electrical resistance evaluate mechanically as isomorphic directional continuum impedances.

\textbf{4. The Rotational Phase Limit ($\hbar$ \& Gyroscopic Inertia):} \\
Rotational kinematic impedance ($Z_{\tau(\circlearrowleft)}$) evaluates the physical resistance of the continuous metric to localized topological vorticity ($\nabla \times \mathbf{v}_{\psi}$). In classical rotational mechanics, impedance evaluates as the ratio of applied rotational energy (torque, $E$) to the resulting continuous angular phase velocity ($\omega$).
\begin{equation}
    Z_{\tau(\circlearrowleft)} = \frac{E}{\omega}
\end{equation}

Substituting the absolute continuous quantum phase relation ($E = \hbar\omega$) geometrically isolates the baseline rotational fluid resistance:
\begin{align}
    Z_{\tau(\circlearrowleft)} &= \frac{\hbar\omega}{\omega} \nonumber \\
    Z_{\tau(\circlearrowleft)} &= \hbar \tag{\ref{eq:planck_impedance}}
\end{align}
Planck's reduced constant ($\hbar$) does not evaluate as an abstract quantum of action; it evaluates explicitly as the absolute Rotational Kinematic Impedance of the $\boldsymbol{\tau}$-fluid. The equation $E = \hbar\omega$ is   the classical definition of mechanical impedance ($Z = E/\omega$).

Furthermore, as the internal fluid circulation of a topological boundary approaches the absolute acoustic limit ($|\nabla \times \mathbf{v}_{\psi}| \to c$), the rotational admittance mathematically vanishes ($\alpha_{\circlearrowleft} \to 0$). Because modal impedance evaluates   as the inverse of fluid compliance ($Z_i \equiv 1/\alpha_i$), the localized dynamic rotational impedance diverges to absolute infinity:
\begin{equation}
    \lim_{|\nabla \times \mathbf{v}_{\psi}| \to c} Z_{\tau(\circlearrowleft)} \to \infty \tag{\ref{eq:infinite_rotational_impedance}}
\end{equation}
In legacy quantum mechanics, this infinite rotational impedance evaluates as the rigid quantization of intrinsic Spin (e.g., $\hbar/2$). Because the vacuum evaluates with infinite structural resistance to additional localized vorticity, fundamental particles evaluate mechanically as absolutely rigid gyroscopes. They cannot be spun up or slowed down through classical interaction, evaluating phenomenologically as Chiral Symmetry Breaking and Gyroscopic Spin Stiffness. 

In natural units, evaluating $\hbar = 1$ does not diminish this mechanical limit to a dimensionless conversion artifact; rather, it mathematically normalizes the baseline rotational impedance of the macroscopic vacuum to unity, against which all localized topological spin is quantized.

\begin{tcolorbox}[recallbox, title={\ref{box:universal_impedance}. The Kinematic Impedance Tensor ($[\boldsymbol{Z}_{\tau}]$)}]
\begin{align}
    Z_{\tau(\perp)} &= \frac{h}{q^2} \equiv R_K \tag{\ref{eq:von_klitzing_impedance}} \\
    Z_{\tau(\parallel)} &= \hbar^2 \left( \frac{d^2 E}{dk^2} \right)^{-1} \equiv m^* \tag{\ref{eq:effective_mass_impedance}} \\
    Z_{\tau(\perp)} &\propto n(r) = \left( 1 - \frac{v_0^2}{c^2} \right)^{-1} \tag{\ref{eq:refractive_impedance}} \\
    Z_{\tau(\circlearrowleft)} &= \frac{E}{\omega} \equiv \hbar \tag{\ref{eq:planck_impedance}} \\
    \lim_{|\nabla \times \mathbf{v}_{\psi}| \to c} Z_{\tau(\circlearrowleft)} &\to \infty \implies \text{Spin Quantization } (\hbar/2) \tag{\ref{eq:infinite_rotational_impedance}}
\end{align}
\end{tcolorbox}



% ==========================================
% PHASE III: TOPOLOGICAL MASS GENERATION
% ==========================================

\part*{PHASE III: TOPOLOGICAL MASS GENERATION}






% ----------------------------------------------------------------------
\section{The Kinematic Origin of Inertia (Hydrodynamic Rest Mass)}
\label{sec:deriv_hydrodynamic_mass}

The kinematic resistance of a localized boundary to spatial translation (inertia) evaluates as the macroscopic scalar fluid drag of the continuum, decoupling inertial mass from localized topology.

\textbf{1. The Massless Void ($m_{intrinsic}$):} \\
A fundamental topological boundary evaluates as a localized geometric cavitation (void) generated by circulating metric fluid. Because the localized coordinate volume ($V_\tau$) is devoid of continuous metric fluid ($\rho_{void} = 0$), the intrinsic mass of the topological boundary evaluates to absolute zero:
\begin{equation}
    m_{intrinsic} = 0 \tag{\ref{eq:intrinsic_mass}}
\end{equation}
A geometric void possesses no internal mechanical mechanism to resist spatial acceleration.

\textbf{2. The Effective Kinematic Mass ($m_\tau$):} \\
In classical continuum mechanics, when a massless spherical cavity accelerates through an incompressible fluid of uniform baseline density ($\rho_{\tau}$), the applied force evaluates against the inertia of the surrounding fluid displaced by the cavity's trajectory. 

The effective observable rest mass of the boundary ($m_\tau$) evaluates as the sum of its intrinsic mass and the hydrodynamic added mass ($m_{added}$) derived via the Massieu Transform:
\begin{equation}
    m_\tau = m_{intrinsic} + m_{added} 
\end{equation}
Substituting the massless void limit zeroes the intrinsic term. Inertial rest mass evaluates exclusively as the entrained mass of the macroscopic continuum:
\begin{equation}
    m_\tau = \rho_{\tau} V_{added} \tag{\ref{eq:kinematic_mass}}
\end{equation}

\textbf{3. Mass-Volume Equivalence \& The Bernoulli Deficit:} \\
The geometric volume of the hydrodynamic added mass evaluates to half the rest volume of the primary spherical boundary ($V_{added} = \frac{1}{2} V_\tau$). This Archimedean geometric fraction ($1/2$) is physically identical to the Bernoulli dynamic pressure coefficient ($\frac{1}{2}\rho v^2$), proving that rest mass is fundamentally a measurement of the structural pressure deficit. 

Substituting this volumetric strain limit isolates the absolute relation between observable topological mass and localized metric volume:
\begin{align}
    m_\tau &= \rho_{\tau} \left( \frac{1}{2} V_\tau \right)  \\
    m_\tau &= \frac{1}{2} \rho_{\tau} V_\tau \tag{\ref{eq:mass_volume_equivalence}}
\end{align}

\textbf{4. The Illusion of Inertia:} \\
Observable localized rest mass evaluates as a geometric macroscopic property. The phenomenon of inertia is simply the localized evaluation of spatial translation against the structural yielding limits of the continuous vacuum metric, driven entirely by buoyancy within a Bernoulli pressure gradient.

\begin{tcolorbox}[recallbox, title={\ref{box:hydrodynamic_mass}. The Kinematic Origin of Inertia}]
\begin{align}
    m_{intrinsic} &= \rho_{void} V_\tau = 0 \tag{\ref{eq:intrinsic_mass}} \\
    m_\tau &= m_{added} = \rho_{\tau} V_{added} \tag{\ref{eq:kinematic_mass}} \\
    m_\tau &= \frac{1}{2} \rho_{\tau} V_\tau \tag{\ref{eq:mass_volume_equivalence}}
\end{align}
\end{tcolorbox}




% ----------------------------------------------------------------------
\section{Navier-Stokes Regularization (Topological Pinch-Off)}
\label{sec:deriv_ns_reg}

The finite-time singularity of the 3D Navier-Stokes equations evaluates as mechanically regularized by the macroscopic acoustic limit of the metric. The singularity is prevented by an absolute viscosity threshold.

\textbf{1. The Kinematic Boundary Condition:} \\
In a turbulent fluid cascade, geometric enstrophy divergence dictates that the rotational velocity ($v$) of a shrinking fluid eddy ($r \to 0$) increases continuously. The metric imposes a macroscopic kinematic boundary condition: internal fluid circulation cannot exceed the acoustic phase limit ($v \le c$).

When the localized fluid velocity strikes the geometric limit ($v = c$), the radius of the vortex evaluates as constrained to its Compton focal length ($L = \hbar / m c$). 

\textbf{2. Absolute Dynamic Viscosity ($\mu$):} \\
In continuous fluid mechanics, dynamic viscosity ($\mu$) evaluates as the ratio of absolute Shear Stress ($\tau_{shear}$) to the local Velocity Gradient ($\nabla v$).

The maximum shear stress the continuum can support before topological failure evaluates as the absolute Bernoulli dynamic pressure limit of the macroscopic vacuum ($\tau_{shear} = \frac{1}{2} \rho_{\tau} c^2$). 
The velocity gradient of the topological vortex evaluates as the kinematic limit ($c$) divided by the constrained boundary radius ($\nabla v = c / L$). Substituting the Compton length ($L = \hbar/mc$) yields:
\begin{equation}
    \nabla v = \frac{c}{\left(\frac{\hbar}{mc}\right)} = \frac{mc^2}{\hbar} 
\end{equation}

Dividing the absolute shear stress by the boundary velocity gradient isolates the absolute dynamic viscosity of the localized metric:
\begin{align}
    \mu &= \frac{\tau_{shear}}{\nabla v}  \\
    \mu &= \frac{\frac{1}{2} \rho_{\tau} c^2}{\left( \frac{mc^2}{\hbar} \right)}  \\
    \mu &= \frac{1}{2} \frac{\rho_{\tau} \hbar}{m} \tag{\ref{eq:vacuum_viscosity}}
\end{align}

\textbf{3. Prandtl Mixing Length Equivalence:} \\
Because the Compton focal length establishes the geometric boundary of the acoustic cavity ($L = \hbar/mc$), rearranging the spatial relation yields $\hbar/m = cL$. Substituting this continuous spatial identity directly back into the absolute viscosity uncovers the classical mechanical origin of mass:
\begin{equation}
    \mu = \frac{1}{2} \rho_{\tau} c L \tag{\ref{eq:prandtl_equivalence}}
\end{equation}

The macroscopic vacuum evaluates exactly according to Prandtl’s Mixing Length Theory for turbulent eddy viscosity ($\mu_t \sim \rho v l$). The continuous vacuum evaluates as a shear-thickening fluid. As the internal circulation accelerates toward the kinematic limit ($c$) at the localized boundary length ($L$), the localized viscosity of the continuum spikes to an absolute geometric maximum.

\textbf{4. The Critical Reynolds Number ($Re_c$):} \\
The macroscopic Critical Reynolds Number ($Re_c$) evaluating topological vortex pinch-off is defined by the standard fluid ratio of inertial forces to viscous forces:
\begin{equation}
    Re_c = \frac{\rho_{\tau} v L}{\mu} 
\end{equation}

Substituting the kinematic limit ($v=c$) and the derived Prandtl dynamic viscosity ($\mu = \frac{1}{2} \rho_{\tau} c L$) geometrically evaluates the topological transition state:
\begin{align}
    Re_c &= \frac{\rho_{\tau} c L}{\frac{1}{2} \rho_{\tau} c L}  \\
    Re_c &= 2 \tag{\ref{eq:critical_reynolds}}
\end{align}

The continuum mechanics of the $\boldsymbol{\tau}$-fluid dictate that as a vortex approaches the $v=c$ limit, the localized Reynolds number evaluates to exactly $2$. In classical fluid mechanics, $Re < 10$ defines absolute Stokes Flow (creeping flow), where viscous dampening fundamentally prevents runaway turbulent cascades. 

The finite-time singularity is mechanically impossible. The intense localized shear-thickening viscosity of the continuum forces the runaway cascade to smoothly pinch off into a stable, quantized topological particle.

\begin{tcolorbox}[recallbox, title={\ref{box:ns_reg}. Navier-Stokes Regularization}]
\begin{align}
    \mu &= \frac{1}{2} \frac{\rho_{\tau} \hbar}{m} \tag{\ref{eq:vacuum_viscosity}} \\
    \mu &= \frac{1}{2} \rho_{\tau} c L \tag{\ref{eq:prandtl_equivalence}} \\
    Re_c &= \frac{\rho_{\tau} c L}{\mu} = 2 \tag{\ref{eq:critical_reynolds}}
\end{align}
\end{tcolorbox}








% ----------------------------------------------------------------------
\section{Origin of Inertia (The Pressure Gradient \& Strain Compactness)}
\label{sec:deriv_fluid_inertia}

Apparent mass and kinematic inertia evaluate as the localized scalar pressure deficit generated by circulating $\boldsymbol{\tau}$-fluid traversing a continuous rheological gradient.

\textbf{1. Relativistic Strain Compactness (The Density-Pressure Identity):} \\
Because the continuous metric evaluates with a 3D volumetric mass density ($\rho_{\tau}$), accelerating the localized circulation velocity ($v_\phi$) subjects the effective fluid boundary to kinematic packing.

Substituting the Bernoulli compliance limit ($\alpha_s = \sqrt{P_{static}/P_c}$) into the Lorentz packing factor ($\rho_{eff} = \rho_{\tau} / \alpha_s$) isolates the fluid-mechanical equation of state governing metric density:
\begin{equation}
    \rho_{eff}(r) = \rho_{\tau} \left( \frac{P_{static}(r)}{P_c} \right)^{-1/2} = \rho_{\tau} \left( 1 - \frac{P_{dyn}(r)}{P_c} \right)^{-1/2} \tag{\ref{eq:fluid_inertia_gradient}}
\end{equation}
Effective mass density evaluates as the inverse of available static pressure. As dynamic kinematic circulation ($P_{dyn}$) consumes the structural capacity of the vacuum ($P_c$), the localized fluid density compacts ($\rho_{eff}$).

\textbf{2. The Dimensional Tension Bridge ($T_{1D}$ to $P_c$):} \\
The absolute 3D hydrostatic pressure capacity of the vacuum evaluates fundamentally from its 1D topological limits. The absolute 1D linear tension ($T_{1D}$) of a transverse acoustic metric string evaluates as its 1D metric stiffness ($\epsilon_{\tau}$) scaled by the squared kinematic limit:
\begin{equation}
    T_{1D} = \epsilon_{\tau} c^2 \label{eq:1d_tension}
\end{equation}
Because 3D pressure evaluates dimensionally as 1D tension distributed across a 2D cross-sectional area, distributing this absolute linear tension across the minimum structural quantum boundary (the Gravitational Yield Area, $A_{yield}$) evaluates the absolute 3D volumetric energy density ($\rho_{\tau} c^2$). Applying the continuous hydrodynamic fraction isolates the 3D acoustic yield pressure:
\begin{align}
    P_c &= \frac{1}{2} \frac{T_{1D}}{A_{yield}} \nonumber \\
    P_c &= \frac{1}{2} \left( \frac{\epsilon_{\tau}}{A_{yield}} \right) c^2 \nonumber \\
    P_c &= \frac{1}{2} \rho_{\tau} c^2 \tag{\ref{eq:dimensional_tension_bridge}}
\end{align}

\textbf{3. The Isotropic Tension Identity ($\mathbb{T}_{local}$):} \\
Because the macroscopic metric evaluates as a 3D elastic continuum under geometric tautness, its baseline static pressure equates to an isotropic volumetric tension ($\mathbb{T}_c = -P_c$). Dimensionally, this tension evaluates as a continuous Cauchy stress ($N/m^2$), mechanically equivalent to localized energy density ($J/m^3$). 

Applying the Bernoulli Conservation Law ($P_{static} = P_c - P_{dyn}$) translates the localized pressure deficit into localized geometric tension:
\begin{align}
    \mathbb{T}_{local} &= - P_{static} \nonumber \\
    \mathbb{T}_{local} &= - (P_c - P_{dyn}) \nonumber \\
    \mathbb{T}_{local} &= -P_c + P_{dyn} \nonumber \\
    \mathbb{T}_{local} &= \mathbb{T}_c + P_{dyn} \tag{\ref{eq:tension_identity}}
\end{align}
The internal kinetic circulation of the topological knot creates dynamic pressure ($P_{dyn}$). To maintain equilibrium, the localized spatial metric tightens. The total static tension of the particle ($\mathbb{T}_{local}$) evaluates as the baseline tension of the vacuum ($\mathbb{T}_c$) plus the dynamic load of the circulating fluid. The coordinate space pulls inward, causing the surrounding fluid to compact into higher mass density.

\textbf{4. The Geometric Equation of State (Pure Tension):} \\
To isolate the geometric structural density of the metric, the dynamic kinetic pressure ($P_{dyn} = \mathbb{T}_{local} - \mathbb{T}_c$) is algebraically substituted out of the Density-Pressure Identity:
\begin{align}
    \rho_{eff}(r) &= \rho_{\tau} \left( 1 - \frac{\mathbb{T}_{local}(r) - \mathbb{T}_c}{\mathbb{T}_c} \right)^{-1/2} \nonumber \\
    \rho_{eff}(r) &= \rho_{\tau} \left( 1 - \left( \frac{\mathbb{T}_{local}(r)}{\mathbb{T}_c} - 1 \right) \right)^{-1/2} \nonumber \\
    \rho_{eff}(r) &= \rho_{\tau} \left( 2 - \frac{\mathbb{T}_{local}(r)}{\mathbb{T}_c} \right)^{-1/2} \tag{\ref{eq:tension_density}}
\end{align}
The effective mass density of the continuum evaluates exclusively as the ratio of localized tension to baseline tension. The continuum sustains exactly double its baseline tension ($2\mathbb{T}_c$) before the structural geometry yields.

\textbf{5. The Compliant Boundary ($\lambda_c$):} \\
The metric evaluates the topological defect at its observable macroscopic boundary (the Compton wavelength, $r = \lambda_c$). At this radius, localized circulation evaluates as sub-luminal ($v_\phi \ll c$), and the dynamic pressure consumes a minority fraction of the yield capacity ($P_{dyn} \ll P_c$).

Expanding the root via the binomial approximation ($(1-x)^{-1/2} \approx 1 + x/2$) isolates the origin of observable mass:
\begin{equation}
    \rho_{eff}(\lambda_c) \approx \rho_{\tau} + \rho_{\tau} \left( \frac{P_{dyn}(\lambda_c)}{2P_c} \right) \tag{\ref{eq:compliant_inertia}}
\end{equation}
The observable mass of a particle evaluates as the baseline background fluid ($\rho_{\tau}$) plus a thermodynamic mass excess generated by localized dynamic tension dragging the coordinate space inward.

\textbf{6. The Dimensionless Singularity Contradiction:} \\
For an irrotational topological circulation ($v_\phi \propto 1/r$), evaluating the core as a zero-dimensional point ($r \to 0$) forces the localized velocity to diverge to infinity ($v_\phi \to \infty$). A dimensionless singularity mandates a dynamic tension exceeding the absolute structural capacity of the vacuum ($P_{dyn} \gg P_c$), driving the localized geometric tension past the absolute yield limit ($\mathbb{T}_{local} \gg 2\mathbb{T}_c$).

Substituting this limit into the Density-Pressure Identity yields an imaginary state:
\begin{equation}
    \rho_{eff}(0) = \rho_{\tau} (1 - \infty)^{-1/2} \implies \mathbb{I} \text{ (Imaginary Density)} \nonumber
\end{equation}
A macroscopic continuum cannot sustain an imaginary mass density. The fluid mathematics mandate geometric boundaries to constrain the scalar field to real values ($\mathbb{R}$).

\textbf{7. The Cavitation Yield Limit ($l_p$):} \\
Before dynamic tension crosses into the imaginary regime, internal fluid circulation strikes the acoustic limit ($v_\phi \to c$). The dynamic load equals the total metric capacity ($P_{dyn} \to P_c$), and the localized tension strikes the absolute yield limit ($\mathbb{T}_{local} \to 2\mathbb{T}_c$). The effective density diverges to a real infinity:
\begin{equation}
    \lim_{P_{dyn} \to P_c} \rho_{eff}(l_p) \to \infty \tag{\ref{eq:infinite_core_inertia}}
\end{equation}
The localized spatial coordinate winches inward under extreme geometric tension until it structurally yields. The fluid cavitates at a finite spatial radius ($l_p$). The continuum shatters into a topological void before crossing into an imaginary state, constructing an impenetrable boundary to stabilize the defect.

\begin{tcolorbox}[recallbox, title={\ref{box:fluid_inertia}. Origin of Inertia (Strain Compactness \& Tension)}]
\begin{align}
    \rho_{eff}(r) &= \rho_{\tau} \left( 1 - \frac{P_{dyn}(r)}{P_c} \right)^{-1/2} \tag{\ref{eq:fluid_inertia_gradient}} \\
    T_{1D} &= \epsilon_{\tau} c^2 \implies P_c = \frac{1}{2} \frac{T_{1D}}{A_{yield}} \tag{\ref{eq:dimensional_tension_bridge}} \\
    \mathbb{T}_{local} &= \mathbb{T}_c + P_{dyn} \tag{\ref{eq:tension_identity}} \\
    \rho_{eff}(r) &= \rho_{\tau} \left( 2 - \frac{\mathbb{T}_{local}(r)}{\mathbb{T}_c} \right)^{-1/2} \tag{\ref{eq:tension_density}} \\
    \rho_{eff}(\lambda_c) &\approx \rho_{\tau} + \rho_{\tau} \left( \frac{P_{dyn}(\lambda_c)}{2P_c} \right) \tag{\ref{eq:compliant_inertia}} \\
    \lim_{P_{dyn} \to P_c} \rho_{eff}(l_p) &\to \infty \label{eq:infinite_core_inertia}
\end{align}
\end{tcolorbox}


% ----------------------------------------------------------------------
\section{The Volumetric Inverse-Cube Identity}
\label{sec:deriv_volumetric_inverse_cube}

Because matter is defined as a cavitated topological boundary, its physical limits must mathematically reconcile the Euclidean spatial footprint of the boundary with the Riemannian macroscopic fluid strain displaced by its void. 

\textbf{1. The Bounding Volume ($V_R$):} \\
The geometric bounding volume ($V_R$) of a topological vortex evaluates as a 3D Euclidean sphere bounded by the continuous kinematic limit of its circulation (the Compton radius, $R = \hbar/mc$):
\begin{align}
    V_R &= \frac{4}{3}\pi R^3 \nonumber \\
    V_R &= \frac{4}{3}\pi \left( \frac{\hbar}{m_\tau c} \right)^3 \tag{\ref{eq:bounding_volume}}
\end{align}

\textbf{2. The Inverse-Cube Substitution:} \\
To evaluate the direct geometric relationship between the external bounding box ($V_R$) and the internal topological void ($V_{\tau}$), we substitute the derived Kinematic Origin of Inertia ($m_\tau = \frac{1}{2} \rho_{\tau} V_{\tau}$) directly into the bounding volume equation:
\begin{equation}
    V_R = \frac{4}{3}\pi \left( \frac{\hbar}{(\frac{1}{2}\rho_{\tau} V_{\tau}) c} \right)^3 
\end{equation}

Distributing the cubic exponent across the denominator isolates the displaced void volume ($V_{\tau}$) and extracts the hydrodynamic fraction:
\begin{equation}
    V_R = \frac{4}{3}\pi \left( \frac{8\hbar^3}{\rho_{\tau}^3 V_{\tau}^3 c^3} \right) 
\end{equation}

Multiplying both sides of the equation by the cubed fluid displacement ($V_{\tau}^3$) algebraically extracts the topological variables from the macroscopic continuum constants:
\begin{align}
    V_R \cdot V_{\tau}^3 &= \frac{4}{3}\pi \left( \frac{8\hbar^3}{\rho_{\tau}^3 c^3} \right)  \\
    V_R \cdot V_{\tau}^3 &= \frac{32\pi}{3} \left( \frac{\hbar}{\rho_{\tau} c} \right)^3 \tag{\ref{eq:volumetric_inverse_cube}}
\end{align}

\textbf{3. Volumetric Equilibrium:} \\
Because Planck's reduced constant ($\hbar$), the baseline fluid density ($\rho_{\tau}$), and the kinematic acoustic limit ($c$) evaluate as absolute macroscopic invariants of the metric, the algebraic right side of the equation evaluates as a universal geometric constant. 

Localized inertia ($m_\tau$) algebraically annihilates from the fundamental boundary condition. The geometric bounding volume of fundamental matter evaluates proportional to the inverse-cube of its displaced fluid volume. As the spatial radius of a topological boundary compresses, the displaced macroscopic void must geometrically dilate by a rigid cubic ratio to conserve the absolute structural limits of the continuum.

\begin{tcolorbox}[recallbox, title={\ref{box:volumetric_inverse_cube}. The Volumetric Inverse-Cube Identity}]
\begin{equation}
    V_R \cdot V_{\tau}^3 = \frac{32\pi}{3} \left( \frac{\hbar}{\rho_{\tau} c} \right)^3 \tag{\ref{eq:volumetric_inverse_cube}}
\end{equation}
\end{tcolorbox}



% ----------------------------------------------------------------------
\section{The Admittance-Modified Navier-Stokes Equation}
\label{sec:deriv_admittance_ns}

To evaluate the acceleration of a fundamental topological boundary, the continuous momentum of the metric evaluates against the localized fluid admittance ($\boldsymbol{\alpha}$). 

\textbf{1. The Relativistic Density Tensor ($\boldsymbol{\rho}_{eff}$):} \\
By classical continuum mechanics, the material derivative representing fluid acceleration ($\frac{D\mathbf{v}}{Dt}$) evaluates against the effective localized fluid density ($\rho_{eff}$). Kinematic packing dictates that effective density scales against macroscopic Lorentz strain.

Because localized kinematic phase space is constrained by the absolute acoustic bottleneck ($c^2$), fluid compliance operates as an anisotropic viscoelastic tensor ($[\boldsymbol{\alpha}]$). The total scalar admittance ($\alpha_s$) evaluates as the absolute volumetric compliance of the metric, derived   from the Pythagorean capacity of the total velocity field ($v_{tot}^2 = v_{\parallel}^2 + v_{\perp}^2 + v_{\circlearrowleft}^2$):
\begin{equation}
    \alpha_s = \sqrt{1 - \frac{v_{tot}^2}{c^2}} \nonumber
\end{equation}
The localized effective density evaluates as the baseline 3D macroscopic mass density ($\rho_{\tau}$) dilated by the inverse of the full Admittance Tensor ($[\boldsymbol{\alpha}]^{-1}$):
\begin{equation}
    \boldsymbol{\rho}_{eff} = \rho_{\tau} [\boldsymbol{\alpha}]^{-1} \tag{\ref{eq:fluid_inertia_gradient}}
\end{equation}

\textbf{2. The Quantum Reynolds Limit (Hadronization):} \\
In fluid dynamics, the Reynolds number ($Re$) dictates the transition from laminar flow to chaotic turbulence via the ratio of inertial mass density to internal viscous shear ($Re = \rho v L / \mu$). 

To evaluate the localized Quantum Reynolds Number ($Re_q$) of a topological boundary, substitute the scalar effective density ($\rho_{eff} = \rho_{\tau}/\alpha_s$), the kinematic circulation limit of the core ($v = c$), and the absolute dynamic viscosity of the continuum derived via the Prandtl boundary limit ($\mu = \frac{1}{2} \rho_{\tau} c L$):
\begin{align}
    Re_q &= \frac{\rho_{eff} c L}{\mu} \nonumber \\
    Re_q &= \frac{\left( \frac{\rho_{\tau}}{\alpha_s} \right) c L}{\frac{1}{2} \rho_{\tau} c L} \nonumber
\end{align}

The macroscopic continuum constants ($\rho_{\tau}, c, L$) annihilate from the numerator and denominator, isolating the transition limit of the metric:
\begin{equation}
    Re_q = \frac{2}{\alpha_s} \tag{\ref{eq:quantum_reynolds}}
\end{equation}

At the topological ground state ($\alpha_s = 1$), the Reynolds number evaluates to exactly $Re_q = 2$, defining absolute Stokes flow. When kinematic energy drives the effective density upward, the total scalar admittance collapses ($\alpha_s \to 0$), and the localized Reynolds number diverges ($Re_q \to \infty$). The electromagnetic continuum shatters into chaotic turbulence. To relieve localized strain, the fluid divides its volume, nucleating additional boundaries (hadronization).

\textbf{3. The Equivalence Principle (Dimensional Immunity):} \\
The Navier-Stokes equation evaluates equilibrium as Force Densities ($\text{N/m}^3$). There is a dimensional distinction between Body Forces and Surface Forces.
\begin{itemize}[nosep]
    \item \textbf{Body Forces ($\mathbf{F}_{body}$):} Gravity ($-\nabla \Phi$) and Inertia act upon the mass of the fluid. Because their kinematic units evaluate as accelerations ($\text{m/s}^2$), they must be multiplied by the effective mass density ($\boldsymbol{\rho}_{eff}$) to evaluate dimensionally as Force Densities ($\text{kg/m}^3 \times \text{m/s}^2 = \text{N/m}^3$).
    \item \textbf{Surface Forces ($\mathbf{F}_{surface}$):} Mechanical compression, planar shear, and internal viscosity act upon the geometric boundary. Their spatial gradients evaluate as Force Densities ($\text{N/m}^3$) independent of internal mass density.
\end{itemize}

Evaluating the Navier-Stokes momentum equation using these physical dimensions yields:
\begin{equation}
    \boldsymbol{\rho}_{eff} \frac{D\mathbf{v}}{Dt} = -\boldsymbol{\rho}_{eff} \nabla \Phi + \Sigma \mathbf{F}_{surface} \nonumber
\end{equation}
Substituting the tensor compactness identity ($\boldsymbol{\rho}_{eff} = \rho_{\tau} [\boldsymbol{\alpha}]^{-1}$):
\begin{equation}
    \left( \rho_{\tau} [\boldsymbol{\alpha}]^{-1} \right) \frac{D\mathbf{v}}{Dt} = -\left( \rho_{\tau} [\boldsymbol{\alpha}]^{-1} \right) \nabla \Phi + \Sigma \mathbf{F}_{surface} \nonumber
\end{equation}

Multiplying the equation by the dimensionless Admittance Tensor ($[\boldsymbol{\alpha}]$) evaluates the tensor against its inverse ($[\boldsymbol{\alpha}][\boldsymbol{\alpha}]^{-1} = \mathbb{I}$). The tensor annihilates from both the acceleration and gravitational terms, attaching itself exclusively to the surface forces:
\begin{equation}
    \rho_{\tau} \frac{D\mathbf{v}}{Dt} = -\rho_{\tau} \nabla \Phi + [\boldsymbol{\alpha}] \Sigma \mathbf{F}_{surface} \tag{\ref{eq:admittance_ns}}
\end{equation}
Gravity evaluates as immune to localized kinematic strain. Because gravity evaluates as the ambient macroscopic scalar pressure gradient of the continuous metric itself, inertial density cancels with gravitational density. This derives the Weak Equivalence Principle from continuum mechanics; acceleration due to gravity evaluates independent of internal mass or boundary velocity, while external physical forces remain throttled by $[\boldsymbol{\alpha}]$.

\textbf{4. Dimensional Phase-Space Mapping (Cauchy-Helmholtz Decomposition):} \\
Because the full Cauchy stress tensor ($\boldsymbol{\sigma}$) mathematically separates normal volumetric pressure onto its main diagonal and deviatoric shear onto its off-diagonals, the Admittance Tensor ($[\boldsymbol{\alpha}]$) operates as a selective geometric filter. It performs a physical Helmholtz decomposition, projecting the orthogonal kinematic compliance limits into uniform macroscopic Force Densities ($\text{N/m}^3$):

\begin{itemize}[leftmargin=*]
    \item \textbf{Normal Stress [$\text{m/s}$]:} Longitudinal translation evaluates as dynamic pressure ($P = \frac{1}{2}\rho_{\tau} v_{\parallel}^2$). The longitudinal admittance ($\alpha_{\parallel}$) scales strictly the main diagonal of the stress tensor. The spatial gradient ($-\nabla$) converts this normal stress into irrotational compressive force density ($\nabla \times \nabla P = 0$).
    \item \textbf{Transverse Shear [$\text{rad/m}$]:} Transverse spatial strain evaluates as lateral angular deformation across distance. The transverse admittance ($\alpha_{\perp}$) scales the symmetric off-diagonal deviatoric shear stresses ($\tau_{ij}$). The divergence ($\nabla \cdot$) converts this off-diagonal spatial geometry into lateral force density.
    \item \textbf{Rotational Vorticity [$\text{rad/s}$]:} Localized proper time evaluates as internal angular vorticity ($\boldsymbol{\omega} = \nabla \times \mathbf{v}$). The rotational admittance ($\alpha_{\circlearrowleft}$) scales the antisymmetric off-diagonal shear components. At the incompressible core ($\nabla \cdot \mathbf{v} = 0$), the viscous vector identity reduces to solenoidal rotational force density ($-\mu \nabla \times \boldsymbol{\omega}$).
\end{itemize}

Expanding the tensor multiplication maps the localized surface forces perfectly to their respective modal admittances:
\begin{equation}
    \rho_{\tau} \frac{D\mathbf{v}}{Dt} = -\rho_{\tau} \nabla \Phi + \alpha_{\parallel}(-\nabla P) + \alpha_{\perp}(\nabla \cdot \mathbb{T}_{shear}) + \alpha_{\circlearrowleft}(\mu \nabla^2 \mathbf{v}) \tag{\ref{eq:propulsion_master}}
\end{equation}

\textbf{5. The Absolute Kinematic Equation of Motion:} \\
To isolate the pure geometric acceleration of the fluid boundary, the absolute background density ($\rho_{\tau}$) factors out of the classical surface operators via explicit kinematic substitutions. Furthermore, by the derived Pressure-Curvature Identity ($\Phi = -P_{dyn}/\rho_{\tau}$), the gravitational body force unifies as an ambient macroscopic kinematic pressure gradient. Because the baseline metric operates under isotropic tension ($\mathcal{T} = -P$), the gradient of positive dynamic pressure evaluates mechanically as the absolute kinematic tension gradient ($\nabla \mathcal{P}_{ambient} \equiv \nabla \mathcal{T}_{ambient} = -\nabla \Phi$).

\begin{itemize}[leftmargin=*]
    \item \textbf{Kinematic Pressure ($\mathcal{P}$):} Dynamic pressure evaluates as $P = \rho_{\tau} \mathcal{P}$. The gradient isolates the density: $-\nabla P = \rho_{\tau}(-\nabla \mathcal{P}_{dyn})$.
    \item \textbf{Kinematic Shear ($\mathbb{T}_{kin}$):} Because the absolute metric shear modulus evaluates as $\mu_s = \rho_{\tau} c^2$, the deviatoric stress isolates the density ($\mathbb{T}_{shear} = \rho_{\tau} 2c^2 \varepsilon' \implies \rho_{\tau} \mathbb{T}_{kin}$). The divergence evaluates as: $\nabla \cdot \mathbb{T}_{shear} = \rho_{\tau}(\nabla \cdot \mathbb{T}_{kin})$.
    \item \textbf{Kinematic Viscosity ($\nu$):} Absolute dynamic viscosity evaluates via the Prandtl threshold ($\mu = \rho_{\tau} \frac{1}{2}cL$). Substituting this limit defines the mass-independent kinematic viscosity ($\nu = \frac{1}{2}cL$). The viscous gradient evaluates as: $\mu \nabla^2 \mathbf{v} = \rho_{\tau}(\nu \nabla^2 \mathbf{v})$.
\end{itemize}

Dividing the entire Navier-Stokes momentum equation by the absolute macroscopic density ($\rho_{\tau}$) annihilates mass from the local frame evaluation:
\begin{equation}
    \frac{D\mathbf{v}}{Dt} = \nabla \mathcal{P}_{ambient} - \alpha_{\parallel}\nabla \mathcal{P}_{dyn} + \alpha_{\perp}(\nabla \cdot \mathbb{T}_{kin}) + \alpha_{\circlearrowleft}(\nu \nabla^2 \mathbf{v}) \nonumber
\end{equation}

Applying the continuous Bernoulli Conservation Law exposes the fundamental symmetry of acceleration. Because the absolute kinematic yield limit of the vacuum ($\mathcal{P}_c = c^2/2$) is constant, its spatial gradient evaluates to zero. Taking the spatial derivative of the kinematic phase exchange ($\mathcal{P}_{static} + \mathcal{P}_{dyn} = \mathcal{P}_c$) isolates the Equivalence Principle:
\begin{equation}
    \nabla \mathcal{P}_{static} + \nabla \mathcal{P}_{dyn} = 0 \implies -\nabla \mathcal{P}_{dyn} = \nabla \mathcal{P}_{static}
\end{equation}

The localized dynamic pressure gradient generated by mechanical propulsion ($-\nabla \mathcal{P}_{dyn}$) evaluates mathematically identically to climbing a macroscopic static pressure gradient ($\nabla \mathcal{P}_{static}$). Unifying the pressure gradients into the shared static scalar field defines the Absolute Kinematic Equation of Motion:
\begin{equation}
    \frac{D\mathbf{v}}{Dt} = \nabla \mathcal{P}_{ambient} + \alpha_{\parallel}\nabla \mathcal{P}_{static} + \alpha_{\perp}(\nabla \cdot \mathbb{T}_{kin}) + \alpha_{\circlearrowleft}(\nu \nabla^2 \mathbf{v}) \tag{\ref{eq:kinematic_ns}}
\end{equation}

Gravity and mechanical propulsion evaluate as unified thermodynamic properties. The continuum distinguishes them   by geometric topological interaction: the ambient environmental tension gradient ($\nabla \mathcal{T}_{ambient}$) evaluates independently of boundary strain, drawing the coordinate space inward, while the localized surface pressure gradient ($\nabla \mathcal{P}_{static}$) remains structurally throttled by longitudinal admittance ($\alpha_{\parallel}$). Localized acceleration evaluates exclusively as scale-invariant geometric phase-space evolution. 

Furthermore, because this phase-space evolution is governed by the Eulerian Material Derivative ($\frac{D\mathbf{v}}{Dt} = \frac{\partial \mathbf{v}}{\partial t} + (\mathbf{v} \cdot \nabla)\mathbf{v}$), the continuous kinematic advection of the $\boldsymbol{\tau}$-fluid evaluates inherently on the left side of the equality, dictating exactly how the fluid physically transports its own momentum through these static gradients.

\begin{tcolorbox}[recallbox, title={\ref{box:admittance_ns}. Admittance-Modified Navier-Stokes}]
\begin{align}
    \alpha_s &= \sqrt{1 - \frac{v_{tot}^2}{c^2}} \nonumber \\
    \boldsymbol{\rho}_{eff} &= \rho_{\tau} [\boldsymbol{\alpha}]^{-1} \tag{\ref{eq:fluid_inertia_gradient}} \\
    Re_q &= \frac{2}{\alpha_s} \tag{\ref{eq:quantum_reynolds}} \\
    \rho_{\tau} \frac{D\mathbf{v}}{Dt} &= -\rho_{\tau} \nabla \Phi + [\boldsymbol{\alpha}] \Sigma \mathbf{F}_{surface} \tag{\ref{eq:admittance_ns}} \\
    \rho_{\tau} \frac{D\mathbf{v}}{Dt} &= -\rho_{\tau} \nabla \Phi + \alpha_{\parallel}(-\nabla P) + \alpha_{\perp}(\nabla \cdot \mathbb{T}_{shear}) + \alpha_{\circlearrowleft}(\mu \nabla^2 \mathbf{v}) \tag{\ref{eq:propulsion_master}} \\
    -\nabla \mathcal{P}_{dyn} &= \nabla \mathcal{P}_{static} \implies \text{Equivalence Principle} \nonumber \\
    \frac{D\mathbf{v}}{Dt} &= \nabla \mathcal{P}_{ambient} + \alpha_{\parallel}\nabla \mathcal{P}_{static} + \alpha_{\perp}(\nabla \cdot \mathbb{T}_{kin}) + \alpha_{\circlearrowleft}(\nu \nabla^2 \mathbf{v}) \tag{\ref{eq:kinematic_ns}}
\end{align}
\end{tcolorbox}


% ----------------------------------------------------------------------
\section{The Localized Metric (Conformal Duality \& Volumetric Strain)}
\label{sec:deriv_localized_metric_strain}

The fundamental metric evaluates via absolute Conformal Duality. The kinematic boundary conditions of a localized topological defect (the micro-scale) evaluate as the exact spatial inversion of the macroscopic expanding cosmos.

\textbf{1. Conformal Inversion (Topological Circulation):} \\
For the macroscopic continuum, the absolute velocity field approaches the acoustic phase limit at the causal horizon ($v \to c$ as $r \to \infty$). For a localized acoustic cavity, this boundary condition inverts. The absolute acoustic phase limit ($v \to c$) occurs exactly at the localized coordinate core ($r \to 0$), and decays to zero as it merges with the quiescent vacuum. 

Applying this geometric boundary swap ($1 - e^{-x} \implies e^{-x}$) natively defines the convergent microscopic phase velocity:
\begin{equation}
    v_{obs}(r) = c \cdot e^{-r/R_{h,p}} \tag{\ref{eq:vobs_conformal}}
\end{equation}

\textbf{2. The True Topological Radius \& Strain-Redshift Identity:} \\
By the Equivalence Principle, the fluid potential energy of this localized circulation ($\Phi_{local} = -v_{obs}^2/2$) dictates the spatial metric tensor ($g_{rr}$). The true physical Euclidean radius ($r_{true}$) evaluates as the differential integration of this strain ($dr_{true} = dr / \sqrt{1 - v_{obs}^2/c^2}$).

Because the localized gravitational redshift ($z$) evaluates via the time-time curvature tensor ($1+z = 1/\sqrt{-g_{tt}}$), substituting the localized circulation potential dictates that the spatial topological strain evaluates identically as the localized gravitational redshift equivalent:
\begin{equation}
    \varepsilon(r) = 1+z(r) = \frac{1}{\sqrt{1 - e^{-2r/R_{h,p}}}} \tag{\ref{eq:strain_redshift_identity}}
\end{equation}

\textbf{3. The Microscopic Conformal Mirror:} \\
Applying the continuous kinematic redshift identity $1+z_{local} = (1-v_{obs}/c)^{-1}$ to the inverted velocity field evaluates the internal geometric depth of the state:
\begin{align}
    1+z_{local} &= \frac{1}{1 - e^{-r/R_{h,p}}}  \\
    e^{-r/R_{h,p}} &= 1 - \frac{1}{1+z_{local}} = \frac{z_{local}}{1+z_{local}} 
\end{align}
Inverting both sides mathematically isolates the positive geometric exponent:
\begin{equation}
    e^{r/R_{h,p}} = 1 + \frac{1}{z_{local}} 
\end{equation}
Taking the natural logarithm natively derives the True Microscopic Distance ($r_{true, micro}$):
\begin{equation}
    r_{true, micro}(z) = R_{h,p} \ln\left( 1 + \frac{1}{z_{local}} \right) \tag{\ref{eq:conformal_r}}
\end{equation}
The true microscopic distance evaluates as the exact conformal mirror ($z \leftrightarrow z^{-1}$) of the macroscopic cosmos ($r_{true, cosmic} = R_{h,0} \ln(1+z)$). 

\textbf{4. The Integration of Inertia ($m_\tau$):} \\
By Eulerian volumetric conservation, the total continuous fluid displaced by the topological boundary ($V_\tau$) evaluates as the spherical integral of the excess spatial strain ($\varepsilon(r) - 1$) over the baseline unperturbed metric. Substituting the redshift identity ($\varepsilon(r) - 1 = z(r)$) isolates the volumetric displacement. 

By the derived Kinematic Origin of Inertia, observable rest mass evaluates as the hydrodynamic added mass ($m_\tau = \frac{1}{2} \rho_{\tau} V_\tau$). Substituting the spatial redshift integral yields the absolute continuum definition of localized inertia:
\begin{equation}
    m_\tau = \frac{1}{2} \rho_{\tau} \int_{l_p}^\infty 4\pi r^2 z(r) \, dr \tag{\ref{eq:mass_integral_redshift}}
\end{equation}
Fundamental mass evaluates mathematically as exactly half the spatial integration of the macroscopic topological redshift gradient.

\textbf{5. Topological Magnification (The Quantum Bit):} \\
Because the microscopic metric acts as a conformal mirror, the Apparent Rest Mass ($m_0$) structurally evaluates as the fundamental Planck Bit ($m_p$) divided by the geometric optical magnification of the compressed fluid surrounding the defect ($Mag$):
\begin{align}
    Mag &= \frac{R_{h,p}}{l_p}  \\
    Mag &= \frac{m_p}{m_0} \implies m_0 = \frac{m_p}{Mag} \tag{\ref{eq:conformal_mag}}
\end{align}

\textbf{6. The Planck Core Energy Limit ($r \to l_p$):} \\
To evaluate the physical reality of the topological core, we evaluate the exponential circulation limit as the radius approaches the Planck length. Because $l_p \ll R_{h,p}$, the exponential evaluates via the first-order Taylor expansion ($e^{-x} \approx 1 - x$):
\begin{equation}
    \varepsilon(l_p) = \frac{1}{\sqrt{1 - (1 - \frac{2l_p}{R_{h,p}})}} = \sqrt{\frac{R_{h,p}}{2l_p}} 
\end{equation}
Applying this geometric strain limit isolates the absolute thermodynamic energy density at the face of the topological core:
\begin{equation}
    \lim_{r \to l_p} E(r) \approx E_0 \sqrt{\frac{R_{h,p}}{2l_p}} \implies \text{Absolute Kinematic Lockout} \tag{\ref{eq:planck_core_energy_limit}}
\end{equation}
The conformal spatial inversion evaluates mechanically as a volumetric fluid wedge. It drives localized kinetic energy toward absolute infinity as the spatial coordinate narrows, forcing the circulating fluid to violently shatter against the Planck boundary, locking the localized energy into stable, quantized topological mass.

\begin{tcolorbox}[recallbox, title={\ref{box:localized_metric_strain}. The Localized Metric (Conformal Duality)}]
\begin{align}
    v_{obs}(r) &= c \cdot e^{-r/R_{h,p}} \tag{\ref{eq:vobs_conformal}} \\
    \varepsilon(r) &= 1+z(r) = \frac{1}{\sqrt{1 - e^{-2r/R_{h,p}}}} \tag{\ref{eq:strain_redshift_identity}} \\
    r_{true, micro}(z) &= R_{h,p} \ln\left( 1 + \frac{1}{z_{local}} \right) \tag{\ref{eq:conformal_r}} \\
    m_\tau &= \frac{1}{2} \rho_{\tau} \int_{l_p}^\infty 4\pi r^2 z(r) \, dr \tag{\ref{eq:mass_integral_redshift}} \\
    Mag &= \frac{R_{h,p}}{l_p} = \frac{m_p}{m_0} \implies m_0 = \frac{m_p}{Mag} \tag{\ref{eq:conformal_mag}} \\
    \lim_{r \to l_p} E(r) &\approx E_0 \sqrt{\frac{R_{h,p}}{2l_p}} \tag{\ref{eq:planck_core_energy_limit}}
\end{align}
\end{tcolorbox}


% ----------------------------------------------------------------------
\section{Confinement Pressure \& Work}
\label{sec:deriv_bag_pressure}

Because the topological boundary circulates at the acoustic limit ($v \to c$), it generates an extreme localized Bernoulli dynamic pressure deficit. To balance this thermodynamic deficit and geometrically secure the acoustic cavity, the continuum exerts an inward hydrostatic centripetal proper acceleration evaluating as $a_p = c/2\pi t_{loc}$. Because proper time evaluates geometrically as $t_{loc} = R_{h,p}/c$, the boundary acceleration isolates to:
\begin{align}
    a_p &= \frac{c}{2\pi (R_{h,p}/c)} \\
    a_p &= \frac{c^2}{2\pi R_{h,p}}
\end{align}

The geometric rest mass boundary evaluates as $R_{h,p} = \hbar/mc$. Substituting this dimension isolates the total hydrostatic restorative Force ($F = m a_p$) and surface Area ($A = 4\pi R_{h,p}^2$) constraints:
\begin{align}
    F &= m \left( \frac{c^2}{2\pi (\hbar/mc)} \right) \\
    F &= m \left( \frac{m c^3}{2\pi \hbar} \right) = \frac{m^2 c^3}{2\pi \hbar}
\end{align}
\begin{align}
    A &= 4\pi \left( \frac{\hbar}{mc} \right)^2 \\
    A &= \frac{4\pi \hbar^2}{m^2 c^2}
\end{align}

Dividing the restorative Force by the spherical surface Area isolates the absolute Confinement Pressure (Bag Pressure, $P_p$) counteracting the Bernoulli deficit:
\begin{align}
    P_p &= \frac{F}{A} \\
    P_p &= \frac{\left( \frac{m^2 c^3}{2\pi \hbar} \right)}{\left( \frac{4\pi \hbar^2}{m^2 c^2} \right)} \\
    P_p &= \left( \frac{m^2 c^3}{2\pi \hbar} \right) \left( \frac{m^2 c^2}{4\pi \hbar^2} \right) \\
    P_p &= \frac{m^4 c^5}{8\pi^2 \hbar^3}
\end{align}

Multiplying this Confinement Pressure by the physical Volume of the acoustic cavity ($V = \frac{4}{3}\pi R_{h,p}^3$) isolates the Volumetric Work ($W_v$) required to maintain topological confinement:
\begin{align}
    W_v &= P_p \cdot V \\
    W_v &= \left( \frac{m^4 c^5}{8\pi^2 \hbar^3} \right) \left( \frac{4}{3}\pi R_{h,p}^3 \right)
\end{align}
Substituting the boundary limit $R_{h,p}^3 = \hbar^3 / m^3 c^3$:
\begin{align}
    W_v &= \left( \frac{m^4 c^5}{8\pi^2 \hbar^3} \right) \left( \frac{4\pi \hbar^3}{3 m^3 c^3} \right) \\
    W_v &= \frac{4\pi m^4 c^5 \hbar^3}{24\pi^2 m^3 c^3 \hbar^3} \\
    W_v &= \frac{4 m c^2}{24 \pi} \\
    W_v &= \frac{m c^2}{6\pi}
\end{align}

\begin{tcolorbox}[recallbox, title={\ref{box:bag_pressure}. Confinement Pressure \& Work}]
\begin{equation}
    P_p = \frac{m^4 c^5}{8\pi^2 \hbar^3} \tag{\ref{eq:pp}}
\end{equation}
\end{tcolorbox}





% ----------------------------------------------------------------------
\section{The Unified Topological Engine}
\label{sec:deriv_unified_engine}

Fundamental matter evaluates as a continuous hydrodynamic engine operating within the closed continuous metric. To satisfy Eulerian continuity ($\nabla \cdot \mathbf{v} = 0$), the geometric boundary acts as a topological stator, driving a continuous, localized fluid circulation loop that natively generates both scalar pressure deficits and rotational impedance.

\textbf{1. The Fundamental Dipole (Leptons):} \\
By the Poincaré-Hopf theorem, a continuous tangent vector field on a spherical boundary mandates two points where the rotational vector evaluates to zero ($\mathbf{v}_{\circlearrowleft} = 0$). These singular poles define the topological axis of the closed circulation loop. 

As the circulating fluid traverses the polar core, it curls orthogonally to establish the equatorial perimeter ($\mathbf{v}_{\circlearrowleft} = c e^{-\bar{\lambda}_c/R_h}$). To maintain   zero volumetric divergence, the fluid evaluates as a continuous, unbroken torus of internal topological vorticity ($\mathbf{v}_{\circlearrowleft}$). 
\begin{equation}
    \nabla \cdot \mathbf{v}_{\circlearrowleft} = 0 \tag{\ref{eq:engine_transduction}}
\end{equation}
The fundamental lepton evaluates physically as this continuous hydrodynamic unknot. The scalar Bernoulli deficit of the closed loop establishes rest mass, the absolute phase limit establishes the geometric boundary, and the off-diagonal tensor coupling to the ambient transverse plane establishes the electromagnetic field. The geometric orientation of the polar axis establishes the absolute topological spin state.

\textbf{2. The Composite Trefoil (Baryons):} \\
To satisfy volumetric continuity for fractional lobes, exactly three orthogonal fluid tubes must topologically interlock. They weave around a central void to form a closed Trefoil knot. 

The high-pressure circulation of one fractional tube evaluates directly into the flow path of the next. Deep within the core, the internal streams evaluate as parallel and laminar ($\nabla \times \mathbf{v}_{core} = 0$), providing frictionless internal circulation (Asymptotic Freedom). At the geometric perimeter, the intersecting tubes force violent Kolmogorov turbulence ($\nabla \times \mathbf{v}_{surface} \neq 0$), weaving the rigid structural mass boundary (Topological Confinement).

\textbf{3. Hydrostatic Supremacy (The Legendre Phase Transformation):} \\
To prevent the macroscopic static pressure from crushing the void, the circulating engine must generate an outward restorative dynamic pressure ($P_{out}$). To evaluate the total work performed across the boundary without encountering a Cartesian radial singularity ($r \to 0$), the power density Lagrangian ($\mathcal{L}$) undergoes a Legendre transformation. 

The independent Cartesian spatial coordinate ($\gamma_z$) is swapped for the conjugate rotational phase angle ($\theta$), projecting the energy into the macroscopic Hamiltonian ($\mathcal{H}$):
\begin{equation}
    \mathcal{H}(\theta, \tau_{\theta z}) = \theta \tau_{\theta z} - \mathcal{L}(\gamma_z) \label{eq:power_legendre}
\end{equation}
This transformation mathematically dictates that the localized 1D inward hydrostatic crush ($P_{in}$) is geometrically mapped against a continuous angular differential. Because the topological boundary evaluates as a closed circulation loop, the total restorative force evaluates as the boundary integral across one full geometric phase cycle:
\begin{align}
    P_{out} &= \oint_{\partial S} P_{in} \, d\theta  \\
    P_{out} &= P_{in} \int_{0}^{2\pi} d\theta \nonumber \\
    P_{out} &= P_{in} \Big[ \theta \Big]_{0}^{2\pi} \nonumber \\
    P_{out} &= 2\pi P_{in} \nonumber
\end{align}
The topological phase loop mechanically amplifies the localized dynamic pressure by exactly $2\pi$. This internal restoring force overpowers the macroscopic vacuum crush, guaranteeing absolute hydrostatic boundary stability:
\begin{equation}
    \frac{P_{out}}{P_{in}} = 2\pi \tag{\ref{eq:engine_supremacy}}
\end{equation}

\textbf{4. The Planck Redline ($l_p \implies m_p$):} \\
By the Inverse-Cube Identity ($V_R \propto 1/m^3$), as mass increases, the Euclidean void shrinks. The topological stator tightens, forcing the circulating fluid to accelerate over a steeper spatial gradient. 

As the void compresses toward the Planck Length ($l_p$), the circulating fluid approaches the absolute acoustic phase limit of the continuum ($v \to c$). The physical cavitation lockout occurs when the core flow evaluates at $c$, defining the fundamental Planck Mass ($m_p$) as the maximum stable geometric displacement of a single topological fluid engine. 
\begin{equation}
    \lim_{V_R \to l_p^3} v_{core} \to c \implies m_{max} = m_p \tag{\ref{eq:engine_redline}}
\end{equation}
Shrinking the stator further requires the fluid to exceed $c$ to maintain the hydrostatic lock, which the continuum algebraically forbids.

\textbf{5. Hydrostatic Instability ($m^4$ Decay Ratio):} \\
To satisfy continuous mass conservation, the Euclidean core volume ($V_R$) must thermodynamically balance the geometric volume of the circulating hydrodynamic halo ($V_\tau$). 
By the Kinematic Origin of Inertia, observable mass evaluates as the hydrodynamic added mass:
\begin{equation}
    m_\tau = \frac{1}{2} \rho_{\tau} V_\tau \implies V_\tau = \frac{2m_\tau}{\rho_{\tau}} \tag{\ref{eq:mass_volume_equivalence}}
\end{equation}
The Euclidean bounding void evaluates as the macroscopic 3D spherical box using the topological scale radius ($R_h = \frac{\hbar}{m_\tau c}$):
\begin{equation}
    V_R = \frac{4}{3}\pi \left(\frac{\hbar}{m_\tau c}\right)^3 = \frac{4\pi\hbar^3}{3m_\tau^3 c^3} 
\end{equation}
The geometric ratio of the required circulation footprint to the available core void evaluates as the algebraic quotient:
\begin{equation}
    \frac{V_\tau}{V_R} = \frac{\frac{2m_\tau}{\rho_{\tau}}}{\frac{4\pi\hbar^3}{3m_\tau^3 c^3}} = \left(\frac{3c^3}{2\pi\rho_{\tau}\hbar^3}\right) m_\tau^4 \tag{\ref{eq:engine_decay}}
\end{equation}
As mass increases, the required geometry to house the circulating topological halo grows to the fourth power of the mass, while the primary displaced core void shrinks. At extreme mass scales, the localized fluid pressure of this $m^4$ circulation halo exceeds the restorative yield strength of the metric. The stator is ripped open by its own localized dynamic pressure, fracturing the boundary into lower-mass geometries. This continuous fluid-dynamic yield derives the absolute mechanism of particle decay.

\begin{tcolorbox}[recallbox, title={\ref{box:unified_engine}. The Unified Topological Engine}]
\begin{align}
    \nabla \cdot \mathbf{v}_{\circlearrowleft} &= 0 \tag{\ref{eq:engine_transduction}} \\
    \mathcal{H}(\theta, \tau_{\theta z}) &= \theta \tau_{\theta z} - \mathcal{L}(\gamma_z) \tag{\ref{eq:power_legendre}} \\
    \frac{P_{out}}{P_{in}} &= 2\pi \tag{\ref{eq:engine_supremacy}} \\
    \lim_{V_R \to l_p^3} v_{core} &\to c \implies m_{max} = m_p \tag{\ref{eq:engine_redline}} \\
    \frac{V_\tau}{V_R} &= \left(\frac{3c^3}{2\pi\rho_{\tau}\hbar^3}\right) m_\tau^4 \tag{\ref{eq:engine_decay}}
\end{align}
\end{tcolorbox}


% ----------------------------------------------------------------------
\section{Fluid Circulation \& The Topological Origin of Spin}
\label{sec:deriv_quantum_spin}

In continuous fluid mechanics, localized shear evaluates mathematically as fluid vorticity ($\boldsymbol{\omega}_{fluid} = \nabla \times \mathbf{u} \neq 0$). Fundamental matter evaluates as a stable topological defect within the tensioned metric, sustaining continuous fluid circulation to prevent unraveling.

\textbf{1. Kinematic Angular Momentum \& The Impedance Limit ($\hbar$):} \\
A localized vortex defines a permanent state of internal topological vorticity ($\nabla \times \mathbf{v}_{\psi}$). The total intrinsic angular momentum ($\mathbf{S}$) evaluates as the volume integral of the fluid's kinematic momentum density crossed with the radial geometric boundary:
\begin{equation}
    \mathbf{S} = \iiint_V \rho_{fluid} (\mathbf{r} \times \mathbf{v}_{\circlearrowleft}) \, dV \nonumber
\end{equation}
At the topological core, the localized fluid evaluates at the maximum acoustic circulation limit ($v_{\circlearrowleft} = c$) at the spatial boundary radius ($r = R_h$). Integrating the internal momentum density for the entire displaced fluid mass ($m$) evaluates the macroscopic scalar angular momentum of the core ($S = mcR_h$). 

Because the localized angular momentum cannot physically exceed the absolute Rotational Impedance of the macroscopic vacuum ($Z_{\tau(\circlearrowleft)} \equiv \hbar$), the spatial boundary is structurally restricted to the exact inverse acoustic phase radius ($R_h = \hbar / mc$). Substituting this boundary limit isolates the absolute continuum scalar:
\begin{align}
    S &= m c \left( \frac{\hbar}{mc} \right) \nonumber \\
    S &= \hbar \tag{\ref{eq:transverse_shear_spin}}
\end{align}

\textbf{2. The Baseline Circulation Limit:} \\
By Kelvin's Circulation Theorem, the total circulation ($\Gamma$) of a closed fluid path evaluates as the closed line integral of the fluid velocity ($\mathbf{v}$) around the contour ($C$):
\begin{equation}
    \Gamma = \oint_C \mathbf{v}_{\circlearrowleft} \cdot d\mathbf{l} 
\end{equation}
For a localized acoustic cavity, the fluid at the boundary is locked to the continuous acoustic limit ($v_{\circlearrowleft} = c$). Evaluating the internal vorticity, the line integral simplifies to the limit velocity multiplied by the physical circumference of the boundary ($C = 2\pi R_h$):
\begin{align}
    \Gamma &= c \oint_C d\mathbf{l}  \\
    \Gamma &= c (2\pi R_h) 
\end{align}
Substituting the explicit topological scale radius ($R_h = \hbar / mc$):
\begin{align}
    \Gamma &= c (2\pi) \left( \frac{\hbar}{m c} \right)  \\
    \Gamma &= \frac{2\pi \hbar}{m} 
\end{align}
Because the reduced Planck constant expands algebraically as $\hbar = h / 2\pi$, substituting this isolates the absolute macroscopic circulation:
\begin{equation}
    \Gamma = \frac{h}{m} \tag{\ref{eq:fluid_circulation}}
\end{equation}

\textbf{3. The Spinor Topology ($4\pi$ Metric Untangling):} \\
To translate fluid circulation into quantized angular momentum ($S$), the circulation must be distributed over the geometric radian slice of the boundary. 

If the core existed in a discrete, empty void, it would evaluate via a standard $2\pi$ rotation. However, the metric evaluates as a continuous medium. The topological core is permanently connected to the macroscopic environment via the continuous hydrostatic pressure gradient. To rotate the topological stator without snapping the connecting fluid metric, the phase geometry must evaluate as a Dirac Spinor. 

The fluid requires exactly two full rotational phase wraps ($4\pi$ radians) to untangle the surrounding metric and return the boundary to its identity state. The intrinsic Angular Momentum ($S_{fermion}$) evaluates as the mass density multiplied by the circulation over this $4\pi$ untangling slice:
\begin{align}
    S_{fermion} &= m \left( \frac{\Gamma}{4\pi} \right)  \\
    S_{fermion} &= \frac{m}{4\pi} \left( \frac{h}{m} \right)  \\
    S_{fermion} &= \frac{h}{4\pi} 
\end{align}
Substituting $h = 2\pi\hbar$ isolates the absolute derivation of the Fermionic spin anomaly:
\begin{align}
    S_{fermion} &= \frac{2\pi\hbar}{4\pi}  \\
    S_{fermion} &= \frac{\hbar}{2} \tag{\ref{eq:spinor_topology}}
\end{align}
Fermionic half-integer spin evaluates not as an intrinsic quantum property, but as the continuous fluid requirement to double-wrap the topological stator to maintain geometric equilibrium with the continuous $\boldsymbol{\tau}$-fluid metric.

\begin{tcolorbox}[recallbox, title={\ref{box:quantum_spin}. Fluid Circulation \& The Topological Origin of Spin}]
\begin{align}
    \mathbf{S} = \iiint_V \rho (\mathbf{r} \times \mathbf{v}_{\circlearrowleft}) \, dV &\implies S = \hbar \tag{\ref{eq:transverse_shear_spin}} \\
    \Gamma &= \oint_C \mathbf{v}_{\circlearrowleft} \cdot d\mathbf{l} = \frac{h}{m} \tag{\ref{eq:fluid_circulation}} \\
    S_{fermion} &= m \left( \frac{\Gamma}{4\pi} \right) = \frac{\hbar}{2} \tag{\ref{eq:spinor_topology}}
\end{align}
\end{tcolorbox}

% ----------------------------------------------------------------------
\section{Kinetic-Potential Strain Conservation}
\label{sec:deriv_volumetric_energy}

Absolute energy evaluates as the displaced fluid volume of the continuous metric. Potential and kinetic energies evaluate as physical, geometric partitions of the localized total scalar admittance ($\alpha_s$).

\textbf{1. Kinetic Strain (The Dynamic Bow Shock):} \\
A topological rest boundary displaces a hydrodynamic added mass volume ($V_{added}$). Under spatial translation, the boundary compresses the surrounding fluid metric. 

The kinematic drag factor ($\gamma$) evaluates as the inverse of the localized total scalar admittance ($\alpha_s$):
\begin{equation}
    \gamma = \frac{1}{\alpha_s} \tag{\ref{eq:strain_lorentz}}
\end{equation}

The physical volume of dynamic kinematic translation ($V_{kin}$) evaluates as the effective dilated volume ($\gamma V_{added}$) minus the invariant baseline rest volume ($V_{added}$). Substituting the admittance ratio expands the kinetic displacement volume as a function of fluid compression:
\begin{align}
    V_{kin} &= \gamma V_{added} - V_{added}  \\
    V_{kin} &= \left( \frac{1}{\alpha_s} \right) V_{added} - V_{added}  \\
    V_{kin} &= \left( \frac{1}{\alpha_s} - 1 \right) V_{added} \tag{\ref{eq:kinetic_volume}}
\end{align}
Multiplying by the baseline macroscopic kinetic energy density of the fluid ($\frac{1}{2}\rho_{\tau} c^2$) yields the macroscopic kinetic energy limit ($E_k = (\gamma - 1)m_0 c^2$).

\textbf{2. Potential Energy (Static Elastic Deficit):} \\
Translating a topological boundary against a macroscopic pressure gradient ($\nabla \Phi$) performs thermodynamic work against the continuous restorative force of the metric ($\mathbf{F}$). The physical volume of fluid displaced into this static elastic strain ($V_{pot}$) evaluates via the energetic work integral, scaled by the baseline kinetic energy density ($\frac{1}{2}\rho_{\tau} c^2$):
\begin{equation}
    V_{pot} = \frac{2}{\rho_{\tau} c^2} \int_{x_1}^{x_2} \mathbf{F} \cdot d\mathbf{x} \tag{\ref{eq:potential_volume}}
\end{equation}
Potential energy evaluates as the geometric volume of fluid displaced into the macroscopic elastic tension of the continuous environment. 

\textbf{3. Isochoric Transmutation (Conservation of Energy):} \\
By the Eulerian continuity equation, the total displaced fluid volume ($V_{total}$) of a closed topological system evaluates as an absolute constant:
\begin{equation}
    V_{total} = V_{added} + V_{kin} + V_{pot} = \text{Constant} 
\end{equation}

Taking the proper time derivative of the absolute volume isolates the volumetric phase exchange. Because the intrinsic added mass of the primary defect evaluates as topologically invariant ($\dot{V}_{added} = 0$), the derivative expands as:
\begin{align}
    \frac{d}{dt} (V_{total}) &= \frac{d}{dt} (V_{added} + V_{kin} + V_{pot}) \nonumber \\
    0 &= \dot{V}_{added} + \dot{V}_{kin} + \dot{V}_{pot} \nonumber \\
    0 &= 0 + \dot{V}_{kin} + \dot{V}_{pot} \nonumber \\
    \dot{V}_{kin} &= -\dot{V}_{pot} \tag{\ref{eq:isochoric_transmutation}}
\end{align}

When a boundary decelerates against a gravitational gradient, the localized fluid compression of its dynamic bow shock relaxes ($\dot{V}_{kin} < 0$). This fluid compliance shifts into the surrounding metric, inflating the static elastic strain of the gravitational well ($\dot{V}_{pot} > 0$). 

\vspace{0.2cm}
\textbf{The Topographical Shift:} \\
Conservation of energy evaluates as the continuous isochoric fluid exchange between localized dynamic boundary strain and macroscopic static elastic tension. In a gravitational orbit, the metric compression shifts from the longitudinal vector of the falling mass ($\mathbf{v}_{translation}$) into the topographical static pressure gradient of the planet ($\nabla \mathcal{P}_{ambient}$).

\begin{tcolorbox}[recallbox, title={\ref{box:volumetric_energy}. Kinetic-Potential Strain Conservation}]
\begin{align}
    \gamma &= \frac{1}{\alpha_s} \tag{\ref{eq:strain_lorentz}} \\
    V_{kin} &= \left( \frac{1}{\alpha_s} - 1 \right) V_{added} \tag{\ref{eq:kinetic_volume}} \\
    V_{pot} &= \frac{2}{\rho_{\tau} c^2} \int_{x_1}^{x_2} \mathbf{F} \cdot d\mathbf{x} \tag{\ref{eq:potential_volume}} \\
    \dot{V}_{kin} &= -\dot{V}_{pot} \tag{\ref{eq:isochoric_transmutation}}
\end{align}
\end{tcolorbox}



% ======================================================================
% PHASE IV: PHASE IV: WAVE KINEMATICS & ELECTROMAGNETISM
% ======================================================================
\part*{PHASE IV: WAVE KINEMATICS \& ELECTROMAGNETISM}








% ----------------------------------------------------------------------
\section{The Unified Fluid Lagrangian \& Modal Admittance}
\label{sec:deriv_unified_lagrangian}

The continuous mechanics of the metric evaluate as a unified thermodynamic fluid action. Because the Helmholtz decomposition partitions the total kinematic field into orthogonal geometries, the fluid metric evaluates as an anisotropic viscoelastic tensor governed by Modal Admittance.

\textbf{1. The Absolute Kinematic Capacity ($v_{tot}^2 \le c^2$):} \\
The total fluid velocity vector partitions into three orthogonal driving fields: the longitudinal scalar gradient ($\nabla h$), the transverse vector circulation ($\nabla \times \mathbf{A}$), and the localized internal topological vorticity ($\nabla \times \mathbf{v}_{\psi}$). 

Because the longitudinal scalar gradient evaluates   irrotational fluid compression, it tracks the kinematic velocities driven by both the ambient macroscopic volume pressure (Gravity, $\nabla \mathcal{P}_{ambient}$) and the localized dynamic surface pressure (Propulsion, $\nabla \mathcal{P}_{local}$). The total squared velocity ($v_{tot}^2$) evaluates the localized kinematic capacity of the medium, bounded by $c^2$:
\begin{equation}
    v_{tot}^2 = (\nabla h)^2 + (\nabla \times \mathbf{A})^2 + (\nabla \times \mathbf{v}_{\psi})^2 \le c^2 \tag{\ref{eq:absolute_pythagorean_constraint}}
\end{equation}
In the macroscopic vacuum, the kinematic velocity evaluates near zero ($v_{tot}^2 \ll c^2$). Where the kinematic fields maximize, the fluid saturates the capacity limit ($v_{tot}^2 \to c^2$).

\textbf{2. Modal Admittance ($\alpha_{i}$) \& Fluid Compliance:} \\
Because the kinematic components strain the fluid in orthogonal dimensions, they consume the absolute phase-space limit ($c$) independently. The fluid admittance evaluates as a discrete modal limit for each geometric strain:
\begin{align}
    \alpha_{\parallel} &= \sqrt{1 - \frac{(\nabla h)^2}{c^2}} \nonumber \\
    \alpha_{\perp} &= \sqrt{1 - \frac{(\nabla \times \mathbf{A})^2}{c^2}} \nonumber \\
    \alpha_{\circlearrowleft} &= \sqrt{1 - \frac{(\nabla \times \mathbf{v}_{\psi})^2}{c^2}} \tag{\ref{eq:modal_admittance}}
\end{align}

\textbf{3. Topological Confinement (The Infinite Deborah Number):} \\
The Deborah number ($De$) defines the phase transition between fluid compliance and solid rigidity. As geometric compliance ($\alpha$) drops to zero, the structural relaxation time diverges ($De \to \infty$). 

Where the internal rotational velocity saturates the medium ($|\nabla \times \mathbf{v}_{\psi}| \to c$), the rotational admittance mathematically vanishes:
\begin{equation}
    |\nabla \times \mathbf{v}_{\psi}| \to c \implies \alpha_{\circlearrowleft} \to 0 \implies De_{rot} \to \infty \tag{\ref{eq:deborah_rotational}}
\end{equation}
Where rotational compliance evaluates to zero, the metric loses capability to yield to shear and behaves as a rigid solid. Because macroscopic translational velocities remain below the kinematic limit ($(\nabla h)^2 \ll c^2$), translational admittances remain open ($\alpha_{\parallel} \approx 1$). A continuous field evaluates as an unyielding solid exclusively in its saturated axis.

\textbf{4. The Localized Compliance Limit:} \\
Squaring the modal admittances and summing their states reveals the viscoelastic geometry:
\begin{align}
    \Sigma \alpha_i^2 &= \left(1 - \frac{(\nabla h)^2}{c^2}\right) + \left(1 - \frac{(\nabla \times \mathbf{A})^2}{c^2}\right) + \left(1 - \frac{(\nabla \times \mathbf{v}_{\psi})^2}{c^2}\right) \nonumber \\
    \Sigma \alpha_i^2 &= 3 - \frac{v_{tot}^2}{c^2} \tag{\ref{eq:compliance_sum}}
\end{align}
Evaluating this sum where the kinematic capacity saturates ($v_{tot} \to c$):
\begin{equation}
    \Sigma \alpha_i^2 = 3 - \frac{c^2}{c^2} = 2 \tag{\ref{eq:absolute_compliance_state}}
\end{equation}
To maintain the phase limit, the localized metric possesses two fully compliant degrees of kinematic freedom; the saturated dimension structurally freezes ($\alpha = 0$) as an energetic boundary.

\textbf{5. The Dynamic Field Lagrangian ($\mathcal{L}_{dynamic}$):} \\
The dynamic action evaluates the kinetic energy density ($\mathcal{K}$) minus the structural potential energy density ($\mathcal{P}$).

The kinetic energy evaluates as the sum of the orthogonal velocity strains scaled by respective modal admittances:
\begin{equation}
    \mathcal{K} = \frac{1}{2} \rho_{\tau} \Big[ \alpha_{\parallel}(\nabla h)^2 + \alpha_{\perp}(\nabla \times \mathbf{A})^2 + \alpha_{\circlearrowleft}(\nabla \times \mathbf{v}_{\psi})^2 \Big] \nonumber
\end{equation}

The structural potential energy consists of the absolute macroscopic vacuum tension ($\rho_{\tau} c^2$) and the localized topological acoustic cavitation ($\frac{1}{2} \rho_{\tau} c^2$). The absolute structural potential ($\mathcal{P}$) evaluates as the difference:
\begin{equation}
    \mathcal{P} = \rho_{\tau} c^2 - \frac{1}{2} \rho_{\tau} c^2 = \frac{1}{2} \rho_{\tau} c^2 \nonumber
\end{equation}

Subtracting the potential energy from the kinetic energy constructs the field Lagrangian:
\begin{equation}
    \mathcal{L}_{dynamic} = \frac{1}{2} \rho_{\tau} \Big[ \alpha_{\parallel}(\nabla h)^2 + \alpha_{\perp}(\nabla \times \mathbf{A})^2 + \alpha_{\circlearrowleft}(\nabla \times \mathbf{v}_{\psi})^2 \Big] - \frac{1}{2} \rho_{\tau} c^2 \tag{\ref{eq:modal_lagrangian}}
\end{equation}

\textbf{6. The Absolute Kinematic Lagrangian ($\mathcal{L}_{kin}$):} \\
To isolate the pure geometric action of the phase space, the absolute background mass density ($\rho_{\tau}$) is factored out of the dynamic Lagrangian. 

Dividing the dynamic Lagrangian by $\rho_{\tau}$ mathematically annihilates the mass term, yielding the Absolute Kinematic Lagrangian in native spatiotemporal dimensions ($\text{m}^2/\text{s}^2$):
\begin{equation}
    \mathcal{L}_{kin} = \frac{1}{2} \Big[ \alpha_{\parallel}(\nabla h)^2 + \alpha_{\perp}(\nabla \times \mathbf{A})^2 + \alpha_{\circlearrowleft}(\nabla \times \mathbf{v}_{\psi})^2 \Big] - \frac{1}{2} c^2 \tag{\ref{eq:kinematic_lagrangian}}
\end{equation}
The geometric action of the continuum evaluates exclusively as scale-invariant phase-space evolution, independent of the numerical magnitude of the macroscopic fluid density.

\textbf{7. The Universal Angular Phase-Space Lagrangian:} \\
Defining the total velocity $v_{tot}$ via spherical coordinates extracts the geometric phase space for any arbitrary magnitude:
\begin{align}
    (\nabla h)^2 &= v_{tot}^2 \cos^2(\theta) \nonumber \\
    (\nabla \times \mathbf{A})^2 &= v_{tot}^2 \sin^2(\theta)\cos^2(\phi) \nonumber \\
    (\nabla \times \mathbf{v}_{\psi})^2 &= v_{tot}^2 \sin^2(\theta)\sin^2(\phi) \nonumber
\end{align}

Substituting these continuous velocity components into the Modal Admittances yields the dynamic compliance states:
\begin{align}
    \alpha_{\parallel} &= \sqrt{1 - \frac{v_{tot}^2}{c^2} \cos^2(\theta)} \nonumber \\
    \alpha_{\perp} &= \sqrt{1 - \frac{v_{tot}^2}{c^2} \sin^2(\theta)\cos^2(\phi)} \nonumber \\
    \alpha_{\circlearrowleft} &= \sqrt{1 - \frac{v_{tot}^2}{c^2} \sin^2(\theta)\sin^2(\phi)} \nonumber
\end{align}

Substituting these explicit angular components back into the Kinematic Lagrangian ($\mathcal{L}_{kin}$) factors the total kinematic intensity ($v_{tot}^2$) from the topology. Writing the full explicit expansion reveals the pure geometric phase space of the continuous fluid:
\begin{align}
    \mathcal{L}_{kin}(v_{tot}, \theta, \phi) = \frac{1}{2} v_{tot}^2 \Big[ &\sqrt{1 - \frac{v_{tot}^2}{c^2} \cos^2(\theta)} \cos^2(\theta) \nonumber \\
    &+ \sqrt{1 - \frac{v_{tot}^2}{c^2} \sin^2(\theta)\cos^2(\phi)} \sin^2(\theta)\cos^2(\phi) \nonumber \\
    &+ \sqrt{1 - \frac{v_{tot}^2}{c^2} \sin^2(\theta)\sin^2(\phi)} \sin^2(\theta)\sin^2(\phi) \Big] \nonumber \\
    &- \frac{1}{2} c^2 \tag{\ref{eq:angular_lagrangian}}
\end{align}

\begin{tcolorbox}[recallbox, title={\ref{box:unified_lagrangian}. The Unified Fluid Lagrangian \& Modal Admittance}]
\begin{align}
    v_{tot}^2 &= (\nabla h)^2 + (\nabla \times \mathbf{A})^2 + (\nabla \times \mathbf{v}_{\psi})^2 \le c^2 \tag{\ref{eq:absolute_pythagorean_constraint}} \\
    \alpha_{\parallel} &= \sqrt{1 - \frac{(\nabla h)^2}{c^2}}, \quad \alpha_{\perp} = \sqrt{1 - \frac{(\nabla \times \mathbf{A})^2}{c^2}}, \quad \alpha_{\circlearrowleft} = \sqrt{1 - \frac{(\nabla \times \mathbf{v}_{\psi})^2}{c^2}} \tag{\ref{eq:modal_admittance}} \\
    |\nabla \times \mathbf{v}_{\psi}| &\to c \implies \alpha_{\circlearrowleft} \to 0 \implies De_{rot} \to \infty \tag{\ref{eq:deborah_rotational}} \\
    \Sigma \alpha_i^2 &= 3 - \frac{v_{tot}^2}{c^2} \tag{\ref{eq:compliance_sum}} \\
    v_{tot} &\to c \implies \Sigma \alpha_i^2 = 2 \tag{\ref{eq:absolute_compliance_state}} \\
    \mathcal{L}_{dynamic} &= \frac{1}{2} \rho_{\tau} \Big[ \alpha_{\parallel}(\nabla h)^2 + \alpha_{\perp}(\nabla \times \mathbf{A})^2 + \alpha_{\circlearrowleft}(\nabla \times \mathbf{v}_{\psi})^2 \Big] - \frac{1}{2} \rho_{\tau} c^2 \tag{\ref{eq:modal_lagrangian}} \\
    \mathcal{L}_{kin} &= \frac{1}{2} \Big[ \alpha_{\parallel}(\nabla h)^2 + \alpha_{\perp}(\nabla \times \mathbf{A})^2 + \alpha_{\circlearrowleft}(\nabla \times \mathbf{v}_{\psi})^2 \Big] - \frac{1}{2} c^2 \tag{\ref{eq:kinematic_lagrangian}} \\
    \mathcal{L}_{kin}(v_{tot}, \theta, \phi) &= \frac{1}{2} v_{tot}^2 \Big[ \sqrt{1 - \frac{v_{tot}^2}{c^2} \cos^2(\theta)} \cos^2(\theta) \nonumber \\
    &\quad + \sqrt{1 - \frac{v_{tot}^2}{c^2} \sin^2(\theta)\cos^2(\phi)} \sin^2(\theta)\cos^2(\phi) \nonumber \\
    &\quad + \sqrt{1 - \frac{v_{tot}^2}{c^2} \sin^2(\theta)\sin^2(\phi)} \sin^2(\theta)\sin^2(\phi) \Big] \nonumber \\
    &\quad - \frac{1}{2} c^2 \tag{\ref{eq:angular_lagrangian}}
\end{align}
\end{tcolorbox}
% ----------------------------------------------------------------------
\section{The Legendre Hamiltonian Translation}
\label{sec:deriv_legendre_hamiltonian}

The continuous fluid dynamics evaluated by the Unified Lagrangian ($\mathcal{L}_{\boldsymbol{\tau}} = \mathcal{K} - \mathcal{P}$) dictate kinematic action. Extracting static topological observables requires translating the continuous system from velocity phase-space into momentum phase-space via the macroscopic Legendre transformation.

\textbf{1. Macroscopic Conjugate Momentum:} \\
The localized conjugate momentum density ($\mathbf{p}_{macro}$) evaluates as the partial derivative of the baseline Lagrangian with respect to the macroscopic driving fluid velocity ($\mathbf{v}_{macro} \equiv \nabla h$). Because the structural potential energy density ($\mathcal{P}$) is independent of velocity, the derivative isolates the kinematic density:
\begin{align}
    \mathbf{p}_{macro} &= \frac{\partial \mathcal{L}_{\boldsymbol{\tau}}}{\partial \mathbf{v}_{macro}} \nonumber \\
    \mathbf{p}_{macro} &= \frac{\partial}{\partial \mathbf{v}_{macro}} \left( \frac{1}{2} \rho_{\tau} \mathbf{v}_{macro}^2 - \frac{1}{2} \rho_{\tau} c^2 \right) \nonumber \\
    \mathbf{p}_{macro} &= \rho_{\tau} \mathbf{v}_{macro} \tag{\ref{eq:conjugate_momentum}}
\end{align}

\textbf{2. The Baseline Legendre Transform:} \\
The baseline Hamiltonian density ($\mathcal{H}_{\boldsymbol{\tau}}$) evaluates as the Legendre transform of the Lagrangian, defined as the dot product of the macroscopic velocity and conjugate momentum minus the Lagrangian action:
\begin{align}
    \mathcal{H}_{\boldsymbol{\tau}} &= \mathbf{p}_{macro} \cdot \mathbf{v}_{macro} - \mathcal{L}_{\boldsymbol{\tau}} \nonumber \\
    \mathcal{H}_{\boldsymbol{\tau}} &= (\rho_{\tau} \mathbf{v}_{macro}) \cdot \mathbf{v}_{macro} - (\mathcal{K}_{macro} - \mathcal{P}) \nonumber \\
    \mathcal{H}_{\boldsymbol{\tau}} &= \rho_{\tau} \mathbf{v}_{macro}^2 - \mathcal{K}_{macro} + \mathcal{P} \nonumber
\end{align}
Because total macroscopic kinetic energy density is defined as $\mathcal{K}_{macro} = \frac{1}{2}\rho_{\tau} \mathbf{v}_{macro}^2$, the momentum term substitutes exactly as $2\mathcal{K}_{macro}$:
\begin{align}
    \mathcal{H}_{\boldsymbol{\tau}} &= 2\mathcal{K}_{macro} - \mathcal{K}_{macro} + \mathcal{P} \nonumber \\
    \mathcal{H}_{\boldsymbol{\tau}} &= \mathcal{K}_{macro} + \mathcal{P} \tag{\ref{eq:legendre_transform}}
\end{align}

\textbf{3. The Covariant Bernoulli State:} \\
Substituting the baseline kinetic and structural potential energy densities ($\mathcal{K}_{macro} = \frac{1}{2}\rho_{\tau} \mathbf{v}_{macro}^2$ and $\mathcal{P} = \frac{1}{2} \rho_{\tau} c^2$) translates the continuous Lagrangian explicitly into the Bernoulli steady-state equation:
\begin{equation}
    \mathcal{H}_{\boldsymbol{\tau}} = \frac{1}{2} \rho_{\tau} \Big[ \mathbf{v}_{macro}^2 + c^2 \Big] \tag{\ref{eq:bernoulli_hamiltonian}}
\end{equation}
The baseline energy state of the universe evaluates as the absolute sum of macroscopic dynamic pressure ($\frac{1}{2}\rho_{\tau} \mathbf{v}_{macro}^2$) and structural hydrostatic yield pressure ($\frac{1}{2}\rho_{\tau} c^2$).

\textbf{4. Modal Energy Extraction ($\mathcal{H}_{dyn}$ \& $\mathcal{H}_{kin}$):} \\
To extract the dynamic energy state of a discrete topology, the Modal Admittances ($\alpha_{i}$) are applied to the geometric velocity fields. The Dynamic Hamiltonian ($\mathcal{H}_{dyn}$) sums the compliant kinetic strains and the structural potential:
\begin{equation}
    \mathcal{H}_{dyn} = \frac{1}{2} \rho_{\tau} \Big[ \alpha_{\parallel}(\nabla h)^2 + \alpha_{\perp}(\nabla \times \mathbf{A})^2 + \alpha_{\circlearrowleft}(\nabla \times \mathbf{v}_{\psi})^2 \Big] + \frac{1}{2} \rho_{\tau} c^2 \tag{\ref{eq:modal_hamiltonian}}
\end{equation}
To evaluate the geometric phase-space state (Specific Energy), the absolute mass density ($\rho_{\tau}$) is mathematically annihilated from the Hamiltonian. Dividing $\mathcal{H}_{dyn}$ by $\rho_{\tau}$ isolates the Kinematic Hamiltonian ($\mathcal{H}_{kin}$) in native spatiotemporal dimensions ($\text{m}^2/\text{s}^2$):
\begin{equation}
    \mathcal{H}_{kin} = \frac{1}{2} \Big[ \alpha_{\parallel}(\nabla h)^2 + \alpha_{\perp}(\nabla \times \mathbf{A})^2 + \alpha_{\circlearrowleft}(\nabla \times \mathbf{v}_{\psi})^2 \Big] + \frac{1}{2} c^2 \tag{\ref{eq:kinematic_hamiltonian}}
\end{equation}

\textbf{5. The Rotational Saturation Limit ($\mathcal{H}_{rest}$):} \\
Where the internal rotational velocity saturates the medium ($|\nabla \times \mathbf{v}_{\psi}| \to c$), the internal rotational admittance vanishes ($\alpha_{\circlearrowleft} \to 0$).

Evaluating a saturated boundary at macroscopic rest ($\mathbf{v}_{macro} = 0$), the dynamic kinetic summation evaluates identically to zero. The Hamiltonian isolates   to the absolute structural potential:
\begin{equation}
    \mathcal{H}_{rest} = \frac{1}{2} \rho_{\tau} c^2 \tag{\ref{eq:hamiltonian_rest_mass}}
\end{equation}
The internal rotational energy evaluates as permanently frozen into the rest mass potential. 

\textbf{6. The Universal Angular Phase-Space Hamiltonian:} \\
Defining the total velocity $v_{tot}$ via spherical coordinates extracts the geometric topology from the kinematic intensity ($v_{tot}^2$). Substituting the expanded angular components into the mass-independent Kinematic Hamiltonian ($\mathcal{H}_{kin}$) yields:
\begin{align}
    \mathcal{H}_{kin}(v_{tot}, \theta, \phi) = \frac{1}{2} v_{tot}^2 \Big[ &\alpha_{\parallel} \cos^2(\theta) + \alpha_{\perp} \sin^2(\theta)\cos^2(\phi) \nonumber \\
    &+ \alpha_{\circlearrowleft} \sin^2(\theta)\sin^2(\phi) \Big] + \frac{1}{2} c^2 \tag{\ref{eq:angular_hamiltonian}}
\end{align}

\textbf{7. Topological Mass Extraction:} \\
While kinematics evaluates as massless geometry, static topological observation requires coupling the geometry to the absolute baseline density ($\rho_{\tau}$). The resting dynamic Hamiltonian evaluates   as the absolute acoustic yield pressure of the continuous fluid ($\mathcal{H}_{rest} = P_c = \frac{1}{2}\rho_{\tau} c^2$). 

To extract the total rest energy ($E_{rest}$), this pressure is integrated across the spatial volume of the localized fluid void ($V_{\psi}$):
\begin{equation}
    E_{rest} = \int \mathcal{H}_{rest} \, dV = \frac{1}{2} \rho_{\tau} c^2 V_{\psi} \tag{\ref{eq:rest_energy_integral}}
\end{equation}
Applying classical mass-energy equivalence ($m = E/c^2$) translates the continuous resting energy into the discrete observable mass:
\begin{align}
    m_0 &= \frac{\frac{1}{2} \rho_{\tau} c^2 V_{\psi}}{c^2} \nonumber \\
    m_0 &= \frac{1}{2} \rho_{\tau} V_{\psi} \tag{\ref{eq:topological_mass}}
\end{align}
Discrete rest mass evaluates absolutely as the localized volumetric pressure deficit of the metric. Observable inertia is exclusively defined by the volumetric displacement of the void ($V_{\psi}$) scaled against the baseline macroscopic mass density of the continuous vacuum ($\rho_{\tau}$).

\begin{tcolorbox}[recallbox, title={\ref{box:legendre_hamiltonian}. The Legendre Hamiltonian Translation}]
\begin{align}
    \mathbf{p}_{macro} &= \rho_{\tau} \mathbf{v}_{macro} \tag{\ref{eq:conjugate_momentum}} \\
    \mathcal{H}_{\boldsymbol{\tau}} &= \mathcal{K}_{macro} + \mathcal{P} \tag{\ref{eq:legendre_transform}} \\
    \mathcal{H}_{\boldsymbol{\tau}} &= \frac{1}{2} \rho_{\tau} \Big[ \mathbf{v}_{macro}^2 + c^2 \Big] \tag{\ref{eq:bernoulli_hamiltonian}} \\
    \mathcal{H}_{dyn} &= \frac{1}{2} \rho_{\tau} \Big[ \alpha_{\parallel}(\nabla h)^2 + \alpha_{\perp}(\nabla \times \mathbf{A})^2 + \alpha_{\circlearrowleft}(\nabla \times \mathbf{v}_{\psi})^2 \Big] + \frac{1}{2} \rho_{\tau} c^2 \tag{\ref{eq:modal_hamiltonian}} \\
    \mathcal{H}_{kin} &= \frac{1}{2} \Big[ \alpha_{\parallel}(\nabla h)^2 + \alpha_{\perp}(\nabla \times \mathbf{A})^2 + \alpha_{\circlearrowleft}(\nabla \times \mathbf{v}_{\psi})^2 \Big] + \frac{1}{2} c^2 \tag{\ref{eq:kinematic_hamiltonian}} \\
    |\nabla \times \mathbf{v}_{\psi}| &\to c \implies \alpha_{\circlearrowleft} \to 0 \implies \mathcal{H}_{rest} = \frac{1}{2}\rho_{\tau}c^2 \tag{\ref{eq:hamiltonian_rest_mass}} \\
    \mathcal{H}_{kin}(v_{tot}, \theta, \phi) &= \frac{1}{2} v_{tot}^2 \Big[ \alpha_{\parallel} \cos^2(\theta) + \alpha_{\perp} \sin^2(\theta)\cos^2(\phi) \nonumber \\
    &\quad + \alpha_{\circlearrowleft} \sin^2(\theta)\sin^2(\phi) \Big] + \frac{1}{2} c^2 \tag{\ref{eq:angular_hamiltonian}} \\
    E_{rest} &= \int \mathcal{H}_{rest} \, dV = \frac{1}{2} \rho_{\tau} c^2 V_{\psi} \tag{\ref{eq:rest_energy_integral}} \\
    m_0 &= \frac{1}{2} \rho_{\tau} V_{\psi} \tag{\ref{eq:topological_mass}}
\end{align}
\end{tcolorbox}







% ----------------------------------------------------------------------
\section{Topological Vorticity Equilibrium (The Dirac Isomorphism)}
\label{sec:deriv_dirac_isomorphism}

Fermion kinematics evaluate as the continuous mechanics of a gyroscopically stabilized acoustic vortex. The abstract algebraic operators of the standard Dirac sector evaluate explicitly as classical geometric fluid identities.

\textbf{1. The Spinor (Cartan Tetrad):} \\
The state vector ($\psi$) evaluates as the orthogonal geometric orientation of the localized fluid vortex, defined by the Cartan moving frame ($\mathbf{\hat{e}}_i$):
\begin{equation}
    \psi \equiv \mathbf{\hat{e}}_i \tag{\ref{eq:spinor_identity}}
\end{equation}
The Clifford algebra of the $\gamma$-matrices evaluates as the continuous rotation generators coupling this localized internal vorticity to the macroscopic Eulerian grid. 

\textbf{2. The Kinematic Operator:} \\
The kinetic operator ($i\gamma^\mu \partial_\mu$) evaluates as the continuous kinematic curl ($\nabla \times \mathbf{v}_\psi$) of the localized $\boldsymbol{\tau}$-fluid:
\begin{equation}
    i\gamma^\mu \partial_\mu \equiv \nabla \times \mathbf{v}_\psi \tag{\ref{eq:gamma_curl}}
\end{equation}

\textbf{3. Geometric Mass Coupling (Absolute Density):} \\
To transition from a discrete topological energy to a continuous field interaction, the localized rest energy ($E_f = \frac{1}{2} \rho_{\tau} V_{\tau} c^2$) is distributed across the geometric field density ($\bar{\psi} \psi = 1 / V_{\tau}$). 

The mass Lagrangian evaluates as the negative product of the localized energy and the field distribution:
\begin{align}
    \mathcal{L}_{mass} &= - E_f (\bar{\psi} \psi) \nonumber \\
    \mathcal{L}_{mass} &= - \left( \frac{1}{2} \rho_{\tau} V_{\tau} c^2 \right) \left( \frac{1}{V_{\tau}} \right)  \\
    \mathcal{L}_{mass} &= - \frac{1}{2} \rho_{\tau} c^2 \left( \frac{V_{\tau}}{V_{\tau}} \right) 
\end{align}
The specific topological volume ($V_{\tau}$) completely cancels from the architecture. The interaction density isolates to an absolute continuum constant:
\begin{equation}
    \mathcal{L}_{mass} = - \frac{1}{2} \rho_{\tau} c^2 \tag{\ref{eq:absolute_mass_lagrangian}}
\end{equation}
The mass interaction term ($m\bar{\psi}\psi$) evaluates universally as one-half the absolute resting energy density of the macroscopic vacuum, eliminating the requirement for parameterized coupling constants.

\textbf{4. The Vorticity Equilibrium Lagrangian:} \\
The fundamental interaction Lagrangian evaluates the kinetic and potential energy densities of the topological state:
\begin{equation}
    \mathcal{L}_{Dirac} = \bar{\psi} (i\gamma^\mu \partial_\mu) \psi - m\bar{\psi}\psi 
\end{equation}
Substituting the continuous fluid vorticity and absolute mass density identities maps the quantum operator   to the macroscopic continuum state equation:
\begin{equation}
    \mathcal{L}_{Dirac} \equiv \frac{1}{2}\rho_{\tau} |\nabla \times \mathbf{v}_\psi|^2 - \frac{1}{2} \rho_{\tau} c^2 \bar{\psi}\psi \tag{\ref{eq:dirac_fluid_lagrangian}}
\end{equation}
Setting the Lagrangian to zero mathematically dictates that the dynamic kinetic pressure of the internal vortex ($\frac{1}{2}\rho_{\tau} |\nabla \times \mathbf{v}_\psi|^2$) must perfectly balance the absolute macroscopic static pressure required to sustain the zero-pressure acoustic cavity. The localized fluid vortex is structurally mandated to circulate at the exact kinematic phase limit ($c$).

\textbf{5. Antimatter (Chiral Inversion \& Annihilation):} \\
The adjoint ($\bar{\psi}$) evaluates as the topological chiral inverse of the fluid circulation. A geometric boundary and its conjugate evaluate as $W_+ = +1$ and $W_- = -1$ fluid circulations. Their spatial superposition dictates the scalar interference:
\begin{align}
    \bar{\psi}\psi &\equiv (W_+ + W_-) \nonumber \\
    \bar{\psi}\psi &= 0 \tag{\ref{eq:chiral_annihilate}}
\end{align}
This zero-sum geometric interference forces localized circulation to zero ($\mathbf{v}_{tot} = 0$). By Bernoulli's principle, the collapse of internal dynamic pressure instantly restores the absolute macroscopic static pressure of the continuum. The vacuum violently crushes the zero-pressure acoustic cavity shut, causing deterministic structural annihilation.

\begin{tcolorbox}[recallbox, title={\ref{box:dirac_isomorphism}. Topological Vorticity Equilibrium (The Dirac Equation)}]
\begin{align}
    \psi &\equiv \mathbf{\hat{e}}_i \tag{\ref{eq:spinor_identity}} \\
    i\gamma^\mu \partial_\mu &\equiv \nabla \times \mathbf{v}_\psi \tag{\ref{eq:gamma_curl}} \\
    \bar{\psi} \psi = \frac{1}{V_{\tau}} &\implies \mathcal{L}_{mass} = - \frac{1}{2} \rho_{\tau} c^2 \tag{\ref{eq:absolute_mass_lagrangian}} \\
    \mathcal{L}_{Dirac} &\equiv \frac{1}{2}\rho_{\tau} |\nabla \times \mathbf{v}_\psi|^2 - \frac{1}{2} \rho_{\tau} c^2 \bar{\psi}\psi \tag{\ref{eq:dirac_fluid_lagrangian}} \\
    \bar{\psi}\psi &\equiv (W_+ + W_-) = 0 \tag{\ref{eq:chiral_annihilate}} 
\end{align}
\end{tcolorbox}


% ----------------------------------------------------------------------
\section{Continuous Wave Mechanics \& Topological Phase States}
\label{sec:deriv_wave_mechanics}

The Kinematic Lagrangian maps the mechanical wave spectrum, discrete particle identities, and nuclear decay as purely geometric constraints of the continuous vacuum.

\textbf{1. The Continuous Wave Solutions:} \\
The Helmholtz decomposition partitions the continuous motion of the metric into a complete, real-valued mechanical wave spectrum. The dynamic propagation evaluates as the absolute kinematic limit ($c$) scaled by the geometric structural phase angles ($\theta, \phi$) and the temporal propagation phase ($\mathbf{k} \cdot \mathbf{r} - \omega t$).

The three orthogonal fluid wave equations evaluate without imaginary phase states:
\begin{align}
    \mathbf{v}_{\parallel}(r,t) &= c \cos(\theta) \cos(\mathbf{k} \cdot \mathbf{r} - \omega t) \tag{\ref{eq:wave_p}} \\
    \mathbf{v}_{\perp}(r,t) &= c \sin(\theta)\cos(\phi) \cos(\mathbf{k} \cdot \mathbf{r} - \omega t) \tag{\ref{eq:wave_s}} \\
    \mathbf{v}_{\circlearrowleft}(r,t) &= c \sin(\theta)\sin(\phi) \cos(\mathbf{k} \cdot \mathbf{r} - \omega t) \tag{\ref{eq:wave_rot}}
\end{align}

These velocity fields evaluate the complete physical spectrum:
\begin{itemize}
    \item \textbf{Longitudinal ($\mathbf{v}_{\parallel}$):} Evaluates as propagating P-waves. A longitudinal P-wave obeys the D'Alembert wave equation ($\nabla^2 h - \frac{1}{c^2} \frac{\partial^2 h}{\partial t^2} = 0$). Because the longitudinal field natively governs macroscopic metric fluid compression, substituting the ambient kinematic velocity required to balance the dynamic pressure deficit ($v_{\parallel} = \sqrt{2GM/r}$) perfectly evaluates the classical gravitational potential:
    \begin{equation}
        \Phi(r,t) = - \frac{(\mathbf{v}_{\parallel})^2}{2} \label{eq:p_wave_potential}
    \end{equation}
    Dynamic fluctuations in this compressive field evaluate mechanically as propagating Acoustic Gravitational Radiation.
    \item \textbf{Transverse ($\mathbf{v}_{\perp}$):} Evaluates as propagating S-waves (Electromagnetic Radiation).
    \item \textbf{Rotational ($\mathbf{v}_{\circlearrowleft}$):} Evaluates as localized acoustic vortices (Topological Rest Mass).
\end{itemize}

\textbf{2. Kinematics \& The Lorentz Angle ($\theta$):} \\
The Divergence Angle ($\theta$) partitions the kinetic capacity of the fluid between macroscopic longitudinal translation ($\nabla h$) and localized topological structure. For a saturated topological boundary ($v_{tot} \to c$), the macroscopic translation velocity ($v_{macro} \equiv \mathbf{v}_{\parallel}$) evaluates as the longitudinal cosine projection of the absolute limit $c$:
\begin{align}
    v_{macro} &= c \cos(\theta) \nonumber \\
    v_{macro}^2 &= c^2 \cos^2(\theta) \nonumber \\
    \frac{v_{macro}^2}{c^2} &= \cos^2(\theta) 
\end{align}
Applying the Pythagorean identity ($\sin^2\theta = 1 - \cos^2\theta$) isolates the sine projection:
\begin{align}
    \sin(\theta) &= \sqrt{1 - \cos^2(\theta)} \nonumber \\
    \sin(\theta) &= \sqrt{1 - \frac{v_{macro}^2}{c^2}} \tag{\ref{eq:lorentz_angle}}
\end{align}
The term $\sin(\theta)$ evaluates identically to the inverse Lorentz factor ($1/\gamma$) and the longitudinal admittance throttle ($\alpha_{\parallel}$). Relativistic kinematics evaluate as a direct algebraic record of the localized boundary closing its Divergence Angle to accommodate macroscopic velocity through the fluid.

\textbf{3. Topological Charge \& The Phase Map ($\phi$):} \\
For a continuous boundary evaluated at the rigid transverse equator ($\theta = \pi/2$), identity is governed exclusively by the Curl Angle ($\phi$), which partitions energy between rotational mass and transverse electromagnetism. Electric charge ($|Q|$) evaluates as the geometric fraction of kinematic energy dedicated to transverse shear:
\begin{equation}
    |Q| = \cos^2(\phi) \tag{\ref{eq:charge_projection}}
\end{equation}

Because the continuum evaluates as a cohesive elastic volume, stable internal topologies lock into geometric harmonic fractions. Table \ref{tab:wave_phases} maps the fundamental and composite states as angular fluid constraints.

\begin{table}[h]
\centering
\begin{tabular}{l c c c l}
\hline
\textbf{Particle State} & \textbf{Divergence ($\theta$)} & \textbf{Curl ($\phi$)} & \textbf{Charge ($|Q|$)} & \textbf{Mechanical State} \\
\hline
\multicolumn{5}{c}{\textbf{Closed Topological Boundaries (Fundamental Fermions)}} \\
\hline
Electron ($e^-$) & Dipole & $\pi$ & 1 & Transverse Vortex \\
Positron ($e^+$) & Dipole & $-\pi$ & 1 & Transverse Anti-Vortex \\
Up Quark ($u$) & Dipole & $\arcsin(\sqrt{1/3})$ & 2/3 & Primary Triad Lobe \\
Down Quark ($d$) & Dipole & $\arcsin(\sqrt{2/3})$ & 1/3 & Recessive Triad Lobe \\
Top Quark ($t$) & Dipole & $\arcsin(\sqrt{1/3})$ & 2/3 & Extreme Triad Lobe \\
\hline
\multicolumn{5}{c}{\textbf{Composite Topological Boundaries (Baryons)}} \\
\hline
Proton ($p^+$) & Dipole & Trefoil ($uud$) & 1 & Stable Composite Knot \\
Neutron ($n^0$) & Dipole & Trefoil ($udd$) & 0 & Unstable Composite Knot \\
\hline
\multicolumn{5}{c}{\textbf{Open Propagating Waves (Bosons \& Uncontained Strain)}} \\
\hline
Photon ($\gamma$) & $\pi/2$ & $0$ & $1$ (Carrier) & Transverse S-Wave \\
Neutrino ($\nu$) & $\to \pi$ & $\pi/2$ & 0 & Left-Handed Slip-Knot \\
Antineutrino ($\bar{\nu}$) & $\to 0$ & $\pi/2$ & 0 & Right-Handed Slip-Knot \\
\hline
\end{tabular}
\caption{Phase Angle Mapping of Fundamental and Composite Fluid States}
\label{tab:wave_phases}
\end{table}

\textbf{4. The Dynamic Neutrino Limit:} \\
A neutrino evaluates as uncontained longitudinal strain ($\phi = \pi/2$). Its cosine projection is zero, yielding zero electromagnetic charge. It lacks a rigid resting volume and travels at the macroscopic kinematic limit ($v_{macro} \to c$). 

Substituting $v \to c$ into Equation \ref{eq:lorentz_angle} forces the Divergence Angle to the longitudinal axis ($\theta \to \pi$ for left-handed, $\theta \to 0$ for right-handed). Substituting these geometric limits into the Kinematic Lagrangian ($\mathcal{L}_{kin}$) eliminates the localized phase-space action:
\begin{equation}
    \mathcal{L}_{kin}(\pi, \pi/2) \to 0 \quad \text{and} \quad \mathcal{L}_{kin}(0, \pi/2) \to 0 \tag{\ref{eq:phase_neutrino}}
\end{equation}
The neutrino evaluates as a longitudinal wave-packet propagating through the continuum without the geometric transverse shear necessary to interact with electromagnetic cross-sections.

\textbf{5. Antimatter \& Annihilation Rupture:} \\
Antimatter evaluates as the geometric inverse of a localized vortex. Where a matter vortex evaluates as $\nabla \times \mathbf{v}_{\psi}^+ = +c$, an anti-vortex evaluates as $\nabla \times \mathbf{v}_{\psi}^- = -c$. Because the equations evaluate the square of the internal velocities, the phase inversion mathematically neutralizes:
\begin{equation}
    (\nabla \times \mathbf{v}_{\psi}^-)^2 = (-c)^2 = c^2 \tag{\ref{eq:antimatter_state}}
\end{equation}
Antimatter evaluates identical positive mass and metric curvature. 

Annihilation evaluates as localized destructive interference. When a vortex ($+c$) and an anti-vortex ($-c$) overlap, their fluid boundaries mechanically sum:
\begin{align}
    \sum \mathbf{v}_{\psi} &= (+c) + (-c) \nonumber \\
    \sum \mathbf{v}_{\psi} &= 0 \nonumber
\end{align}
Substituting this zero-velocity state into the rotational admittance equation calculates the mechanical destruction of the topological boundary:
\begin{align}
    \alpha_{\circlearrowleft} &= \sqrt{1 - \frac{(\sum \mathbf{v}_{\psi})^2}{c^2}} \nonumber \\
    \alpha_{\circlearrowleft} &= \sqrt{1 - \frac{0^2}{c^2}} \nonumber \\
    \alpha_{\circlearrowleft} &= 1 \tag{\ref{eq:annihilation_limit}}
\end{align}
The infinite Deborah number collapses. The absolute macroscopic static pressure of the continuum violently crushes the zero-pressure cavities shut. The kinetic energy of this fluid collapse rebounds outward into the continuum as compliant transverse electromagnetic shear (gamma radiation).

\textbf{6. Mechanical Yield \& Beta Decay ($\beta^-$):} \\
Nuclear decay evaluates as a localized mechanical fluid yield. A neutron evaluates as an isolated triad ($u, d, d$). Outside the compressive pressure of a nucleus, the triad loses equilibrium. The heavy down lobe mechanically buckles, structurally relaxing into the lower-volume phase of an up quark ($u$). 

The excess transverse phase ($\Delta Q = -1$) and volumetric kinetic energy sheared during this phase transition mechanically extrude from the triad, stabilizing into an independent transverse vortex (the electron) and a purely kinetic longitudinal wake (the antineutrino):
\begin{equation}
    \mathcal{L}_{kin, n}(udd) \to \mathcal{L}_{kin, p}(uud) + \mathcal{L}_{kin, e}(\phi=\pi) + \mathcal{L}_{kin, \bar{\nu}}(\theta \to 0) \tag{\ref{eq:neutron_rupture}}
\end{equation}
Standard Model decay channels evaluate as absolute algebraic records of continuous angular fluid volume seeking structural geometric equilibrium.

\begin{tcolorbox}[recallbox, title={\ref{box:wave_mechanics}. Continuous Wave Mechanics \& Topological Phase States}]
\begin{align}
    \mathbf{v}_{\parallel}(r,t) &= c \cos(\theta) \cos(\mathbf{k} \cdot \mathbf{r} - \omega t) \tag{\ref{eq:wave_p}} \\
    \mathbf{v}_{\perp}(r,t) &= c \sin(\theta)\cos(\phi) \cos(\mathbf{k} \cdot \mathbf{r} - \omega t) \tag{\ref{eq:wave_s}} \\
    \mathbf{v}_{\circlearrowleft}(r,t) &= c \sin(\theta)\sin(\phi) \cos(\mathbf{k} \cdot \mathbf{r} - \omega t) \tag{\ref{eq:wave_rot}} \\
    \Phi(r,t) &= - \frac{(\mathbf{v}_{\parallel})^2}{2} \tag{\ref{eq:p_wave_potential}} \\
    \sin(\theta) &= \sqrt{1 - \frac{v_{macro}^2}{c^2}} \tag{\ref{eq:lorentz_angle}} \\
    |Q| &= \cos^2(\phi) \tag{\ref{eq:charge_projection}} \\
    \theta \to \pi, 0 \implies \mathcal{L}_{kin} &\to 0 \tag{\ref{eq:phase_neutrino}} \\
    (\nabla \times \mathbf{v}_{\psi}^-)^2 &= c^2 \tag{\ref{eq:antimatter_state}} \\
    \sum \mathbf{v}_{\psi} = 0 &\implies \alpha_{\circlearrowleft} = 1 \tag{\ref{eq:annihilation_limit}} \\
    \mathcal{L}_{kin, n}(udd) &\to \mathcal{L}_{kin, p}(uud) + \mathcal{L}_{kin, e}(\phi=\pi) + \mathcal{L}_{kin, \bar{\nu}}(\theta \to 0) \tag{\ref{eq:neutron_rupture}}
\end{align}
\end{tcolorbox}












% ----------------------------------------------------------------------
\section{EM Kinematics \& The LC-Acoustic Isomorphism}
\label{sec:deriv_em_lc_isomorphism}

The phase velocity of a transverse shear wave propagating through the metric evaluates identically via the Newton-Laplace acoustic equation ($c^2 = K/\rho_{\tau}$) and the Maxwell electromagnetic limit ($c^2 = 1/\epsilon_{\tau}\mu_{\tau}$). 

\textbf{1. Magnetic Vorticity \& Electric Strain:} \\
Let the continuous transverse kinematic displacement field of the metric evaluate as $\mathbf{A}_{\perp}$. The magnetic field ($\mathbf{B}$) evaluates as the spatial vorticity of the continuous metric ($\mathbf{B} = \nabla \times \mathbf{A}_{\perp}$). The divergence of a curl evaluates to zero:
\begin{equation}
    \nabla \cdot \mathbf{B} = \nabla \cdot (\nabla \times \mathbf{A}_{\perp}) = 0 \tag{\ref{eq:maxwell_gauss_faraday}a}
\end{equation}
The electric field ($\mathbf{E}$) evaluates as the transverse strain velocity ($\mathbf{E} = -\partial \mathbf{A}_{\perp}/\partial t$). Evaluating the curl of the electric field and commuting the independent spatial and temporal derivative operators isolates the induction state:
\begin{align}
    \nabla \times \mathbf{E} &= \nabla \times \left(-\frac{\partial \mathbf{A}_{\perp}}{\partial t}\right) \nonumber \\
    \nabla \times \mathbf{E} &= -\frac{\partial}{\partial t}(\nabla \times \mathbf{A}_{\perp}) = -\frac{\partial \mathbf{B}}{\partial t} \tag{\ref{eq:maxwell_gauss_faraday}b}
\end{align}

\textbf{2. Volumetric Strain Displacement (Gauss's Law for Electricity):} \\
Absolute metric charge ($e_{\tau}$) evaluates as the spatial volume displaced by transverse strain. The localized charge density ($\rho_e$) evaluates as this displacement per unit volume:
\begin{equation}
    e_{\tau} = \int_V \rho_e \, dV
\end{equation}
Equating this to the geometric surface integral of the electric strain ($\mathbf{E}$), scaled by the metric stiffness ($\epsilon_{\tau}$):
\begin{equation}
    \int_V \rho_e \, dV = \int_S \epsilon_{\tau} \mathbf{E} \cdot d\mathbf{A}
\end{equation}
Applying the Divergence Theorem ($\int_S \mathbf{E} \cdot d\mathbf{A} = \int_V (\nabla \cdot \mathbf{E}) \, dV$) translates the surface integral into a continuous spatial divergence. Equating the integrands isolates the continuous volumetric expansion of the transverse field:
\begin{equation}
    \nabla \cdot \mathbf{E} = \frac{\rho_e}{\epsilon_{\tau}} \tag{\ref{eq:maxwell_gaussE}}
\end{equation}

\textbf{3. The 1D vs. 3D Dimensional Mapping:} \\
Under the Electro-Mechanical Mobility mapping, the geometry underlying 1D metric permittivity evaluates as the absolute 1D Metric Stiffness ($\epsilon_{\tau}$) of a continuous topological flux tube. Metric permeability evaluates as the absolute 1D Metric Compliance ($\mu_{\tau}$) of that topological tether. 

To transition from the 1D metric tether to the 3D macroscopic continuum ($\rho_{\tau}$), the 1D linear density divides by the fundamental transverse geometric yield area ($A_{yield}$) of the continuum strain. Substituting the compliance limit ($1/\mu_{\tau} = T_{1D}$) yields the absolute 3D elastic modulus ($K$):
\begin{align}
    \rho_{\tau} &= \frac{\epsilon_{\tau}}{A_{yield}} \nonumber \\
    K &= \frac{1}{\mu_{\tau} A_{yield}} = \frac{T_{1D}}{A_{yield}} \tag{\ref{eq:lc_constants}}
\end{align}

\textbf{4. The Isomorphic Circuit Elements:} \\
\textit{Capacitance ($C \equiv m_\tau$):} Observable rest mass evaluates as 3D hydrodynamic added mass ($m_\tau = \frac{1}{2} \rho_{\tau} V_\tau$). Because 3D volume divided by 2D area reduces to 1D length ($V_\tau/A_{yield} = L_\tau$), observable inertia evaluates as localized 1D capacitance:
\begin{equation}
    m_\tau = \frac{1}{2} \epsilon_{\tau} L_\tau \tag{\ref{eq:lc_mass}}
\end{equation}

\textit{Transverse Velocity ($v_{\perp}$):} Evaluating the energy of a localized capacitor ($E_{cap}$) and substituting the isomorphic metric ($C \to m_\tau$) maps directly to the classical kinetic energy of a macroscopic mass:
\begin{equation}
    E_{cap} = E_{kin} = \frac{1}{2} m_\tau v_{\perp}^2
\end{equation}
Electrical potential evaluates fundamentally as the transverse kinematic velocity ($v_{\perp}$) of the metric. Substituting this identity into the continuous Bernoulli equation dictates that this localized kinetic energy generates the macroscopic dynamic pressure of the vacuum metric:
\begin{equation}
    P_{dyn} = \frac{1}{2}\rho_{\tau} v_{\perp}^2 \tag{\ref{eq:velocity_pressure}}
\end{equation}

\textit{Inductance ($L_H \equiv 1/k$):} Electrical current evaluates longitudinally as Macroscopic Mechanical Force ($I \to F$), and compliance evaluates as the structural Mechanical Compliance ($L_H \to 1/k$). Because Hookean elastic force dictates $F = kx$, dynamic inductive energy evaluates as static elastic potential energy:
\begin{equation}
    E_{ind} = \frac{1}{2} L_H I^2 = \frac{1}{2} \left(\frac{1}{k}\right) (kx)^2 = \frac{1}{2} k x^2 \tag{\ref{eq:lc_inductance}}
\end{equation}

\textbf{5. The 1D Electromagnetic Projection (Ampère-Maxwell):} \\
Applying the continuous 3D Navier-Cauchy equation for transverse waves yields:
\begin{equation}
    \rho_{\tau} \frac{\partial^2 \mathbf{A}_{\perp}}{\partial t^2} + \mu_s \nabla \times (\nabla \times \mathbf{A}_{\perp}) = \mathbf{J}_{3D}
\end{equation}
Substituting the kinematic fields ($\frac{\partial \mathbf{A}_{\perp}}{\partial t} = -\mathbf{E}$ and $\nabla \times \mathbf{A}_{\perp} = \mathbf{B}$) and the dimensional 1D/3D geometric projections ($\rho_{\tau} = \epsilon_{\tau}/A_{yield}$ and $\mu_s = 1/\mu_{\tau} A_{yield}$):
\begin{align}
    \left( \frac{1}{\mu_{\tau} A_{yield}} \right) \nabla \times \mathbf{B} &= \mathbf{J}_{3D} + \left( \frac{\epsilon_{\tau}}{A_{yield}} \right) \frac{\partial \mathbf{E}}{\partial t} \nonumber \\
    \nabla \times \mathbf{B} &= \mu_{\tau} (A_{yield} \mathbf{J}_{3D}) + \mu_{\tau} \epsilon_{\tau} \frac{\partial \mathbf{E}}{\partial t}
\end{align}
Because electromagnetic current evaluates as the 1D topological projection of the 3D continuum, the current density ($\mathbf{J}_{EM}$) evaluates as the 3D macroscopic current ($\mathbf{J}_{3D}$) projected across the transverse yield area ($\mathbf{J}_{EM} = A_{yield} \mathbf{J}_{3D}$). Substituting this localized continuous current isolates the relation:
\begin{equation}
    \nabla \times \mathbf{B} = \mu_{\tau} \mathbf{J}_{EM} + \mu_{\tau} \epsilon_{\tau} \frac{\partial \mathbf{E}}{\partial t} \tag{\ref{eq:maxwell_curlB}}
\end{equation}

\textbf{6. Kinematic Energy Flux (The Poynting Vector):} \\
The directional transport of electromagnetic energy evaluates geometrically as the macroscopic propagation of continuous transverse kinetic energy through the fluid metric. This propagation evaluates as the Poynting vector ($\mathbf{S}$):
\begin{equation}
    \mathbf{S} = \frac{1}{\mu_{\tau}} (\mathbf{E} \times \mathbf{B}) \tag{\ref{eq:poynting_flux}}
\end{equation}
Because the electric strain velocity ($\mathbf{E}$) and magnetic vorticity ($\mathbf{B}$) are orthogonal transverse fluid geometries, their cross product generates a purely longitudinal propagation vector. 

This evaluates the continuous temporal flux of kinetic energy, independent of bulk mass advection. The transport of kinematic energy evaluates exclusively within the surrounding metric fluid, mathematically orthogonal to the topological boundary anchoring the strain. Macroscopic energy flux does not flow inside the localized topological boundary; it propagates longitudinally through the external macroscopic vacuum metric.

\textbf{7. The Dual Transverse Geometries (Linear vs. Rotational):} \\
Transverse metric kinematics bifurcate into two orthogonal spatial deformations within the 2D plane perpendicular to longitudinal advection. Electrical potential (Voltage) evaluates as the pure translational sliding of the metric, defining the linear transverse shear velocity ($\mathcal{V} \equiv v_{\perp}$). Magnetism evaluates as the localized twisting of the metric, defining the rotational transverse vorticity ($\mathbf{B} \equiv \nabla \times \mathbf{A}_{\perp}$). 

Because these distinct geometries consume the localized phase limit ($c$) independently, their simultaneous presence forces an orthogonal response. The cross-product dynamically squeezes the continuum, generating the longitudinal kinematic flux ($\mathbf{S}$).
\begin{equation}
    \mathcal{V} \equiv v_{\perp} \text{ (Linear)} \quad \text{\&} \quad \mathbf{B} \equiv \nabla \times \mathbf{A}_{\perp} \text{ (Rotational)} \label{eq:dual_transverse}
\end{equation}

\textbf{8. The Universal Voltage Constant \& Absolute Phase Limit:} \\
To map human-defined electrical units to absolute continuum kinematics, invariant mechanical energy ($E_{cap} = E_{kin}$) dictates the dimensional scaling. The conversion factor ($Z$) mapping legacy Capacitance (Farads) to absolute 1D fluid Mass (kg) isolates as the ratio of true metric stiffness ($\epsilon_{\tau}$) to legacy vacuum permittivity ($\epsilon_0$):
\begin{equation}
    Z = \frac{\epsilon_{\tau}}{\epsilon_0} \approx 2.597 \times 10^{71} \text{ kg/F}
\end{equation}
Equating the capacitive and kinematic energies ($\frac{1}{2} C \mathcal{V}^2 = \frac{1}{2} m v_{\perp}^2$) and substituting the mass mapping ($m = C \cdot Z$) isolates the Universal Voltage Constant:
\begin{equation}
    \mathcal{V} = \sqrt{Z} \cdot v_{\perp} \approx (5.096 \times 10^{35}) \cdot v_{\perp} \label{eq:universal_voltage_a}
\end{equation}
The absolute maximum continuous voltage of the metric evaluates when the linear transverse velocity strikes the phase bottleneck ($v_{\perp} = c$):
\begin{equation}
    \mathcal{V}_{max} = \sqrt{Z} \cdot c \approx 1.527 \times 10^{44} \text{ Volts} \label{eq:universal_voltage_b}
\end{equation}
This absolute kinematic ceiling bounds the continuous static domain. Human-scale high-voltage phenomena ($\approx 10^9 \text{ V}$) evaluate to infinitesimally small transverse fluid velocities ($v_{\perp} \approx 10^{-27} \text{ m/s}$), ensuring the localized vacuum metric operates deep within the linear, un-strain-hardened Newtonian regime.

\textbf{9. The Macroscopic Vorticity Tensor (Faraday Matrix):} \\
The classical electromagnetic Faraday tensor ($F^{\mu\nu}$) evaluates structurally as the macroscopic antisymmetric velocity gradient tensor ($\Omega^{\mu\nu}$) of the continuous $\boldsymbol{\tau}$-fluid. By mapping the four-potential to the continuous fluid velocity ($A^\mu = u^\mu$), the unified tensor cleanly separates the dual transverse geometries into distinct dimensional planes based on their orthogonal shear kinematics ($\partial^\mu u^\nu - \partial^\nu u^\mu$):
\begin{equation}
    F^{\mu\nu} \equiv \Omega^{\mu\nu} = \begin{bmatrix} 
    0 & -E_x/c & -E_y/c & -E_z/c \\ 
    E_x/c & 0 & -B_z & B_y \\ 
    E_y/c & B_z & 0 & -B_x \\ 
    E_z/c & -B_y & B_x & 0 
    \end{bmatrix} \tag{\ref{eq:rosetta_faraday}}
\end{equation}

Because the kinematic deformations are transverse (volume-conserving), all normal stress components vanish, driving the principal diagonal to zero. The remaining off-diagonal elements isolate the two distinct fluid behaviors:

\textit{Space-Time Shear (Linear / Voltage):} The intersection of the temporal and spatial dimensions isolates the temporal rate of spatial sliding ($\partial_t u_x$). This rigidly maps the linear transverse velocity ($\mathbf{E}$) to the space-time off-diagonals. 

\textit{Space-Space Shear (Rotational / Magnetism):} The intersection of purely spatial dimensions isolates the classical 2D fluid curl ($\partial_x u_y - \partial_y u_x$). This rigidly maps the localized rotational vorticity ($\mathbf{B}$) to the space-space off-diagonals. 

The macroscopic electromagnetic field evaluates as the unified off-diagonal kinematic shear of the baseline vacuum metric.

\begin{tcolorbox}[recallbox, title={\ref{box:em_lc_isomorphism}. EM Kinematics \& The LC-Acoustic Isomorphism}]
\begin{align}
    \nabla \cdot \mathbf{B} = 0 \quad &\text{\&} \quad \nabla \times \mathbf{E} = -\frac{\partial \mathbf{B}}{\partial t} \tag{\ref{eq:maxwell_gauss_faraday}} \\
    \nabla \cdot \mathbf{E} &= \frac{\rho_e}{\epsilon_{\tau}} \tag{\ref{eq:maxwell_gaussE}} \\
    \rho_{\tau} = \frac{\epsilon_{\tau}}{A_{yield}}, \quad &K = \frac{1}{\mu_{\tau} A_{yield}} \tag{\ref{eq:lc_constants}} \\
    m_\tau &= \frac{1}{2} \epsilon_{\tau} L_\tau \tag{\ref{eq:lc_mass}} \\
    E_{cap} = E_{kin} &= \frac{1}{2} m_{\tau} v_{\perp}^2 \implies P_{dyn} = \frac{1}{2}\rho_{\tau} v_{\perp}^2 \tag{\ref{eq:velocity_pressure}} \\
    E_{ind} = \frac{1}{2} L_H I^2 &= \frac{1}{2} k x^2 \tag{\ref{eq:lc_inductance}} \\
    \nabla \times \mathbf{B} &= \mu_{\tau} \mathbf{J}_{EM} + \mu_{\tau} \epsilon_{\tau} \frac{\partial \mathbf{E}}{\partial t} \tag{\ref{eq:maxwell_curlB}} \\
    \mathbf{S} &= \frac{1}{\mu_{\tau}} (\mathbf{E} \times \mathbf{B}) \tag{\ref{eq:poynting_flux}} \\
    \mathcal{V} \equiv v_{\perp} \text{ (Linear)} \quad &\text{\&} \quad \mathbf{B} \equiv \nabla \times \mathbf{A}_{\perp} \text{ (Rotational)} \tag{\ref{eq:dual_transverse}} \\
    \mathcal{V}_{max} &= \sqrt{Z} \cdot c \approx 1.52 \times 10^{44} \text{ V} \tag{\ref{eq:universal_voltage_b}}
\end{align}
\end{tcolorbox}





% ----------------------------------------------------------------------
\section{Absolute Continuum Constants (UV Regularization \& Stiffness)}
\label{sec:deriv_continuum_constants}

The integration of momentum states evaluates physically to a structural maximum bounded by the macroscopic yield strength of the continuous metric fluid.

\textbf{1. Absolute 3D Volumetric Density ($\rho_{\tau}$):} \\
The 3D macroscopic mass density of the vacuum isolates as a spherical geometric fraction ($9/4\pi$) of the absolute absolute baseline limit ($\rho_{pl} = c^5 / \hbar G^2$). Evaluating this limit physically quantifies the absolute continuous fluid density:
\begin{align}
    \rho_{\tau} &= \frac{9}{4\pi} \frac{c^5}{\hbar G^2} \label{eq:rho_tau_analytical} \\
    \rho_{\tau} &\approx 3.69 \times 10^{96} \text{ kg/m}^3 \tag{\ref{eq:rho_tau_numerical}}
\end{align}

\textbf{2. Absolute Energy Ceiling (The Cavitation Limit):} \\
By the macroscopic elastic equation of state governed by the Bulk Modulus ($K = \frac{1}{3}\rho_{\tau} c^2$), the absolute isotropic hydrostatic tension of the 3D fluid vacuum evaluates as $\sigma_{max} = \frac{1}{3}\rho_{\tau} c^2$.

For a localized geometric wave packet of kinematic energy $E$, the localized dynamic pressure exerted by this circulating geometric strain evaluates as the energy density of the spherical cavity:
\begin{equation}
    P_{dyn} = \frac{3E^4}{4\pi \hbar^3 c^3} \label{eq:dynamic_pressure}
\end{equation}

If the dynamic pressure equals the structural yield tension, the fluid is deterministically forced to cavitate ($P_{dyn} = \sigma_{max}$). Substituting the absolute density ($\rho_{\tau}$) isolates the maximum energy ceiling of the continuum:
\begin{align}
    \frac{3E_{max}^4}{4\pi \hbar^3 c^3} &= \frac{1}{3} \left( \frac{9}{4\pi} \frac{c^5}{\hbar G^2} \right) c^2 \label{eq:cavitation_equivalence} \\
    E_{max}^4 &= \frac{4\pi \hbar^3 c^3}{3} \left( \frac{3 c^7}{4\pi \hbar G^2} \right) \label{eq:energy_isolation_a} \\
    E_{max}^4 &= \frac{\hbar^2 c^{10}}{G^2} \label{eq:energy_isolation_b}
\end{align}
Taking the fourth root mathematically isolates the absolute maximum energy ceiling of the continuum as the exact geometric evaluation of the Planck Energy ($E_p$):
\begin{align}
    E_{max} &= \left( \frac{\hbar^2 c^{10}}{G^2} \right)^{1/4} \label{eq:energy_fourth_root} \\
    E_{max} &= \sqrt{\frac{\hbar c^5}{G}} \tag{\ref{eq:planck_ceiling}}
\end{align}
The dynamic pressure ruptures the metric at this precise geometric limit, generating the void topology.

\textbf{3. Absolute 1D Metric Stiffness ($\epsilon_{\tau}$):} \\
The absolute 1D acoustic stiffness of the untensioned vacuum ($\epsilon_{\tau}$) evaluates deterministically by equating the gravitational weak coupling with the strong confinement limit of the continuous topology. 

Equating the derived continuum Fermi constant ($G_F$) to the proton boundary limit establishes the identity:
\begin{equation}
    \frac{2\sqrt{2}\pi}{9} \frac{\epsilon_{\tau} G^2 \hbar}{c^5} = \frac{81\sqrt{2}\pi\alpha^3}{14} \frac{\hbar^2}{m_{pr}^2 c^2} \label{eq:fermi_proton_identity}
\end{equation}
Algebraically isolating $\epsilon_{\tau}$ and factoring out the absolute macroscopic mass ($m_{pl} = \sqrt{\hbar c / G}$) exposes the fundamental structural hierarchy:
\begin{equation}
    \epsilon_{\tau} = \frac{729 \alpha^3}{28} \left( \frac{m_{pl}}{m_{pr}} \right)^2 \frac{c^2}{G} \label{eq:epsilon_analytical}
\end{equation}
The baseline stiffness evaluates as the continuous linear limit ($c^2/G$) scaled deterministically by the squared mass ratio of the macroscopic horizon to the microscopic core. Evaluating this limit isolates the physical 1D stiffness:
\begin{equation}
    \epsilon_{\tau} \approx 2.30 \times 10^{60} \text{ kg/m} \tag{\ref{eq:epsilon_numerical}}
\end{equation}

\textbf{4. The Gravitational Yield Area ($A_{yield}$):} \\
By the continuous dimensional geometry derived via the LC-Acoustic Isomorphism, the 3D density evaluates as the 1D stiffness divided by the transverse yield area ($\rho_{\tau} = \epsilon_{\tau} / A_{yield}$). Rearranging this equivalence algebraically isolates the baseline spatial area of the metric strain:
\begin{align}
    A_{yield} &= \frac{\epsilon_{\tau}}{\rho_{\tau}} \label{eq:area_isolation} \\
    A_{yield} &= \epsilon_{\tau} \left( \frac{4\pi G^2 \hbar}{9 c^5} \right) \label{eq:area_substitution} \\
    A_{yield} &= \frac{4\pi}{9} \frac{\epsilon_{\tau} G^2 \hbar}{c^5} \tag{\ref{eq:gravitational_yield}}
\end{align}

\begin{tcolorbox}[recallbox, title={\ref{box:continuum_constants}. Absolute Continuum Constants}]
\begin{align}
    \rho_{\tau} &= \frac{9}{4\pi} \frac{c^5}{\hbar G^2} \approx 3.69 \times 10^{96} \text{ kg/m}^3 \tag{\ref{eq:rho_tau_numerical}} \\
    E_{max} &= \sqrt{\frac{\hbar c^5}{G}} \tag{\ref{eq:planck_ceiling}} \\
    \epsilon_{\tau} &= \frac{729 \alpha^3}{28} \left( \frac{m_{pl}}{m_{pr}} \right)^2 \frac{c^2}{G} \approx 2.30 \times 10^{60} \text{ kg/m} \tag{\ref{eq:epsilon_numerical}} \\
    A_{yield} &= \frac{4\pi}{9} \frac{\epsilon_{\tau} G^2 \hbar}{c^5} \tag{\ref{eq:gravitational_yield}}
\end{align}
\end{tcolorbox}

% ----------------------------------------------------------------------
\section{The Topological Coupling Limit ($\alpha_{\circlearrowleft \perp}$) \& Charge}
\label{sec:deriv_fine_structure}

The continuous metric of the vacuum operates as a viscoelastic continuum governed by a geometric compliance limit.  In macroscopic mechanics, dimensionless transverse compliance evaluates as the ratio of longitudinal bulk modulus ($K$) to transverse shear modulus ($G$). For the vacuum metric, this fundamental rotation-to-shear yield ratio ($\alpha_{\circlearrowleft \perp} = K_\tau / G_\tau$) evaluates phenomenologically as the fine-structure constant ($\alpha \approx 1/137$). Because the metric's transverse shear modulus ($G_\tau$) is exactly 137 times stronger than its longitudinal bulk modulus ($K_\tau$), the vacuum is fiercely rigid against transverse rotation relative to longitudinal acoustic compression.

This off-diagonal rotational-transverse component of the Admittance Tensor governs the continuous transfer of energy between the internal topological vorticity ($v_{\circlearrowleft}$) of a particle and the ambient propagating transverse shear ($v_{\perp}$) of the metric.

\textbf{1. Transverse Shear (The Origin of the Coulomb Force):} \\

Because the topological boundary evaluates as an imperfect acoustic container, the internal rotation ($\nabla \times \mathbf{v}_{\psi}$) structurally agitates the surrounding metric. The absolute baseline spatial restorative force of this internal rotation evaluates via the rotational impedance limit ($E = \hbar c / r \implies F = -dE/dr$):
\begin{equation}
    F_{\circlearrowleft}(r) = \frac{\hbar c}{r^2} \tag{\ref{eq:internal_rotational_force}}
\end{equation}

A geometric fraction of this rotational energy undergoes a secondary transduction, shearing into the transverse plane to manifest as the macroscopic electric field. This leakage fraction is throttled by the off-diagonal coupling limit ($\alpha_{\circlearrowleft \perp}$). 

The macroscopic electromagnetic force ($F_{EM}$) evaluates as the internal rotational force ($F_{\circlearrowleft}$) scaled by this transverse limit:
\begin{align}
    F_{EM}(r) &= \alpha_{\circlearrowleft \perp} \cdot F_{\circlearrowleft}(r)  \\
    F_{EM}(r) &= \alpha_{\circlearrowleft \perp} \frac{\hbar c}{r^2} \tag{\ref{eq:electromagnetic_leakage}}
\end{align}
Opposing topological spins induce compounding fluid velocities between interacting boundaries, dropping local static pressure to create Bernoulli attraction. Identical spins induce fluid stagnation, spiking static pressure to create Bernoulli repulsion.

\textbf{2. The Continuous Isolation of Charge ($e^2$):} \\
The empirical Coulomb force evaluates as $F = e^2 / 4\pi \epsilon_0 r^2$. Equating this law to the fluid-dynamic leakage force isolates the geometric origin of charge:
\begin{align}
    \frac{e^2}{4\pi \epsilon_0 r^2} &= \alpha_{\circlearrowleft \perp} \frac{\hbar c}{r^2}  \\
    \frac{e^2}{4\pi \epsilon_0} &= \alpha_{\circlearrowleft \perp} \hbar c  \\
    e &= \sqrt{4\pi \epsilon_0 \hbar c \alpha_{\circlearrowleft \perp}} \tag{\ref{eq:alpha_continuum_derived}}
\end{align}
Electric charge ($e$) evaluates mechanically as the physical magnitude of the internal localized rotation that successfully leaks through the off-diagonal tensor limit ($\alpha_{\circlearrowleft \perp}$) to manifest as macroscopic transverse strain. 

\textbf{3. Macroscopic vs. Topological Transverse Impedance:} \\
The resistance of the metric to this energy transfer evaluates as the ratio of two distinct structural boundaries. The acoustic transverse impedance of the macroscopic metric ($Z_{macro}$) evaluates as the product of the absolute fluid density and the acoustic phase limit:
\begin{equation}
    Z_{macro} = \rho_{\tau} c \equiv Z_0 \tag{\ref{eq:macroscopic_impedance}}
\end{equation}

Because the topological core evaluates as a Fermionic Dirac Spinor, it requires a $4\pi$ geometric double-wrap to maintain continuous equilibrium. Energetic transfer traversing this boundary interfaces through this $4\pi$ stator, scaling the metric transition by exactly one-half ($1/2$).

The off-diagonal tensor limit ($\alpha_{\circlearrowleft \perp}$) evaluates as the mechanical ratio of the macroscopic environmental impedance ($Z_{macro}$) to the localized topological impedance ($Z_{topo}$), governed by the geometric Spinor fraction:
\begin{equation}
    \alpha_{\circlearrowleft \perp} = \frac{1}{2} \left( \frac{Z_{macro}}{Z_{topo}} \right) \tag{\ref{eq:tensor_coupling_derivation}}
\end{equation}

\textbf{4. The Absolute Topological Impedance ($Z_{topo}$):} \\
Algebraically rearranging the tensor ratio isolates the exact mechanical resistance of the topological boundary ($Z_{topo}$):
\begin{align}
    Z_{topo} &= \frac{Z_{macro}}{2 \alpha_{\circlearrowleft \perp}} \nonumber \\
    Z_{topo} &= \frac{\rho_{\tau} c}{2 \alpha_{\circlearrowleft \perp}} \tag{\ref{eq:topological_impedance}}
\end{align}
Substituting the empirical fine-structure limit ($\approx 1/137$) evaluates the absolute transverse resistance of the topological boundary as $Z_{topo} \approx 25,812 \text{ }\Omega$. This purely geometric fluid limit maps flawlessly to the legacy von Klitzing constant ($R_K = h/e^2$), deriving the quantum of resistance entirely independent of empirical charge equivalents.

\begin{tcolorbox}[recallbox, title={\ref{box:fine_structure}. Topological Coupling ($\alpha_{\circlearrowleft \perp}$) \& Electric Charge}]
\begin{align}
    F_{\circlearrowleft}(r) &= \frac{\hbar c}{r^2} \tag{\ref{eq:internal_rotational_force}} \\
    F_{EM}(r) &= \alpha_{\circlearrowleft \perp} \frac{\hbar c}{r^2} \tag{\ref{eq:electromagnetic_leakage}} \\
    e &= \sqrt{4\pi \epsilon_0 \hbar c \alpha_{\circlearrowleft \perp}} \tag{\ref{eq:alpha_continuum_derived}} \\
    Z_{macro} &= \rho_{\tau} c \equiv Z_0 \tag{\ref{eq:macroscopic_impedance}} \\
    \alpha_{\circlearrowleft \perp} &= \frac{1}{2} \left( \frac{Z_{macro}}{Z_{topo}} \right) \tag{\ref{eq:tensor_coupling_derivation}} \\
    Z_{topo} &= \frac{\rho_{\tau} c}{2 \alpha_{\circlearrowleft \perp}} \equiv R_K \tag{\ref{eq:topological_impedance}}
\end{align}
\end{tcolorbox}


% ----------------------------------------------------------------------
\section{1D Fractional Phase Flux (Dimensional Projections)}
\label{sec:deriv_tensor_projections}

Fractional charges evaluate mathematically as the continuous geometric tensor artifacts of projecting an isotropic 3D fluid circulation onto a 1D observation axis.

\textbf{1. The Isotropic 3D Circulation ($SU(3)$ Traceless Strain):} \\
As established, the $SU(3)$ traceless deviatoric strain ($\text{Tr}(\varepsilon) = 0$) evaluates as an absolutely incompressible, volume-preserving topological defect. To maintain gyroscopic stability without fracturing, the continuous phase flux ($\Phi_{\circlearrowleft}$) circulates symmetrically across all $3$ orthogonal spatial dimensions.

\textbf{2. The Longitudinal Dimensional Projection ($\parallel$):} \\
When a linear kinematic probe impacts this 3D isotropic vortex, the momentum transfer evaluates the continuous fluid flux along a single 1D Cartesian vector (the longitudinal observation axis).

The exact continuous geometric projection of a uniform 3D spherical distribution onto a single 1D longitudinal axis evaluates by integrating the squared directional cosine over the full 3D solid angle ($\Omega$). Evaluating the polar integration isolates the absolute 1D longitudinal flux fraction ($\Phi_{\parallel}$):
\begin{align}
    \Phi_{\parallel} &= \Phi_{\circlearrowleft} \left( \frac{1}{4\pi} \int_{0}^{2\pi} d\phi \int_{0}^{\pi} \cos^2 \theta \sin \theta \, d\theta \right) \nonumber \\
    \Phi_{\parallel} &= \frac{1}{3} \Phi_{\circlearrowleft} \tag{\ref{eq:tensor_long}}
\end{align}
Because the longitudinal projection measures exactly $1$ of the $3$ orthogonal spatial dimensions, it evaluates deterministically as exactly one-third of the total 3D circulation.

\textbf{3. The Transverse Dimensional Projection ($\perp$):} \\
The remainder of the isotropic 3D flux projects onto the coupled transverse geometric plane orthogonal to the observation axis. This continuous spatial projection evaluates by integrating the squared sine over the identical solid angle to isolate the transverse flux fraction ($\Phi_{\perp}$):
\begin{align}
    \Phi_{\perp} &= \Phi_{\circlearrowleft} \left( \frac{1}{4\pi} \int_{0}^{2\pi} d\phi \int_{0}^{\pi} \sin^2 \theta \sin \theta \, d\theta \right) \nonumber \\
    \Phi_{\perp} &= \frac{2}{3} \Phi_{\circlearrowleft} \tag{\ref{eq:tensor_trans}}
\end{align}
Because the transverse projection measures exactly $2$ of the $3$ orthogonal spatial dimensions, it evaluates deterministically as exactly two-thirds of the total 3D circulation.

\textbf{4. Absolute Flux Conservation:} \\
Because the geometric sum of the orthogonal spatial projections resolves identically to unity, the total continuous vortex flux is   conserved during the 1D kinematic measurement:
\begin{align}
    \Phi_{\parallel} + \Phi_{\perp} &= \frac{1}{3}\Phi_{\circlearrowleft} + \frac{2}{3}\Phi_{\circlearrowleft} \nonumber \\
    \Phi_{\parallel} + \Phi_{\perp} &= \Phi_{\circlearrowleft} \tag{\ref{eq:projection_conservation}}
\end{align}
The fractional charges ($1/3, 2/3$) evaluate as the orthogonal spatial cross-sections of a single 3D continuous volume, independent of discrete sub-components.

\textbf{5. The 1D Fractional Phase Flux ($\Phi_0$):} \\
The linear phase flux projecting across a single 1D continuous tether ($\Phi_0$) evaluates geometrically as exactly the $1/3$ longitudinal dimensional projection:
\begin{equation}
    \Phi_0 = \Phi_{\parallel} = \frac{1}{3} \Phi_{\circlearrowleft} \nonumber
\end{equation}

Substituting the established Charge-Flux Calibration limit ($\Phi_{\circlearrowleft} \equiv \Phi_e = e\sqrt{2A_{yield}} / \epsilon_{\tau}$) explicitly isolates the continuous linear phase flux of the localized projection:
\begin{align}
    \Phi_0 &= \frac{1}{3} \left( \frac{e\sqrt{2A_{yield}}}{\epsilon_{\tau}} \right) \nonumber \\
    \Phi_0 &= \frac{e\sqrt{2A_{yield}}}{3\epsilon_{\tau}} \tag{\ref{eq:1d_phase_flux}}
\end{align}
This derivation eradicates fractional charge as a discrete particulate property, redefining the subatomic sector universally as the continuous tensor projection of 3D Euclidean geometry.

\begin{tcolorbox}[recallbox, title={\ref{box:tensor_projections}. 1D Fractional Phase Flux (Dimensional Projections)}]
\begin{align}
    \Phi_{\parallel} &= \frac{1}{3} \Phi_{\circlearrowleft} \tag{\ref{eq:tensor_long}} \\
    \Phi_{\perp} &= \frac{2}{3} \Phi_{\circlearrowleft} \tag{\ref{eq:tensor_trans}} \\
    \Phi_{\parallel} + \Phi_{\perp} &= \Phi_{\circlearrowleft} \tag{\ref{eq:projection_conservation}} \\
    \Phi_0 &= \frac{e\sqrt{2A_{yield}}}{3\epsilon_{\tau}} \tag{\ref{eq:1d_phase_flux}}
\end{align}
\end{tcolorbox}

% ----------------------------------------------------------------------
\section{The Kinematic Origin of Charge \& The Magnetic Anomaly}
\label{sec:deriv_kinematic_charge}

Electromagnetism evaluates as continuous geometric fluid shear across the viscoelastic vacuum.

\textbf{1. The Absolute Shear Transduction Limit:} \\
Because the continuous metric acts as a viscoelastic medium, the kinetic energy of the rotational fluid spilling outward across the equatorial perimeter ($\frac{1}{2}\rho_{\tau} v_{\circlearrowleft}^2$) transduces into absolute transverse shear. The conversion of kinetic rotation into an elastic shear wave evaluates via energy equivalence:
\begin{equation}
    \frac{1}{2}\rho_{\tau} v_{\circlearrowleft}^2 = \frac{1}{2}\epsilon_{\tau} (\nabla \times \mathbf{A}_{\perp})^2 \tag{\ref{eq:shear_transduction_limit}}
\end{equation}

\textbf{2. Absolute Continuum Charge ($e_{\tau}$):} \\
Absolute metric charge ($e_{\tau}$) evaluates as the total physical fluid volume displaced by the continuous shear stress. It derives as the closed surface integral of the transverse strain scaled by the baseline 1D metric stiffness ($\epsilon_{\tau}$):
\begin{equation}
    e_{\tau} = \int_S \epsilon_{\tau} (\nabla \times \mathbf{A}_{\perp}) \cdot d\mathbf{A} \tag{\ref{eq:charge_displacement}}
\end{equation}

The continuum electrostatic force evaluates as the geometric interaction of two absolute charge displacements across the spatial stiffness of the metric ($F = e_{\tau}^2 / 4\pi \epsilon_{\tau} r^2$). Equating this continuum force to the macroscopic Bernoulli leakage gradient ($F = \alpha_{\circlearrowleft \perp} \hbar c / r^2$) isolates the kinematic magnitude of the displaced volume:
\begin{align}
    \frac{e_{\tau}^2}{4\pi \epsilon_{\tau} r^2} &= \alpha_{\circlearrowleft \perp} \frac{\hbar c}{r^2} \nonumber \\
    e_{\tau}^2 &= 4\pi \epsilon_{\tau} \hbar c \alpha_{\circlearrowleft \perp} \nonumber \\
    e_{\tau} &= \sqrt{4\pi \epsilon_{\tau} \hbar c \alpha_{\circlearrowleft \perp}} \tag{\ref{eq:absolute_charge_derived}}
\end{align}
Evaluating this kinematic limit via the absolute constants isolates the mechanical scalar magnitude of elementary charge: $e_{\tau} \approx 8.16 \times 10^{16}$.

\textbf{3. The Dimensional Reduction of Charge:} \\
The continuum metric evaluates electrostatic charge exclusively via fundamental dimensions (mass, length, time). Substituting the kinematic dimensions of the variables ($[\epsilon_{\tau}] = \text{kg/m}$, $[\hbar] = \text{kg}\cdot\text{m}^2/\text{s}$, $[c] = \text{m/s}$) into the absolute charge equivalence isolates the structural dimension:
\begin{align}
    [e_{\tau}] &= \sqrt{\left(\frac{\text{kg}}{\text{m}}\right) \left(\frac{\text{kg}\cdot\text{m}^2}{\text{s}}\right) \left(\frac{\text{m}}{\text{s}}\right)} \nonumber \\
    [e_{\tau}] &= \sqrt{\frac{\text{kg}^2 \cdot \text{m}^2}{\text{s}^2}} \nonumber \\
    [e_{\tau}] &= \text{kg}\cdot\text{m/s} \label{eq:charge_dimension}
\end{align}
Elementary charge evaluates structurally as macroscopic kinematic momentum displacing the transverse plane of the fluid.

\textbf{4. Dimensional Force Verification:} \\
Substituting the absolute dimensional charge back into the continuum electrostatic force equation ($F = e_{\tau}^2 / 4\pi \epsilon_{\tau} r^2$) verifies the mechanical interaction:
\begin{align}
    [F] &= \frac{(\text{kg}\cdot\text{m/s})^2}{(\text{kg/m}) \cdot \text{m}^2} \nonumber \\
    [F] &= \frac{\text{kg}^2\cdot\text{m}^2/\text{s}^2}{\text{kg}\cdot\text{m}} \nonumber \\
    [F] &= \text{kg}\cdot\text{m/s}^2 \equiv \text{Newtons} \label{eq:force_dimension}
\end{align}
The structural stiffness ($\epsilon_{\tau}$) and spherical geometry ($4\pi$) cancel symmetrically from the equivalence, reducing the interaction exactly to the established transverse leakage force ($F = \alpha_{\circlearrowleft \perp} \hbar c / r^2$).

\textbf{5. The Absolute Dirac Moment ($\mu_{B\tau}$):} \\
Evaluating the absolute kinematic charge displacement ($e_{\tau}$) and the topological angular momentum limit ($\hbar$) against the boundary mass displacement ($m$) defines the 2D baseline magnetic moment of the continuum:
\begin{equation}
    \mu_{B\tau} = \frac{e_{\tau} \hbar}{2m} \tag{\ref{eq:bohr_magneton}}
\end{equation}

\textbf{6. The Toroidal Phase Spill \& True Magnetic Moment:} \\
Because the topological torus operates as a closed 3D manifold, the transverse fluid coupling limit ($\alpha_{\circlearrowleft \perp}$) is continuously distributed across the entire rotational phase loop of the geometric boundary. 

The geometric anomaly evaluates as the transverse spill capacity divided by the complete angular phase cycle ($2\pi$). Adding this distributed toroidal phase spill to the absolute Dirac loop extracts the true continuous magnetic moment:
\begin{equation}
    a_e = \frac{\alpha_{\circlearrowleft \perp}}{2\pi} \implies \mu_{true} = \mu_{B\tau} \left( 1 + \frac{\alpha_{\circlearrowleft \perp}}{2\pi} \right) \tag{\ref{eq:true_magnetic_moment}}
\end{equation}

\begin{tcolorbox}[recallbox, title={\ref{box:kinematic_charge}. The Kinematic Origin of Charge \& $g-2$}]
\begin{align}
    \frac{1}{2}\rho_{\tau} v_{\circlearrowleft}^2 &= \frac{1}{2}\epsilon_{\tau} (\nabla \times \mathbf{A}_{\perp})^2 \tag{\ref{eq:shear_transduction_limit}} \\
    e_{\tau} &= \int_S \epsilon_{\tau} (\nabla \times \mathbf{A}_{\perp}) \cdot d\mathbf{A} \tag{\ref{eq:charge_displacement}} \\
    e_{\tau} &= \sqrt{4\pi \epsilon_{\tau} \hbar c \alpha_{\circlearrowleft \perp}} \approx 8.16 \times 10^{16} \text{ kg}\cdot\text{m/s} \tag{\ref{eq:absolute_charge_derived}} \\
    \mu_{B\tau} &= \frac{e_{\tau} \hbar}{2m} \tag{\ref{eq:bohr_magneton}} \\
    a_e = \frac{\alpha_{\circlearrowleft \perp}}{2\pi} &\implies \mu_{true} = \mu_{B\tau} \left( 1 + \frac{\alpha_{\circlearrowleft \perp}}{2\pi} \right) \tag{\ref{eq:true_magnetic_moment}}
\end{align}
\end{tcolorbox}

% ----------------------------------------------------------------------
\section{Coulomb's Law (Spherical Bernoulli \& Phase Flux)}
\label{sec:deriv_coulomb_bernoulli}

Electric charge evaluates as the geometric chirality of a localized acoustic cavity. The electrostatic force evaluates as the macroscopic kinematic pressure gradient generated by the Bernoulli effect.

\textbf{1. The Kinematic Velocity Field (Dimensional Projections):} \\
A continuous topological vortex cannot radiate purely longitudinally without violating absolute incompressibility. The total continuous phase flux ($\Phi_{\circlearrowleft}$) radiating from the central defect decomposes exactly into its orthogonal dimensional projections: the longitudinal radial expansion ($\Phi_{\parallel}$) and the transverse localized circulation ($\Phi_{\perp}$):
\begin{equation}
    \Phi_{\parallel} + \Phi_{\perp} = \Phi_{\circlearrowleft} \tag{\ref{eq:projection_conservation}}
\end{equation}
For the macroscopic electric field, the Continuity Equation mandates that the localized fluid velocity ($v$) of these orthogonal projections decreases proportionally to the spherical surface area ($A = 4\pi r^2$):
\begin{equation}
    v_{\parallel}(r) = \frac{\Phi_{\parallel}}{4\pi r^2}, \quad v_{\perp}(r) = \frac{\Phi_{\perp}}{4\pi r^2}
\end{equation}

\textbf{2. The Transverse Dynamic Pressure \& Binomial Expansion:} \\
When two transverse kinematic fields superimpose ($\mathbf{v}_{\perp} = \mathbf{v}_{\perp(1)} + \mathbf{v}_{\perp(2)}$), the baseline kinetic energy density (dynamic pressure) of the metric expands via the binomial theorem against the absolute unadulterated fluid density ($\rho_{\tau}$):
\begin{equation}
    P_{dyn} = \left[ \frac{1}{2}\rho_{\tau} v_{\perp(1)}^2 \right] + \left[ \frac{1}{2}\rho_{\tau} v_{\perp(2)}^2 \right] + \left[ \rho_{\tau} (\mathbf{v}_{\perp(1)} \cdot \mathbf{v}_{\perp(2)}) \right] \label{eq:binomial_energy}
\end{equation}
The independent squared terms represent the isolated self-energy (localized static mass equivalent) of the individual topological defects. The geometric dot product of the cross-term dictates the active interaction energy (attraction/repulsion) circulating around adjacent boundaries.

\textit{Attraction (Opposite Winding Numbers):} If the defects possess opposite chiralities, their adjacent transverse circulations align ($\mathbf{v}_{\perp(1)} \cdot \mathbf{v}_{\perp(2)} > 0$). This kinematic alignment drops the static hydrostatic pressure ($P_{static} \downarrow$) between them:
\begin{equation}
    \mathbf{v}_{\perp(1)} \cdot \mathbf{v}_{\perp(2)} > 0 \implies P_{static} \downarrow \tag{\ref{eq:attraction}}
\end{equation}

\textit{Repulsion (Identical Winding Numbers):} If the defects possess identical chiralities, their adjacent transverse circulations oppose each other ($\mathbf{v}_{\perp(1)} \cdot \mathbf{v}_{\perp(2)} < 0$). This kinematic collision forces fluid stagnation, causing the localized static hydrostatic pressure to spike:
\begin{equation}
    \mathbf{v}_{\perp(1)} \cdot \mathbf{v}_{\perp(2)} < 0 \implies P_{static} \uparrow \tag{\ref{eq:repulsion}}
\end{equation}

\textbf{3. The Volume Integral \& Admittance-Modified Force:} \\
The isolated interaction pressure deficit ($\Delta P_{int}$) evaluates from the unthrottled transverse cross-term:
\begin{equation}
    \Delta P_{int}(r) = - \rho_{\tau} (v_{\perp(1)} v_{\perp(2)}) = - \frac{\rho_{\tau} \Phi_{\perp(1)} \Phi_{\perp(2)}}{16\pi^2 r^4}
\end{equation}

The total raw interaction potential energy ($U_{raw}$) evaluates as the explicit volume integral ($\iiint \Delta P_{int} \, dV$) over the spherical metric from the boundary ($r$) to macroscopic infinity, scaled by a factor of $1/2$ to correct for mutual interaction double-counting:
\begin{align}
    U_{raw}(r) &= \frac{1}{2} \int_{r}^{\infty} \left( - \frac{\rho_{\tau} \Phi_{\perp(1)} \Phi_{\perp(2)}}{16\pi^2 r'^4} \right) 4\pi r'^2 dr'  \\
    U_{raw}(r) &= - \frac{\rho_{\tau} \Phi_{\perp(1)} \Phi_{\perp(2)}}{8\pi r}
\end{align}

Because the Bernoulli pressure gradient operates as a mechanical Surface Force, the Admittance-Modified Navier-Stokes Equation mandates that its spatial application is dimensionally throttled by the Admittance Tensor ($[\boldsymbol{\alpha}] \Sigma \mathbf{F}_{surface}$). Therefore, the absolute fluid-mechanical restoring force ($\mathbf{F}_{fluid}$) evaluates as the spatial gradient ($-\nabla U_{raw}$) multiplied by the off-diagonal Vorticity-Shear Coupling Coefficient ($\alpha_{\circlearrowleft \perp}$):
\begin{equation}
    \mathbf{F}_{fluid} = \alpha_{\circlearrowleft \perp} (-\nabla U_{raw}) = - \frac{\alpha_{\circlearrowleft \perp} \rho_{\tau} \Phi_{\perp(1)} \Phi_{\perp(2)}}{8\pi r^2} \mathbf{\hat{r}} \tag{\ref{eq:bernoulli_force}}
\end{equation}

\textbf{4. The Geometric Elimination of the Coupling Coefficient ($\alpha_{\circlearrowleft \perp}$):} \\
The discrete Coulomb force evaluates geometrically via the continuous kinematic action of the continuum, scaled by the empirical Fine-Structure Constant ($\alpha$):
\begin{equation}
    F_e = \frac{\alpha \hbar c}{r^2}
\end{equation}
Equating this macroscopic fluid Bernoulli force to the empirical kinematic limit eliminates parameterized constants ($k_e, \epsilon_0, q$) from the architecture:
\begin{equation}
    \frac{\alpha_{\circlearrowleft \perp} \rho_{\tau} \Phi_{\perp(e)}^2}{8\pi r^2} = \frac{\alpha \hbar c}{r^2} \tag{\ref{eq:force_equivalence}}
\end{equation}
Because the empirical Fine-Structure Constant ($\alpha$) evaluates phenomenologically exactly as the continuum's absolute off-diagonal coupling limit ($\alpha_{\circlearrowleft \perp} \equiv \alpha$), the impedance factor cancels symmetrically across the equivalence:
\begin{align}
    \Phi_{\perp(e)}^2 &= \frac{8\pi \hbar c}{\rho_{\tau}}
\end{align}

\textbf{5. The Pure Quantum Volumetric Phase Flux ($\Phi_{\perp(e)}$):} \\
Substituting the absolute macroscopic density of the metric ($\rho_{\tau} = \frac{9 c^5}{4\pi G^2 \hbar}$) into the isolated elementary flux:
\begin{align}
    \Phi_{\perp(e)}^2 &= \frac{8\pi \hbar c}{\left(\frac{9 c^5}{4\pi G^2 \hbar}\right)}  \\
    \Phi_{\perp(e)}^2 &= \frac{32\pi^2 G^2 \hbar^2}{9 c^4}  \\
    \Phi_{\perp(e)} &= \frac{4\sqrt{2}\pi}{3} \frac{G \hbar}{c^2} 
\end{align}

Because the geometric term $G\hbar/c^2$ evaluates exactly as the absolute Planck Area ($l_p^2$) multiplied by the kinematic phase velocity ($c$), the discrete elementary charge evaluates as a pure, unscaled continuous volumetric flux:
\begin{equation}
    \Phi_{\perp(e)} = \frac{4\sqrt{2}\pi}{3} c \space l_p^2 \tag{\ref{eq:charge_flux}}
\end{equation}
Evaluating this kinematic limit yields $\Phi_{\perp(e)} \approx 4.63 \times 10^{-61} \text{ m}^3/\text{s}$. Electric charge evaluates structurally as the absolute spatial limit of the continuum ($l_p^2$) expanding at the kinematic limit ($c$). The Fine-Structure Constant ($\alpha$) is entirely removed from the particle's intrinsic geometry, proving it is strictly the macroscopic surface force transmission limit of the surrounding boundary layer.

\begin{tcolorbox}[recallbox, title={\ref{box:coulomb_bernoulli}. Coulomb's Law (Spherical Bernoulli \& Phase Flux)}]
\begin{align}
    P_{dyn} &= \frac{1}{2}\rho_{\tau} v_{\perp(1)}^2 + \frac{1}{2}\rho_{\tau} v_{\perp(2)}^2 + \rho_{\tau} (\mathbf{v}_{\perp(1)} \cdot \mathbf{v}_{\perp(2)}) \tag{\ref{eq:binomial_energy}} \\
    \mathbf{F}_{fluid} &= - \frac{\alpha_{\circlearrowleft \perp}\rho_{\tau} \Phi_{\perp(1)} \Phi_{\perp(2)}}{8\pi r^2} \mathbf{\hat{r}} \tag{\ref{eq:bernoulli_force}} \\
    \frac{\alpha_{\circlearrowleft \perp}\rho_{\tau} \Phi_{\perp(e)}^2}{8\pi r^2} &= \frac{\alpha \hbar c}{r^2} \tag{\ref{eq:force_equivalence}} \\
    \Phi_{\perp(e)} &= \frac{4\sqrt{2}\pi}{3}c \space l_p^2 \approx 4.63 \times 10^{-61} \text{ m}^3/\text{s} \tag{\ref{eq:charge_flux}}
\end{align}
\end{tcolorbox}



% ======================================================================
% PHASE V: WAVE KINEMATICS & QUANTUM MECHANICS
% ======================================================================
\part*{PHASE V: TENSOR TRANSLATION \& FORCE UNIFICATION}






% ----------------------------------------------------------------------
\section{Tensor Evolution \& Geometric Bending (The Rosetta Stone)}
\label{sec:deriv_rosetta_stone}

Tensor curvature and geometric gauge evolution evaluate as the kinematic velocity field, spatial pressure gradients, vorticity transport, and structural incompatibility of the continuous metric fluid.

\textbf{1. The Metric Tensor ($[\boldsymbol{g}_{\mu\nu}]$) \& The Acoustic Horizon:} \\
In continuum mechanics, the metric tensor evaluates natively as the Unruh-Visser Acoustic Metric for a continuous fluid. The curved spacetime of General Relativity evaluates explicitly as the localized degradation of Eulerian fluid compliance ($\alpha_s$). 

The spacetime geometry evaluates as a $4 \times 4$ diagonal tensor. In classical relativity, the diagonal components ($g_{tt}, g_{rr}, g_{\theta\theta}, g_{\phi\phi}$) govern time dilation, radial stretching, and the spherical coordinate surface. 

In Refractive Cosmology, these components evaluate as absolute functions of the Total Scalar Admittance ($\alpha_s$). Whether modeling a localized black hole (compliance dropping to zero at $r_s$) or the macroscopic cosmos (compliance dropping to zero at the causal boundary $R_h = c \cdot t$), the tensor evaluates identical thermodynamic geometry:
\begin{equation}
    [\boldsymbol{g}_{\mu\nu}] = \begin{bmatrix}
    -\alpha_s^2 & 0 & 0 & 0 \\
    0 & \alpha_s^{-2} & 0 & 0 \\
    0 & 0 & r^2 & 0 \\
    0 & 0 & 0 & r^2 \sin^2(\theta)
    \end{bmatrix} \tag{\ref{eq:rosetta_metric_matrix}}
\end{equation}

Expanding the compliance limits ($\alpha_s^2 = 1 - P_{dyn}/P_c$), the invariant spacetime interval ($ds^2 = g_{\mu\nu} dx^\mu dx^\nu$) translates directly into the classical acoustic wave equation for the tensioned $\boldsymbol{\tau}$-fluid:
\begin{equation}
    ds^2 \equiv -\left( 1 - \frac{P_{dyn}}{P_c} \right) c^2 dt^2 + \left( 1 - \frac{P_{dyn}}{P_c} \right)^{-1} dr^2 + r^2 d\Omega^2 \label{eq:rosetta_acoustic_metric}
\end{equation}

\textbf{2. The Covariant Derivative (Material Evolution):} \\
The Covariant Derivative ($\nabla_\mu$) tracks the spatial evolution of a tensor field ($A^\nu$). It evaluates the partial derivative and the connection terms:
\begin{equation}
    \nabla_\mu A^\nu = \partial_\mu A^\nu + \Gamma^\nu_{\mu\lambda} A^\lambda 
\end{equation}
The continuous geometric transport of a localized property across the manifold evaluates as the Eulerian Material Derivative ($D/Dt$):
\begin{equation}
    \frac{D A_i}{Dt} = \frac{\partial A_i}{\partial t} + v_j \partial_j A_i 
\end{equation}
The Covariant Derivative evaluates exactly as the Material Derivative. The tensor is sheared and rotated because it propagates along the macroscopic kinematic velocity field of the fluid.
\begin{equation}
    \nabla_\mu \equiv \frac{D}{Dt} \tag{\ref{eq:rosetta_covariant}}
\end{equation}

\textbf{3. The Christoffel Symbols (Acceleration \& Shear):} \\
The geometric connection bending the metric tensor evaluates as both the kinematic acceleration and the localized deformation rate of the continuum. 

The temporal connection driving gravitational acceleration evaluates as the spatial derivative of the metric ($\Gamma^i_{00} = -\frac{1}{2}\partial^i g_{00}$). Because the temporal metric evaluates as the dimensionless fluid pressure capacity ($g_{00} = -1 + P_{dyn}/P_c$), its spatial derivative isolates the Bernoulli dynamic pressure gradient:
\begin{align}
    \Gamma^i_{00} &= -\frac{1}{2}\partial^i \left( -1 + \frac{P_{dyn}}{P_c} \right) \nonumber \\
    \Gamma^i_{00} &= -\frac{1}{2 P_c} \partial^i P_{dyn} \nonumber
\end{align}
Substituting the acoustic yield pressure ($P_c = \frac{1}{2}\rho_{\tau} c^2$) and recalling that the static pressure gradient   balances the dynamic pressure ($\nabla^i P_{static} = -\partial^i P_{dyn}$) isolates the exact Navier-Stokes kinematic acceleration ($a^i = -c^2 \Gamma^i_{00}$):
\begin{align}
    -c^2 \Gamma^i_{00} &= \frac{c^2}{2 \left(\frac{1}{2}\rho_{\tau} c^2\right)} \partial^i P_{dyn} \nonumber \\
    -c^2 \Gamma^i_{00} &= \frac{1}{\rho_{\tau}} \partial^i P_{dyn} \nonumber \\
    c^2 \Gamma^i_{00} &\equiv -\frac{1}{\rho_{\tau}} \nabla^i P_{static} \tag{\ref{eq:rosetta_christoffel_pressure}}
\end{align}

Conversely, the spatial connections evaluate as the localized fluid velocity gradient tensor.
\begin{equation}
    \Gamma^k_{ij} \equiv \partial_j v_i \tag{\ref{eq:rosetta_christoffel}}
\end{equation}
The Christoffel symbols are not abstract geometric guides; they are the continuous hydrodynamic pressure and velocity gradients physically steering the topology.

\textbf{4. Riemann Curvature (Saint-Venant Compatibility):} \\
The failure of a manifold to close a coordinate loop under parallel transport evaluates as the spatial incompatibility of the fluid. The continuous displacement ($u_i$) of the fluid generates a localized symmetric Cauchy strain tensor ($\varepsilon_{ij}$):
\begin{equation}
    \varepsilon_{ij} = \frac{1}{2}(\partial_j u_i + \partial_i u_j) 
\end{equation}
The effective macroscopic metric tensor ($g_{ij}$) evaluates as the flat Euclidean background ($\delta_{ij}$) deformed by this continuous fluid strain ($g_{ij} = \delta_{ij} + 2\varepsilon_{ij}$). 

Substituting this metric strain equivalence into the Christoffel symbols yields the connection as the spatial derivative of the strain:
\begin{equation}
    \Gamma_{ijk} = \partial_j \varepsilon_{ik} + \partial_i \varepsilon_{jk} - \partial_k \varepsilon_{ij} 
\end{equation}
Substituting this connection into the linear derivative terms of the Riemann tensor ($R_{ijkl} = \partial_i \Gamma_{jkl} - \partial_j \Gamma_{ikl}$) evaluates the second-order spatial derivatives. Subtracting the symmetric cross-terms yields:
\begin{equation}
    R_{ijkl} = \partial_j \partial_l \varepsilon_{ik} + \partial_i \partial_k \varepsilon_{jl} - \partial_j \partial_k \varepsilon_{il} - \partial_i \partial_l \varepsilon_{jk} \nonumber
\end{equation}

An arbitrary strain field ($\varepsilon_{ij}$) exists within a continuous fluid only if it satisfies the Saint-Venant Compatibility Condition ($\nabla \times \boldsymbol{\varepsilon} \times \nabla = 0$). The incompatibility tensor ($\text{Inc}(\boldsymbol{\varepsilon})$) evaluates the identical combination of second-order spatial derivatives. A topological void evaluates physically as a Volterra dislocation. Because the localized fluid strain cannot mathematically close without tearing, Riemann curvature evaluates exclusively as fluid strain incompatibility:
\begin{equation}
    R_{ijkl} \equiv \text{Inc}(\boldsymbol{\varepsilon})_{ijkl} \tag{\ref{eq:rosetta_riemann}}
\end{equation}

\textbf{5. Ricci \& Weyl Curvature (Volumetric vs. Deviatoric Strain):} \\
In General Relativity, the Riemann curvature tensor decomposes natively into the Ricci tensor ($R_{\mu\nu}$) and the Weyl tensor ($C_{\mu\nu\rho\sigma}$). This evaluates explicitly as the classical elastodynamic decomposition of total fluid strain.

The Ricci tensor evaluates the change in spatial volume. It translates   to the Isotropic Volumetric Strain (Dilatation) of the fluid metric, dictating the localized compactification of fluid density:
\begin{equation}
    R_{\mu\nu} \equiv \text{Inc}\left(\frac{1}{3}\text{Tr}(\varepsilon)\delta_{ij}\right) \implies \text{Volumetric Strain} \tag{\ref{eq:rosetta_ricci}}
\end{equation}
The Weyl tensor evaluates tidal curvature that   conserves volume. It translates exactly to the pure Deviatoric Shear Strain of the fluid metric. Gravitational waves evaluate identically as macroscopic transverse deviatoric shear waves:
\begin{equation}
    C_{\mu\nu\rho\sigma} \equiv \text{Inc}(\varepsilon'_{ij}) \implies \text{Deviatoric Shear} \tag{\ref{eq:rosetta_weyl}}
\end{equation}

\textbf{6. The Stress-Energy \& Einstein Tensors (Mechanical Equilibrium):} \\
The source of spacetime curvature ($T_{\mu\nu}$) evaluates as the Cauchy Stress Tensor ($\sigma_{\mu\nu}$) combined with the convective momentum flux of the continuous fluid. Because the Cauchy tensor natively maps normal stress to hydrostatic pressure ($\sigma_{ij} = -P \delta_{ij} + \tau_{ij}$), it physically defines the exact scalar Bernoulli pressure deficit, internal fluid shear, and kinematic momentum density at any localized coordinate.
\begin{equation}
    T_{\mu\nu} \equiv \sigma_{\mu\nu} + \rho u_\mu u_\nu \tag{\ref{eq:rosetta_stress_energy}}
\end{equation}

The Einstein Field Equation ($G_{\mu\nu} = \frac{8\pi G}{c^4} T_{\mu\nu}$) equates structural geometry ($G_{\mu\nu}$) to thermodynamic stress ($T_{\mu\nu}$). In continuous fluid mechanics, this evaluates natively as Newton's Third Law of mechanical equilibrium applied to an elastic medium. The fluid's internal geometric strain must perfectly equal the applied localized thermodynamic stress, scaled by the absolute geometric tension limit of the $\boldsymbol{\tau}$-fluid ($c^4/8\pi G$):
\begin{equation}
    G_{\mu\nu} \equiv \text{Hookean Strain Balance of } T_{\mu\nu} \tag{\ref{eq:rosetta_einstein_tensor}}
\end{equation}

\textbf{7. The Faraday Tensor (Kinematic Vorticity):} \\
The electromagnetic tensor ($F_{\mu\nu} = \partial_\mu A_\nu - \partial_\nu A_\mu$) evaluates as the macroscopic Vorticity Tensor ($\Omega_{\mu\nu}$) of the continuum. With the four-potential evaluating as fluid velocity ($A_\mu = u_\mu$), the tensor isolates the antisymmetric covariant derivative of the macroscopic velocity field.
\begin{equation}
    F_{\mu\nu} \equiv \Omega_{\mu\nu} = \partial_\mu u_\nu - \partial_\nu u_\mu \tag{\ref{eq:rosetta_faraday}}
\end{equation}

\textbf{8. Yang-Mills Evolution (Convective Acceleration):} \\
The non-Abelian $SU(3)$ field tensor ($F_{\mu\nu}$) incorporates a geometric self-interaction commutator:
\begin{equation}
    F_{\mu\nu} = \partial_\mu A_\nu - \partial_\nu A_\mu + g[A_\mu, A_\nu] 
\end{equation}
Geometric gauge self-interaction evaluates as the intrinsic mechanical property of the fluid displacing its own momentum. This evaluates natively as the Convective Acceleration tensor of the continuous fluid ($\mathbf{a}_{conv} = (\mathbf{u} \cdot \nabla)\mathbf{u} = u_j \partial_j u_i$). The non-Abelian gauge interaction evaluates as the macroscopic convective inertia of the continuum.
\begin{equation}
    \partial_\mu F^{\mu\nu} + [A_\mu, F^{\mu\nu}] \equiv (\mathbf{u} \cdot \nabla)\mathbf{u} \tag{\ref{eq:rosetta_yang_mills}}
\end{equation}

\textbf{9. Spinors \& $SU(2)$ (Cartan Moving Frames):} \\
The fundamental angular momentum of fermions evaluates via the orthogonal continuous rotation generators of localized 3D fluid axes undergoing gyroscopic precession. The localized spatial axes of a topological fluid cavity evaluate as a Cartan moving frame (a localized tetrad, $\mathbf{\hat{e}}_1, \mathbf{\hat{e}}_2, \mathbf{\hat{e}}_3$). The evolution of these localized spatial axes evaluates via the Cartan structure equations:
\begin{equation}
    \frac{d\mathbf{\hat{e}}_i}{dt} = \boldsymbol{\omega}_{fluid} \times \mathbf{\hat{e}}_i \tag{\ref{eq:rosetta_cartan}}
\end{equation}
Because the topological axes are tethered to the 3D continuous metric, a $2\pi$ rotation tangles the surrounding fluid geometry. It requires exactly $4\pi$ of geometric phase rotation to untangle the continuous background metric. The $SU(2)$ matrices evaluate explicitly as the orthogonal continuous rotation generators of this topological constraint.

\textbf{10. The Unified Non-Symmetric Field Tensor ($U_{\mu\nu}$):} \\
The total kinematic deformation of the continuous $\boldsymbol{\tau}$-fluid mathematically decomposes into exactly two components: a symmetric strain rate tensor (normal volumetric deformation) and an antisymmetric vorticity tensor (transverse shear). The unification of General Relativity and Electromagnetism evaluates purely as the superposition of the symmetric metric tensor ($g_{\mu\nu}$) and the antisymmetric Faraday tensor ($F_{\mu\nu}$):
\begin{equation}
    U_{\mu\nu} = g_{\mu\nu} + F_{\mu\nu} 
\end{equation}

Translating this into localized Cartesian space coordinates ($1, 2, 3$), the non-symmetric unified field tensor of the baseline metric evaluates as:
\begin{equation}
    U_{\mu\nu} = \begin{bmatrix}
    -\alpha_s^2 & -E_1/c & -E_2/c & -E_3/c \\
    E_1/c & \alpha_s^{-2} & -B_3 & B_2 \\
    E_2/c & B_3 & g_{22} & -B_1 \\
    E_3/c & -B_2 & B_1 & g_{33}
    \end{bmatrix}
\end{equation}

Because the geometric properties share the same continuous fluid tensor, the symmetric diagonals (Gravity) physically bound the antisymmetric off-diagonals (Electromagnetism). As localized gravitational advection drops the scalar admittance ($\alpha_s \to 0$), the fluid metric hardens, physically dampening the phase capacity for transverse electromagnetic shear.


\begin{tcolorbox}[recallbox, title={\ref{box:rosetta_stone}. Tensor Evolution \& Geometric Bending}]
\begin{align}
    [\boldsymbol{g}_{\mu\nu}] &= \text{diag}(-\alpha_s^2, \, \alpha_s^{-2}, \, r^2, \, r^2\sin^2\theta) \tag{\ref{eq:rosetta_metric_matrix}} \\
    ds^2 &\equiv -\left( 1 - \frac{P_{dyn}}{P_c} \right) c^2 dt^2 + \left( 1 - \frac{P_{dyn}}{P_c} \right)^{-1} dr^2 + r^2 d\Omega^2 \tag{\ref{eq:rosetta_acoustic_metric}} \\
    \nabla_\mu &\equiv \frac{D}{Dt} \tag{\ref{eq:rosetta_covariant}} \\
    c^2 \Gamma^i_{00} &\equiv -\frac{1}{\rho_{\tau}} \nabla^i P_{static}, \quad \Gamma^k_{ij} \equiv \partial_j v_i \tag{\ref{eq:rosetta_christoffel_pressure}} \\
    R_{ijkl} &\equiv \text{Inc}(\boldsymbol{\varepsilon})_{ijkl} \tag{\ref{eq:rosetta_riemann}} \\
    R_{\mu\nu} &\equiv \text{Volumetric Strain}, \quad C_{\mu\nu\rho\sigma} \equiv \text{Deviatoric Shear} \tag{\ref{eq:rosetta_ricci}} \\
    T_{\mu\nu} &\equiv \sigma_{\mu\nu} + \rho u_\mu u_\nu \tag{\ref{eq:rosetta_stress_energy}} \\
    G_{\mu\nu} &\equiv \text{Hookean Strain Balance of } T_{\mu\nu} \tag{\ref{eq:rosetta_einstein_tensor}} \\
    F_{\mu\nu} &\equiv \Omega_{\mu\nu} \tag{\ref{eq:rosetta_faraday}} \\
    \partial_\mu F^{\mu\nu} + [A_\mu, F^{\mu\nu}] &\equiv (\mathbf{u} \cdot \nabla)\mathbf{u} \tag{\ref{eq:rosetta_yang_mills}} \\
    \frac{d\mathbf{\hat{e}}_i}{dt} &= \boldsymbol{\omega}_{fluid} \times \mathbf{\hat{e}}_i \tag{\ref{eq:rosetta_cartan}}
\end{align}
\end{tcolorbox}


% ----------------------------------------------------------------------
\section{Continuous Pressure Dynamics (Grand Unification)}
\label{sec:deriv_pressure_gradient}

The four fundamental forces evaluate as identical geometric projections of a single continuous fluid pressure gradient. The absolute nature of this metric gradient is proven by deriving it via two independent mathematical paths.

\textbf{1. Path I: Topological Substitution (The Algebraic Proof):} \\
The localized Eulerian Bag Pressure dictates structural containment at the exact Euclidean boundary of a fundamental mass ($P_p = m^4 c^5 / 8\pi^2 \hbar^3$). 

Treating the geometric radius ($r$) as a continuous spatial variable, we substitute the inverse mass-radius limit ($m = \hbar/rc$) directly into the Bag Pressure equation to extract the continuous spatial pressure gradient:
\begin{align}
    P(r) &= \frac{(\hbar / rc)^4 c^5}{8\pi^2 \hbar^3} \nonumber \\
    P(r) &= \frac{\hbar c}{8\pi^2 r^4} \tag{\ref{eq:pressure_thermo}}
\end{align}

\textbf{2. Path II: Volumetric Strain Energy Density (The Thermodynamic Proof):} \\
In a compressible continuum, hydrostatic pressure ($P$) evaluates physically as volumetric strain energy density ($E/V$). For a continuous acoustic boundary, this internal thermodynamic state is governed by the radiation equation of state ($P = \frac{1}{3} \frac{E}{V}$). 

The topological rest energy ($E$) required to physically integrate continuous force across one full structural phase cycle ($2\pi$) evaluates via the thermodynamic work limit ($E = F \cdot r$). At any spatial radius $r$, this geometric energy limit evaluates to $E(r) = \frac{\hbar c}{2\pi r}$. The 3D Euclidean spherical volume evaluates as $V(r) = \frac{4}{3}\pi r^3$. 

Substituting the structural phase energy and the spherical volume into the acoustic equation of state isolates the continuous spatial pressure gradient:
\begin{align}
    P(r) &= \frac{1}{3} \frac{E(r)}{V(r)}  \\
    P(r) &= \frac{1}{3} \frac{\left( \frac{\hbar c}{2\pi r} \right)}{\left( \frac{4}{3}\pi r^3 \right)} \nonumber \\
    P(r) &= \frac{\hbar c}{8\pi^2 r^4} \tag{\ref{eq:pressure_thermo}}
\end{align}
The thermodynamic derivation perfectly recovers the algebraic derivation, proving the spatial pressure gradient evaluates fundamentally as the continuous volumetric strain compactness of the metric.

\textbf{3. The Unified Force Gradient:} \\
In continuous fluid mechanics, localized force evaluates as the geometric projection of this pressure across a specific spatial cross-section ($F(r) = P(r) \cdot A(r)$). Substituting the isotropic spherical bounding area ($A(r) = 4\pi r^2$) yields the unified continuum force:
\begin{align}
    F(r) &= \left( \frac{\hbar c}{8\pi^2 r^4} \right) (4\pi r^2) \nonumber \\
    F(r) &= \frac{\hbar c}{2\pi r^2} \tag{\ref{eq:pressure_gradient}}
\end{align}

\textbf{4. The Kinematic Mass Limit ($E = mc^2/2\pi$):} \\
To evaluate the mechanical energy ($E$) required to physically integrate this continuous force across its characteristic topological radius, we evaluate the thermodynamic work limit ($E = F \cdot r$):
\begin{equation}
    E(r) = \left( \frac{\hbar c}{2\pi r^2} \right) r = \frac{\hbar c}{2\pi r} 
\end{equation}
By evaluating the thermodynamic energy at the exact localized boundary of a topological mass ($r = \hbar/mc$), the geometric constants perfectly cancel:
\begin{align}
    E &= \frac{\hbar c}{2\pi (\hbar / mc)} \nonumber \\
    E &= \frac{m c^2}{2\pi} \tag{\ref{eq:energy_2pi}}
\end{align}
This geometrically proves that the baseline kinematic rest energy of a massive state ($mc^2$) evaluates as the continuous thermodynamic work of the $\boldsymbol{\tau}$-fluid integrated over one full structural rotation ($2\pi$ radians) of the localized topological boundary.

\textbf{The Fundamental Force Hierarchy (Geometric Substitution):} \\
The continuous fluid force equation mathematically unifies the four empirical fundamental forces. The isolated interaction strengths evaluate as the exact same fluid equation ($F(r)$) projected across four absolute spatial coordinates ($r_i$).

\textbf{I. Gravity (The Absolute Core, $r = l_p$):} \\
Evaluating the spatial gradient at the absolute geometric minimum of the acoustic cavity ($l_p = \sqrt{\hbar G / c^3}$) natively extracts the absolute structural tension limit of the continuum:
\begin{align}
    F_{grav} = F(l_p) &= \frac{\hbar c}{2\pi (\sqrt{\hbar G / c^3})^2} \nonumber \\
    F_{grav} &= \frac{c^4}{2\pi G} \tag{\ref{eq:force_gravity}}
\end{align}

\textbf{II. The Weak Force (The Cavitation Limit, $r = R_W$):} \\
Evaluating the continuous pressure across the Compton boundary of the topological W-Boson ($R_W = \hbar / m_W c$) isolates the volumetric yield strength of the vacuum:
\begin{align}
    F_{weak} = F(R_W) &= \frac{\hbar c}{2\pi (\hbar / m_W c)^2} \nonumber \\
    F_{weak} &= \frac{m_W^2 c^3}{2\pi \hbar} \tag{\ref{eq:force_weak}}
\end{align}

\textbf{III. The Strong Force (The Acoustic Boundary, $r = L_{crit}$):} \\
Evaluating the spatial gradient at the absolute viscoelastic relaxation limit of the continuum ($L_{crit} = \alpha \hbar / m_e c$) extracts the geometric tension envelope for topological phase equilibrium:
\begin{align}
    F_{strong} = F(L_{crit}) &= \frac{\hbar c}{2\pi (\alpha \hbar / m_e c)^2} \nonumber \\
    F_{strong} &= \frac{m_e^2 c^3}{2\pi \alpha^2 \hbar} \tag{\ref{eq:force_strong}}
\end{align}

\textbf{IV. Electromagnetism (The Macroscopic Halo, $r = a_0$):} \\
Evaluating the spatial gradient at the classical Bohr radius ($a_0 = \hbar / \alpha m_e c$) isolates the macroscopic electrostatic restoring tension:
\begin{align}
    F_{EM_{base}} = F(a_0) &= \frac{\hbar c}{2\pi (\hbar / \alpha m_e c)^2} \nonumber \\
    F_{EM_{base}} &= \frac{\alpha^2 m_e^2 c^3}{2\pi \hbar} 
\end{align}
Because the identical fluid boundary dictates the force, the ratio of the macroscopic baseline boundary force to the localized strong boundary force geometrically isolates perfectly:
\begin{equation}
    \frac{F_{EM_{base}}}{F_{strong}} = \frac{\alpha^2}{\alpha^{-2}} = \alpha^4 \tag{\ref{eq:force_em_ratio}}
\end{equation}

\textbf{V. The Macroscopic Baseline (The MOND Floor, $r = R_{h,0}$):} \\
Evaluating the spatial gradient at the absolute macroscopic scale radius of the universe ($R_{h,0}$) extracts the baseline kinematic boundary condition of the cosmic vacuum:
\begin{equation}
    F_{cosmic} = F(R_{h,0}) = \frac{\hbar c}{2\pi R_{h,0}^2} 
\end{equation}
To evaluate the kinematic acceleration floor this geometric tension exerts, we divide the force by the characteristic topological mass limit of the horizon boundary ($m_h = \hbar / R_{h,0} c$):
\begin{align}
    a_0 &= \frac{F(R_{h,0})}{m_h} \nonumber \\
    a_0 &= \left( \frac{\hbar c}{2\pi R_{h,0}^2} \right) \left( \frac{R_{h,0} c}{\hbar} \right) \nonumber \\
    a_0 &= \frac{c^2}{2\pi R_{h,0}} \tag{\ref{eq:a0z}}
\end{align}
The quantum mechanical constant ($\hbar$) mathematically annihilates, flawlessly recovering the exact macroscopic empirical MOND acceleration threshold derived from the geometric Unruh effect. This rigorously proves that galactic rotation limits are governed by the identical continuous fluid mechanics that dictate subatomic particle confinement.

\textbf{VI. Macroscopic Combinatorial Aggregates (Newtonian Gravity):} \\
The unified spatial gradient   describes the continuous limit of a singular topological defect. For a macroscopic object, the physical mass evaluates as a combinatorial superposition of localized fundamental boundaries.

The fundamental fluid mass limit ($M_0$) required to evaluate classical aggregate interactions dictates a reduced scalar:
\begin{equation}
    M_0 = \sqrt{\frac{\hbar c}{2\pi G}} 
\end{equation}

The total number of fundamental interacting boundaries ($N$) within a macroscopic mass ($M$) evaluates directly as the ratio of the aggregate mass to this fundamental fluid limit ($N = M / M_0$). Because the continuous metric interacts with itself across the spatial void, the field superimposes combinatorially. The total macroscopic force ($F_{macro}$) between two aggregate masses ($M_1$ and $M_2$) separated by a coordinate radius ($r$) evaluates as the combinatorial product of their individual boundaries ($N_1 \cdot N_2$) applied to the unified continuum force at that exact radius:
\begin{align}
    F_{macro} &= N_1 \cdot N_2 \cdot F(r) \\
    F_{macro} &= \left( \frac{M_1}{M_0} \right) \left( \frac{M_2}{M_0} \right) \left( \frac{\hbar c}{2\pi r^2} \right) \nonumber
\end{align}

Combine the aggregate mass terms:
\begin{equation}
    F_{macro} = \frac{M_1 M_2}{M_0^2} \left( \frac{\hbar c}{2\pi r^2} \right) \nonumber
\end{equation}

Substituting the squared fundamental mass limit ($M_0^2 = \hbar c / 2\pi G$) isolates the classical macroscopic force:
\begin{align}
    F_{macro} &= \frac{M_1 M_2}{\left( \frac{\hbar c}{2\pi G} \right)} \left( \frac{\hbar c}{2\pi r^2} \right) \nonumber \\
    F_{macro} &= M_1 M_2 \left( \frac{2\pi G}{\hbar c} \right) \left( \frac{\hbar c}{2\pi r^2} \right) \nonumber
\end{align}

The quantum mechanical fluid constants ($\hbar c / 2\pi$) geometrically annihilate across the numerator and denominator, flawlessly recovering the classical macroscopic measurement:
\begin{equation}
    F_{macro} = \frac{G M_1 M_2}{r^2} 
\end{equation}
Classical universal gravitation mathematically evaluates   as the macroscopic combinatorial scaling of the continuous quantum fluid boundary tension. It is not a fundamental physical law, but the direct macroscopic aggregate of the unified continuum boundary force.

\textbf{Note on Hadronic Yield (Static Confinement vs. Dynamic Fracture):} \\
While the evaluated spatial gradient ($F(L_{crit}) \approx 634 \text{ N}$) isolates the absolute static confinement pressure required by the macroscopic vacuum to structurally contain the acoustic boundary, this evaluates distinctly from the empirical Ultimate Tensile Strength of the Strong Interaction ($\approx 1 \text{ GeV/fm} \approx 160 \text{ N}$). 

Static confinement evaluates the baseline restorative crush of the continuous metric. Conversely, empirical string tension evaluates the dynamic extraction force required to stretch a composite topology (such as a Baryonic Trefoil knot) until its accumulated Hookean elastic strain breaches the $140$ MeV symmetric pair-production threshold, causing the solid-state metric to undergo brittle fracture and nucleate a pion. Confinement evaluates as the static boundary envelope, while empirical string tension evaluates as the localized kinematic fracture yield of that boundary.

\textbf{5. Empirical Numerical Evaluation \& The Coulomb Equivalence:} \\
Evaluating the unified fluid gradient ($F(r) = \hbar c / 2\pi r^2$) against standard empirical constants ($c, G, \hbar$) isolates the absolute physical tension of the metric at the respective topological boundaries:
\begin{itemize}[nosep]
    \item \textbf{Gravity ($l_p$):} $F(l_p) \approx 1.93 \times 10^{43} \text{ N}$. The absolute geometric fracture limit of spacetime.
    \item \textbf{Weak Force ($R_W$):} $F(R_W) \approx 8.35 \times 10^8 \text{ N}$. The immense localized hydrostatic pressure driving the violent $10^{-25}$ second cavitation rupture of the W-boson.
    \item \textbf{Strong Force ($L_{crit} \equiv r_e$):} $F(L_{crit}) \approx 634 \text{ N}$. The stable, macroscopic mechanical "string tension" bounding hadronic confinement.
    \item \textbf{Electromagnetism ($a_0$):} $F(a_0) \approx 1.80 \times 10^{-6} \text{ N}$. The total continuous baseline fluid pressure of the metric at the perimeter of the hydrogen atom.
\end{itemize}

Because the unified fluid gradient ($F(r)$) evaluates the strict macroscopic baseline force, the localized internal rotational force ($F_{\circlearrowleft}$) evaluates via the Legendre transformation as one full structural phase loop ($2\pi \cdot F(r) = \hbar c / r^2$). 

By the established Topological Transverse Leakage identity ($\alpha$), the observable macroscopic electric field evaluates as this internal rotation leaking through the boundary. Applying this geometric limit to the Bohr radius natively extracts the classical observable Coulomb force ($F_c$):
\begin{align}
    F_c(a_0) &= \alpha \cdot F_{\circlearrowleft}(a_0) \nonumber \\
    F_c(a_0) &= \alpha \cdot \Big[ 2\pi \cdot F(a_0) \Big] \nonumber \\
    F_c(a_0) &= \alpha \frac{\hbar c}{a_0^2} \approx 8.23 \times 10^{-8} \text{ N} 
\end{align}
This geometrically derived continuum limit flawlessly recovers the classical empirical measurement ($\frac{e^2}{4\pi \epsilon_0 a_0^2} \approx 8.23 \times 10^{-8} \text{ N}$),   proving that Coulomb's Law evaluates exclusively as the macroscopic tensor leakage of the fundamental fluid pressure gradient.

\begin{tcolorbox}[recallbox, title={\ref{box:pressure_gradient}. Continuous Pressure Dynamics}]
\begin{align}
    P(r) &= \frac{1}{3} \frac{E}{V} = \frac{\hbar c}{8\pi^2 r^4} \tag{\ref{eq:pressure_thermo}} \\
    F(r) &= P(r) \cdot A(r) = \frac{\hbar c}{2\pi r^2} \tag{\ref{eq:pressure_gradient}} \\
    E(r) &= F(r) \cdot r = \frac{\hbar c}{2\pi r} \implies E = \frac{mc^2}{2\pi} \tag{\ref{eq:energy_2pi}} \\
    F(l_p) &= \frac{c^4}{2\pi G} \tag{\ref{eq:force_gravity}} \\
    F(R_W) &= \frac{m_W^2 c^3}{2\pi \hbar} \tag{\ref{eq:force_weak}} \\
    F(L_{crit}) &= \frac{m_e^2 c^3}{2\pi \alpha^2 \hbar} \tag{\ref{eq:force_strong}} \\
    F(a_0) &= \frac{\alpha^2 m_e^2 c^3}{2\pi \hbar} \implies \frac{F_{EM_{base}}}{F_{strong}} = \alpha^4 \tag{\ref{eq:force_em_ratio}}
\end{align}
\end{tcolorbox}

% ----------------------------------------------------------------------
\section{Generalized Hooke's Law \& $SU(3)$ Symmetry}
\label{sec:deriv_hookes_su3}

A baryon evaluates geometrically as a 3D isochoric fluid cavity. The 8 fundamental geometric degrees of freedom of a volume-preserving spatial deformation tensor evaluate mathematically as the exact $SU(3)$ symmetry group of the Strong Interaction.

\textbf{1. Tensor Decomposition (Volumetric vs. Deviatoric):} \\
By classical elastodynamics, the total continuous macroscopic strain tensor ($\varepsilon_{ij}$) evaluates as the linear superposition of its isotropic volumetric expansion and its shape-distorting deviatoric shear ($\varepsilon'_{ij}$):
\begin{equation}
    \varepsilon_{ij} = \underbrace{\frac{1}{3}\text{Tr}(\varepsilon)\delta_{ij}}_{\text{Volumetric Strain}} + \underbrace{\varepsilon'_{ij}}_{\text{Deviatoric Strain}} 
\end{equation}
The corresponding Cauchy stress tensor ($\sigma_{ij}$) evaluates via Generalized Hooke's Law, partitioning identical geometric boundaries governed by the Bulk Modulus ($K$) and the Shear Modulus ($\mu_s$):
\begin{equation}
    \sigma_{ij} = K \text{Tr}(\varepsilon)\delta_{ij} + 2\mu_s \varepsilon'_{ij} 
\end{equation}

\textbf{2. The Incompressible Annihilation (Zero-Trace):} \\
As established for the stable Baryon cavity, the localized fluid circulates at the absolute acoustic limit ($v \to c$). The fluid evaluates as absolutely incompressible ($\nabla \cdot \mathbf{u} = 0$). This kinematic limit mandates the volumetric trace evaluates to exactly zero:
\begin{equation}
    \text{Tr}(\varepsilon) = \varepsilon_{11} + \varepsilon_{22} + \varepsilon_{33} = 0 \tag{\ref{eq:isochoric_trace}}
\end{equation}
Substituting this zero-divergence limit directly into the continuous strain decomposition mathematically annihilates the isotropic volumetric component ($\varepsilon_{ij} = \varepsilon'_{ij}$). For an incompressible topological cavity, the total geometric strain equates entirely to traceless deviatoric shear.

\textbf{3. The Deviatoric Stress State (Hookean $SU(3)$):} \\
Substituting the zero-trace volumetric boundary condition into Generalized Hooke's Law isolates the absolute structural restorative stress of the Baryon cavity:
\begin{align}
    \sigma_{ij} &= K(0)\delta_{ij} + 2\mu_s \varepsilon'_{ij}  \\
    \sigma_{ij} &= 2\mu_s \varepsilon'_{ij} \tag{\ref{eq:deviatoric_stress}}
\end{align}
Because the Deviatoric Strain Tensor ($\varepsilon'_{ij}$) evaluates mathematically as a $3 \times 3$ symmetric, traceless matrix, setting the trace to zero removes $1$ volumetric degree of freedom from the standard $9$. The continuous fluid deformation is now governed perfectly by exactly $8$ independent geometric generators.

\textbf{4. The Diagonal Generators (Longitudinal Strain):} \\
The continuous longitudinal strains evaluate as diagonal, traceless matrices. To form a complete mathematical basis, these generators must evaluate geometrically orthogonal and normalize to the trace inner product ($\text{Tr}(\lambda_a \lambda_b) = 2\delta_{ab}$).

\textit{Basis 1 ($\lambda_3$):} 
Evaluating a longitudinal stretch along the $x$-axis ($\varepsilon_{11} = a$) while conserving volume within the $x$-$y$ plane dictates $\varepsilon_{22} = -a$ and $\varepsilon_{33} = 0$. Applying the normalization constraint ($\text{Tr}(\lambda_3^2) = 2$) isolates $a = 1$:
\begin{equation}
    \lambda_3 = \begin{pmatrix} 1 & 0 & 0 \\ 0 & -1 & 0 \\ 0 & 0 & 0 \end{pmatrix} 
\end{equation}

\textit{Basis 2 ($\lambda_8$):} 
The second longitudinal state must conserve 3D volume ($\text{Tr}(\lambda_8) = 0$) and evaluate orthogonal to $\lambda_3$ ($\text{Tr}(\lambda_3 \lambda_8) = 0$). Let the matrix evaluate as $\text{diag}(x, y, z)$.
\begin{align}
    x + y + z &= 0 \nonumber \\
    (1)(x) + (-1)(y) + (0)(z) &= 0 \implies x = y \nonumber
\end{align}
Substituting $x = y$ into the isochoric constraint dictates $x + x + z = 0$, yielding $z = -2x$. Applying the normalization constraint ($\text{Tr}(\lambda_8^2) = 2$) natively isolates the exact geometric fraction ($x = 1/\sqrt{3}$):
\begin{equation}
    \lambda_8 = \frac{1}{\sqrt{3}} \begin{pmatrix} 1 & 0 & 0 \\ 0 & 1 & 0 \\ 0 & 0 & -2 \end{pmatrix} 
\end{equation}
The $1/\sqrt{3}$ fraction evaluates as the absolute algebraic normalization requirement for a 3D orthogonal volume-conserving fluid strain.

\textbf{5. The Off-Diagonal Generators (Shear \& Vorticity):} \\
The remaining $6$ geometric degrees of freedom evaluate as off-diagonal spatial transitions, partitioned into symmetric and antisymmetric components:

\textit{Symmetric Transverse Shear:} The real, symmetric matrices evaluate the $3$ degrees of isochoric planar sliding (e.g., $\lambda_1, \lambda_4, \lambda_6$).
\begin{equation}
    \lambda_1 = \begin{pmatrix} 0 & 1 & 0 \\ 1 & 0 & 0 \\ 0 & 0 & 0 \end{pmatrix} 
\end{equation}

\textit{Antisymmetric Vorticity:} The imaginary, antisymmetric matrices evaluate the $3$ degrees of localized geometric fluid rotation (e.g., $\lambda_2, \lambda_5, \lambda_7$). The imaginary unit ($i$) evaluates as the geometric requirement for rotational transformations to evaluate as Hermitian observables ($\lambda^\dagger = \lambda$).
\begin{equation}
    \lambda_2 = \begin{pmatrix} 0 & -i & 0 \\ i & 0 & 0 \\ 0 & 0 & 0 \end{pmatrix} 
\end{equation}
The $8$ Gell-Mann matrices evaluate explicitly as the orthogonal basis vectors for this exact macroscopic deviatoric fluid strain tensor. 

\textbf{6. Kinematic Non-Commutativity (Finite Strain):} \\
Infinitesimal geometric wave strains superimpose linearly. However, extreme spatial deformation at the localized core evaluates via the non-linear Finite Eulerian (Almansi) Strain Tensor:
\begin{equation}
    \varepsilon_{ij} = \frac{1}{2} \left( \frac{\partial u_i}{\partial x_j} + \frac{\partial u_j}{\partial x_i} - \frac{\partial u_k}{\partial x_i} \frac{\partial u_k}{\partial x_j} \right) \tag{\ref{eq:su3_finite_strain}}
\end{equation}
Because the orthogonal displacement vectors ($u_k$) are algebraically multiplied against each other within the quadratic tensor product, the strain field couples directly with its own localized displacement gradients. 

Finite rotational shears do not commute. The operation of shearing the fluid followed by rotation yields a fundamentally different geometric state than rotating followed by shearing. This geometric incompatibility dictates that the matrix multiplication of the continuous generators evaluates as non-commutative ($[\lambda_a, \lambda_b] \neq 0$), yielding the structural Lie bracket:
\begin{equation}
    [\lambda_a, \lambda_b] = 2i f_{abc} \lambda_c \tag{\ref{eq:non_abelian_commutator}}
\end{equation}
The non-Abelian self-interaction of Quantum Chromodynamics evaluates not as abstract color exchange, but as the exact non-linear kinematic product of finite 3D Eulerian fluid strain.

\begin{tcolorbox}[recallbox, title={\ref{box:hookes_su3}. Generalized Hooke's Law \& $SU(3)$}]
\begin{align}
    \text{Tr}(\varepsilon) &= 0 \implies \varepsilon_{ij} = \varepsilon'_{ij} \tag{\ref{eq:isochoric_trace}} \\
    \sigma_{ij} &= 2\mu_s \varepsilon'_{ij} \implies \text{Deviatoric } SU(3) \text{ Basis} \tag{\ref{eq:deviatoric_stress}} \\
    \varepsilon_{ij} &= \frac{1}{2} \left( \frac{\partial u_i}{\partial x_j} + \frac{\partial u_j}{\partial x_i} - \frac{\partial u_k}{\partial x_i} \frac{\partial u_k}{\partial x_j} \right) \tag{\ref{eq:su3_finite_strain}} \\
    [\lambda_a, \lambda_b] &= 2i f_{abc} \lambda_c \tag{\ref{eq:non_abelian_commutator}}
\end{align}
\end{tcolorbox}




\end{document}